\documentclass[11pt, a4paper]{article}

% --- Packages ---
\usepackage[margin=2.2cm, top=2.2cm, bottom=2.2cm]{geometry}
\usepackage{amsmath, amssymb}
\usepackage{booktabs, array}
\usepackage{enumitem}
\usepackage{xcolor}
\usepackage{titlesec}
\usepackage{parskip}

% --- Colors & Styling (matching the main formulario) ---
\definecolor{primary}{RGB}{30, 30, 30}
\definecolor{secondary}{RGB}{95, 95, 95}
\definecolor{accent}{RGB}{184, 108, 58}

\titleformat{\section}
  {\normalfont\Large\bfseries\scshape\color{primary}}
  {}{0em}
  {}
  [{\color{accent}\titlerule[0.8pt]}]
\titlespacing*{\section}{0pt}{1.4em}{0.6em}

\titleformat{\subsection}
  {\normalfont\large\bfseries\itshape\color{secondary}}
  {}{0em}
  {}
\titlespacing*{\subsection}{0pt}{1em}{0.4em}

\newcommand{\source}[1]{\hfill\textnormal{[#1]}}

% --- Document ---
\begin{document}

\begin{center}
    {\LARGE \bfseries \color{accent} Exercise Companion}\\
    \vspace{0.2cm}
    {\color{accent}\rule{0.55\linewidth}{0.8pt}}
\end{center}
\vspace{0.35cm}

This is a raw, unabridged transcription of the exercise sets from Joseph Edwards's \textit{A Treatise on the Integral Calculus with Applications, Examples and Problems}, Volume I (Macmillan and Co., 1921), companion to the theory-only formulario. Wording, numbering, and citations are reproduced exactly as printed; nothing has been reworded, abbreviated, or solved. Exercises are grouped under the article ranges they follow in the original text, so each set sits under the theory it tests.

\newpage

\begin{center}
    {\LARGE \bfseries \color{accent} Chapter II: Standard Forms --- Exercises}\\
    \vspace{0.2cm}
    {\color{accent}\rule{0.55\linewidth}{0.8pt}}
\end{center}
\vspace{0.35cm}

\section*{Exercises I --- Reversal of Differentiation \& the Logarithmic Rule (Arts. 27--42)}

Write down the indefinite integrals of:

\begin{enumerate}[leftmargin=2em]

\item $x^{10}$, $x^{-10}$, $1$, $0$, $x^{\frac{7}{5}}$, $x^{-\frac{5}{7}}$, $\sqrt[3]{x^{-2}}$, $\dfrac{1}{\sqrt{x}}$, $x^7 \times x^{-\frac{9}{2}}$,

\item $a\sqrt{x} + \dfrac{b}{\sqrt{x}}$, $a\sqrt[p]{x} + \dfrac{b}{\sqrt[p]{x}}$, $(ax^{\frac{1}{p}}+b)(cx^{\frac{1}{q}}+d)$, $\dfrac{ax^2+bx+c}{x^2}$.

\item $\dfrac{(ax^2+bx+c)(ax^{-2}+bx^{-1}+c)}{x}$, $\dfrac{1}{a-x}$, $\dfrac{1}{(a-x)^2}$, $\dfrac{1}{(a-x)^p}$.

\item $\dfrac{1}{a-x}+\dfrac{1}{a+x}$, $\dfrac{x}{a+x}$, $\dfrac{1}{(a+x)^2}+\dfrac{1}{(a-x)^2}$, $\dfrac{(x-a)(x+a)(x^2+a^2)}{x^4}$.

\item Calculate $\displaystyle\int_0^2 x^{0.3}dx$; $\displaystyle\int_3^5 x^{-\frac{1}{3}}dx$; $\displaystyle\int_1^2 \dfrac{1}{2x+3}dx$.

\item Calculate $\displaystyle\int_2^3 (a+bx^3)^2\, 3bx^2\, dx$ for the values $a=2$, $b=5$.

\item $\displaystyle\int_a^{2a} \left(\sqrt{\dfrac{a}{x}} + \sqrt{\dfrac{x}{a}}\right)^2 dx$.

\item If the retardation of a particle be 2 foot-seconds per second, and its initial velocity 10 f.s., when and where will it come to a stop?

\item Given $pv=$ constant, and that $p=40$ when $v=10$, calculate $\displaystyle\int_{10}^{20} p\, dv$. What does this integration mean?

\item Calculate $\displaystyle\int (x-1)(x-2)(x-3)(x-4)\, dx$ between limits (a) 0 and 1, (b) 1 and 2, (c) 2 and 3, (d) 3 and 4, (e) 4 and 5. Explain the signs which occur in the results. Illustrate by a graph.

\item Write down the indefinite integrals of:
\begin{enumerate}[label=(\roman*)]
\item $(ae^x+b)^n e^x$,
\item $\dfrac{ce^x}{ae^x+b}$,
\item $\left(ax+\dfrac{b}{x}+c\right)^n\left(\dfrac{ax^2-b}{x^2}\right)$,
\item $(ax^p+be^x)^n(pax^{p-1}+be^x)$.
\end{enumerate}

\item Integrate
\begin{enumerate}[label=(\roman*)]
\item $\displaystyle\int \dfrac{ae^{ax}+be^{bx}}{e^{ax}+e^{bx}}dx$,
\item $\displaystyle\int \cot 2x\, dx$,
\item $\displaystyle\int \tanh x\, dx$,
\item $\displaystyle\int \dfrac{2ax^n+b}{(ax^{2n}+bx^n+c)^p}x^{n-1}dx$.
\end{enumerate}

\item Integrate
\begin{enumerate}[label=(\roman*)]
\item $\displaystyle\int \dfrac{dx}{(1+x^2)\tan^{-1}x}$,
\item $\displaystyle\int \dfrac{dx}{(1+x^2)(\tan^{-1}x)^n}$,
\item $\displaystyle\int \dfrac{(\sin^{-1}x)^n}{\sqrt{1-x^2}}dx$,
\item $\displaystyle\int \dfrac{1}{\sin^{-1}x}\dfrac{dx}{\sqrt{1-x^2}}$,
\item $\displaystyle\int \dfrac{\text{vers}^{-1}x}{\sqrt{2x-x^2}}dx$.
\end{enumerate}

\item Integrate
\begin{enumerate}[label=(\roman*)]
\item $\displaystyle\int \dfrac{dx}{x\log x}$,
\item $\displaystyle\int \dfrac{1}{x\log x}\dfrac{1}{\log\log x}dx$,
\item $\displaystyle\int \dfrac{1}{x\log x}\dfrac{1}{\log\log x}\dfrac{1}{(\log\log\log x)^n}dx$,
\item $\displaystyle\int \dfrac{dx}{x\, l(x)l^2(x)l^3(x)\dots l^r(x)[l^{r+1}(x)]^n}$,
\end{enumerate}
where $l^r x$ represent $\log\log\log\dots\log x$, the log being repeated $r$ times.

\end{enumerate}

\section*{Exercises II --- The Preliminary Standard-Forms Table (Arts. 43--45)}

Write down the indefinite integrals of the following functions:

\begin{enumerate}[leftmargin=2em]

\item $\dfrac{1}{x+1}$, $\dfrac{x-a}{x+a}$, $\dfrac{x}{x^2+a^2}$, $\dfrac{x+a}{x^2+a^2}$, $\dfrac{x^2}{x^3+a^3}$, $\dfrac{x^{n-1}}{x^n+a^n}$.

\item $2^x$, $2x$, $\dfrac{2}{x}$, $x^2$, $x^3+3^x$, $a+b^x+c^{2x}+d^{3x}$.

\item $\cos^2\dfrac{x}{2}$, $\sin^2\dfrac{x}{2}$, $\cot x + \tan x$, $\cos x\left(\dfrac{1}{\sin x}+\dfrac{1}{\sin^2 x}\right)$.

\item $\dfrac{1}{\sqrt{9-x^2}}$, $\dfrac{1}{\sqrt{9+x^2}}$, $\dfrac{1}{\sqrt{x^2-9}}$, $\dfrac{1}{9+x^2}$, $\dfrac{1}{9-x^2}$, $\dfrac{1}{x^2-9}$.

\item $\dfrac{1}{x\sqrt{x^2-4}}$, $\dfrac{x+1}{x\sqrt{x^2-4}}$, $\dfrac{ax+b}{\sqrt{c^2-x^2}}$, $\dfrac{ax+b}{\sqrt{x^2-c^2}}$, $\dfrac{ax+b}{\sqrt{x^2+c^2}}$.

\item $\dfrac{1}{\sqrt{x-x^2}}$, $\dfrac{1}{x\sqrt{3x^2-27}}$, $\dfrac{1}{\sqrt{27-3x^2}}$, $\dfrac{x^2-4}{x^2+4}$, $\dfrac{x^2+4}{x^2-4}$.

\item Write down the indefinite integrals of:
\begin{enumerate}[label=(\roman*)]
\item $\displaystyle\int \sin^{-3}x\cos x\, dx$,
\item $\displaystyle\int \cot x \sec^2 x\, dx$,
\item $\displaystyle\int (e^x+a)^n e^x\, dx$,
\item $\displaystyle\int (x^3+a^3)^n x^2\, dx$,
\item $\displaystyle\int (ax^3+bx+c)^n(3ax^2+b)\, dx$.
\end{enumerate}

\item Write down the indefinite integrals of:
\begin{enumerate}[label=(\roman*)]
\item $\displaystyle\int \dfrac{1}{1+x^2}\dfrac{1}{\tan^{-1}x}dx$,
\item $\displaystyle\int \dfrac{dx}{\sqrt{1-x^2}(\sin^{-1}x)^2}$,
\item $\displaystyle\int \dfrac{dx}{x(\log x)^3}$.
\end{enumerate}

\item Evaluate
\begin{enumerate}[label=(\roman*)]
\item $\displaystyle\int_1^2 \dfrac{2x+1}{(x^2+x+1)}dx$,
\item $\displaystyle\int_1^2 \dfrac{2x+1}{(x^2+x+1)^2}dx$.
\end{enumerate}

\item Draw the graph of $64(x-2)(x-3)(2x-5)$, and show that the area between $x=2$ and $x=2.5$, bounded by the curve and the $x$-axis, is $32\left[(x-2)^2(x-3)^2\right]_2^{2.5}$, i.e.\ $=2$ square units. Verify by multiplying out and integrating each term.

\item Write down the values of:
\begin{enumerate}[label=(\roman*)]
\item $\displaystyle\int_0^x \dfrac{e^{3x}+e^{6x}}{e^x+e^{-x}}dx$,
\item $\displaystyle\int_0^x \dfrac{e^{(n+1)x}-e^{(n-1)x}}{\sinh x}dx$,
\item $\displaystyle\int_0^1 \dfrac{e^{2x}+e^{4x}}{e^{3x}}dx$,
\item $\displaystyle\int_a^b \cosh(\log x)\, dx$.
\end{enumerate}

\item Evaluate
\begin{enumerate}[label=(\roman*)]
\item $\displaystyle\int_0^{\pi/2} \cos x\, dx$,
\item $\displaystyle\int_0^{\pi/2}\cos^2 x\, dx$,
\item $\displaystyle\int_0^{\pi/4}\cos 2x\, dx$,
\item $\displaystyle\int_0^x (\cosh x + \cos x)\, dx$.
\end{enumerate}

\item Evaluate
\begin{enumerate}[label=(\roman*)]
\item $\displaystyle\int_0^{\pi/4n}\sec^2 nx\, dx$,
\item $\displaystyle\int_0^{\pi/4}\sec x\tan x\, dx$,
\item $\displaystyle\int_1^{\sqrt3}\dfrac{dx}{1+x^2}$,
\item $\displaystyle\int_0^1 \dfrac{dx}{\sqrt{1-x^2}}$.
\end{enumerate}

\item Evaluate
\begin{enumerate}[label=(\roman*)]
\item $\displaystyle\int \dfrac{x^n-a^n}{x-a}dx$,
\item $\displaystyle\int \dfrac{x^n}{x-a}dx$,
\item $\displaystyle\int \dfrac{x^5-2x+1}{x-1}dx$,
\item $\displaystyle\int \dfrac{x^6}{x-1}dx$,
\item $\displaystyle\int \dfrac{x^3-6x^2+11x-6}{x^2-3x+2}dx$.
\end{enumerate}

\end{enumerate}

\section*{Exercises III --- Miscellaneous (Art. 46 onward)}

\begin{enumerate}[leftmargin=2em]

\item Write down the indefinite integrals of:
\begin{enumerate}[label=(\arabic*)]
\item $\dfrac{(x+a)(x+b)(x+c)}{x^4}$,
\item $\dfrac{x^3+a^3+b^3-3abx}{x+a+b}$,
\item $\sqrt[p-q]{x^{\frac{q+r}{r-p}}}\, \sqrt[q-r]{x^{\frac{r+p}{p-q}}}\, \sqrt[r-p]{x^{\frac{p+q}{q-r}}}$,
\item $\dfrac{a\cos x - b\sin x}{a\sin x+b\cos x+c}$,
\item $e^{a\log x}+e^{x\log a}$,
\item $\dfrac{\tan^{-1}\frac{x}{3}}{9+x^2}$,
\item $\dfrac{1}{\sin x\cos x}$,
\item $\dfrac{\cos x(1+\sin x)}{\sin^2 x}$,
\item $\dfrac{1}{1-\cos x}$,
\item $\sin\left(x+\dfrac{\pi}{4}\right)$,
\item $\dfrac{x^2+\sin^2 x}{x^2+1}\sec^2 x$,
\item $(1+\sin x\cos x)\sec^2 x$,
\item $\tan x(1+\sec x)$,
\item $\dfrac{a\sin^3 x+b\cos^3 x}{\sin^2 x\cos^2 x}$,
\item $(\cos x-\sin x)(2+\sin 2x)\sec^2 x\,\text{cosec}^2 x$,
\item $(a+\tan x)(b+\tan x)\sec^2 x$,
\item $\dfrac{1}{x[1+(\log x)^2]}$,
\item $\dfrac{1}{x}\cos(\log x)$,
\item $\dfrac{x^5+1}{x-1}$,
\item $\dfrac{e^x}{a^2 e^{2x}+1}$.
\end{enumerate}

\item
\begin{enumerate}[label=(\alph*)]
\item If $f(x)=\dfrac{1}{1-x}$, prove that $\displaystyle\int fff(x)\, dx = \dfrac{x^2}{2}$.
\item If $f(x)=a+bx$, prove that $\displaystyle\int f^r(x)\, dx = a\dfrac{b^r-1}{b-1}x+b^r\dfrac{x^2}{2}$; where $f^r(x)$ means $fff\dots f(x)$, the functionality symbol $f$ occurring $r$ times.
\end{enumerate}

\item Show by expansion that
\begin{enumerate}[label=(\alph*)]
\item $\displaystyle\int \log(1+x)\, dx = \dfrac{x^2}{1.2}-\dfrac{x^3}{2.3}+\dfrac{x^4}{3.4}-\dfrac{x^5}{4.5}+\dots = (1+x)\log(1+x)-x$,
\item $\displaystyle\int \log\dfrac{1+x+x^2}{1-x+x^2}\, dx$
$=2\left\{\dfrac{x^2}{1.2}-\dfrac{2x^4}{3.4}+\dfrac{x^6}{5.6}+\dfrac{x^8}{7.8}-\dfrac{2x^{10}}{9.10}+\dfrac{x^{12}}{11.12}+\dfrac{x^{14}}{13.14}-\dfrac{2x^{16}}{15.16}+\dots\right\}$,
\item $\displaystyle\int (1+x)^n dx = x+\binom{n}{1}\dfrac{x^2}{2}+\binom{n}{2}\dfrac{x^3}{3}+\binom{n}{3}\dfrac{x^4}{4}+\binom{n}{4}\dfrac{x^5}{5}+\dots$,
\end{enumerate}
where $\displaystyle\binom{n}{r} \equiv \dfrac{n(n-1)\dots(n-r+1)}{1.2\dots r}$.

\item Prove by Differentiation or Integration from the Binomial Expansion of $(1+x)^n$, where $n$ is a positive integer,
\begin{enumerate}[label=(\alph*)]
\item $1C_1+2C_2+3C_3+\dots+nC_n = n2^{n-1}$,
\item $1.2C_2+2.3C_3+3.4C_4+\dots+(n-1)nC_n = n(n-1)2^{n-2}$,
\item $1C_1+3C_3+5C_5+\dots = n2^{n-2}$,
\item $1^2C_1+2^2C_2+3^2C_3+\dots+n^2C_n = n(n+1)2^{n-2}$,
\item $1^3C_1+2^3C_2 x+3^3C_3x^2+\dots+n^3C_nx^{n-1}$
$=x\{n^2x^2+(3n-1)x+1\}(1+x)^{n-3}$,
\item $1C_0+2C_1+3C_2+\dots+(n+1)C_n = 2^{n-1}(n+2)$,
\item $\dfrac{C_0}{1}+\dfrac{C_1}{2}+\dfrac{C_2}{3}+\dots+\dfrac{C_n}{n+1} = \dfrac{2^{n+1}-1}{n+1}$,
\item $\dfrac{C_0}{1.2}+\dfrac{C_1}{2.3}+\dfrac{C_2}{3.4}+\dots+\dfrac{C_n}{(n+1)(n+2)} = \dfrac{2^{n+2}-n-3}{(n+1)(n+2)}$,
\item $\dfrac{C_0}{1.2.3}-\dfrac{C_1}{2.3.4}+\dfrac{C_2}{3.4.5}-\dots+\dfrac{(-1)^n C_n}{(n+1)(n+2)(n+3)} = \dfrac{1}{2(n+3)}$,
\item $\dfrac{3.4}{1.2}C_0 - \dfrac{4.5}{2.3}C_1 + \dfrac{5.6}{3.4}C_2 - \dots + (-1)^n\dfrac{(n+3)(n+4)}{(n+1)(n+2)}C_n = 2\dfrac{3n+5}{(n+1)(n+2)}$.
\end{enumerate}

\item Prove from the expansions of $\sin x$ and $\cos x$ in powers of $x$ that $\displaystyle\int_0^x \sin x\, dx = 1-\cos x$ and that $\displaystyle\int_0^x \cos x\, dx = \sin x$.

\item Prove from the expansion of $\exp x$ that
$\displaystyle\int_{-\infty}^x \exp x\, dx = \exp x \quad [\exp x \equiv e^x]$.

\item Prove that
\begin{enumerate}[label=(\alph*)]
\item $\displaystyle\int_1^x \dfrac{(1+x)^n}{x}dx = \log x + \binom{n}{1}(x-1)+\binom{n}{2}\dfrac{x^2-1}{2}+\binom{n}{3}\dfrac{x^3-1}{3}+\dots$,
\item $\displaystyle\int_2^x \dfrac{(x+1)^n}{x-1}dx = 2^n\log(x-1)+\binom{n}{1}2^{n-1}[(x-1)-1]$
$+\binom{n}{2}2^{n-2}\dfrac{(x-1)^2-1}{2}+\binom{n}{3}2^{n-3}\dfrac{(x-1)^3-1}{3}+\dots$.
\end{enumerate}

\item Show that
$\displaystyle\int (a+bx)(x+c)^n dx = b\dfrac{(x+c)^{n+2}}{n+2}+(a-bc)\dfrac{(x+c)^{n+1}}{n+1}$.

\item Show that
$\displaystyle\int_a^b (1+x)^n dx = (b-a)+\binom{n}{1}\dfrac{b^2-a^2}{2}+\binom{n}{2}\dfrac{b^3-a^3}{3}+\dots+\binom{n}{n}\dfrac{b^{n+1}-a^{n+1}}{n+1}$.

\item If $\phi(x)$ be a rational integral algebraic function of $x$, show that
$\displaystyle\int_c^x \phi(x)(x-c)^n dx = (x-c)^{n+1}\left[\dfrac{\phi(c)}{n+1}+\dfrac{\phi'(c)}{n+2}\dfrac{x-c}{1!}+\dfrac{\phi''(c)}{n+3}\dfrac{(x-c)^2}{2!}+\dots\right]$.

\item By considering $\displaystyle\int (x-a)^p(x-b)^q dx$, show that the difference of the series ($p$ and $q$ being positive integers)
\[
\dfrac{(x-b)^{p+q+1}}{p+q+1}+p(b-a)\dfrac{(x-b)^{p+q}}{p+q}+\dfrac{p(p-1)}{1.2}(b-a)^2\dfrac{(x-b)^{p+q-1}}{p+q-1}+\dots,
\]
\[
\dfrac{(x-a)^{p+q+1}}{p+q+1}-q(b-a)\dfrac{(x-a)^{p+q}}{p+q}+\dfrac{q(q-1)}{1.2}(b-a)^2\dfrac{(x-a)^{p+q-1}}{p+q-1}-\dots
\]
is independent of $x$.

\item Verify by differentiation that
\begin{enumerate}[label=(\arabic*)]
\item $\displaystyle\int \dfrac{dx}{1+x^4} = \dfrac{1}{4\sqrt2}\log\dfrac{1+x\sqrt2+x^2}{1-x\sqrt2+x^2}+\dfrac{1}{2\sqrt2}\tan^{-1}\dfrac{x\sqrt2}{1-x^2}$,
\item $\displaystyle\int \dfrac{\sqrt{1-x^2}-x}{\sqrt{1-x^2}[1+x\sqrt{1-x^2}]}dx = 2\tan^{-1}(x+\sqrt{1-x^2})$.
\end{enumerate}

\item If $\phi(x)=A_0x^n+A_1x^{n-1}+A_2x^{n-2}+\dots+A_n$, prove that
\[
\int \dfrac{\phi(x)}{x-h}dx = B_0\dfrac{x^n}{n}+B_1\dfrac{x^{n-1}}{n-1}+B_2\dfrac{x^{n-2}}{n-2}+\dots+\phi(h)\log(x-h)
\]
where $B_r = hB_{r-1}+A_r$. Write down the values of $B_0$, $B_1$, $B_2$.

\item Show that $\displaystyle\int \dfrac{\phi(x)}{x-h}dx$ for rational integral algebraic forms of $\phi(x)$ may also be expressed as
\[
\phi(h)\log(x-h)+\phi'(x)\dfrac{(x-h)}{1!}+\phi''(x)\dfrac{(x-h)^2}{2!}+\phi'''(x)\dfrac{(x-h)^3}{3!}+\dots.
\]
Prove that
\[
\int \dfrac{e^x-e^h}{x-h}dx = e^h\left[\dfrac{x-h}{1^2\cdot1!}+\dfrac{(x-h)^2}{2^2\cdot2!}+\dfrac{(x-h)^3}{3^2\cdot3!}+\dfrac{(x-h)^4}{4^2\cdot3!}+\dots\right].
\]

\item Prove that $\displaystyle\int_a^b \dfrac{\log x}{x}dx = \dfrac{1}{2}\log(ab)\log\left(\dfrac{b}{a}\right)$.

\item If
\[
J_n(x) = \dfrac{x^n}{2^n n!}\left[1-\dfrac{x^2}{2(2n+2)}+\dfrac{x^4}{2.4(2n+2)(2n+4)}-\dots\right]
\]
(i.e.\ Bessel's function), prove that
\begin{enumerate}[label=(\arabic*)]
\item $\displaystyle\int_0^x J_1(x)\, dx = 1-J_0(x)$,
\item $\displaystyle\int_0^x [J_{n-1}(x)-J_{n+1}(x)]\, dx = 2J_n(x)\quad(n>0)$.
\end{enumerate}

\item Prove that
\[
\lambda\xi^{-\lambda}\int_0^\xi x^{\lambda-1}(1-x^\mu)^\nu dx = F\left(-\nu,\dfrac{\lambda}{\mu},\dfrac{\lambda}{\mu}+1,\xi^\mu\right)
\]
when $F(\alpha,\beta,\gamma,x)$ denotes the hypergeometric series
\[
1+\dfrac{\alpha.\beta}{1.\gamma}x+\dfrac{\alpha\,\overline{\alpha+1}}{1.2}\dfrac{\beta\,\overline{\beta+1}}{\gamma\,\overline{\gamma+1}}x^2+\dfrac{\alpha\,\overline{\alpha+1}\,\overline{\alpha+2}}{1.2.3}\cdot\dfrac{\beta\,\overline{\beta+1}\,\overline{\beta+2}}{\gamma\,\overline{\gamma+1}\,\overline{\gamma+2}}x^3+\dots.
\]
\source{I. C. S., 1898.}

\item Assuming that the speed of the current in a river at a distance $x$ from the bank follows the law
\[
v=v_0+kx(a-x),
\]
where $a$ is the breadth of the river and $v_0$ and $k$ are constants, find by integration how far down stream a man will be carried who rows 4 miles an hour, pointing the boat's head always straight at the opposite bank, so as to cross in the least time possible: the width of the river being half a mile, the banks being straight and parallel, and the speed of the current being 2 miles an hour near the banks, and 3 miles an hour in mid-stream. \source{I. C. S., 1905.}

\item Find the moment of inertia of a rectangle of sides $2a$, $2b$ about a line joining the mid-points of the opposite sides of length $2a$.

The section of a ship at the water line is 120 feet long. If the middle line be divided into six equal portions, the ordinates of the boundary of the area at the middle points of the segments are given by the following table:

\begin{center}
\begin{tabular}{@{}l|cccccc@{}}
\toprule
Distances from end & 10 & 30 & 50 & 70 & 90 & 110 \\
\midrule
Ordinates & 10 & 20 & 20 & 18 & 14 & 8 \\
\bottomrule
\end{tabular}
\end{center}

Draw a figure showing the section, remembering that it is symmetrical on both sides of the middle line. By approximate methods of summation (treating the segments as rectangles) find the value of $Ak^2$, the moment of inertia of the section about the middle line, and the height of the metacentre above the centre of buoyancy if the ship displaces 36,000 cubic feet of water.
\[
\left[\text{The height required} = \dfrac{\text{Moment of Inertia}}{\text{Displacement}}.\right]
\]
\source{I. C. S., 1908.}

\item A substance $A$ transforms into a substance $B$, the rate of transformation in grammes per second at the time $t$ being equal to $ax$, where $x$ denotes the number of grammes of $A$ existing at that instant. In like manner $B$ transforms into a third substance $C$, the rate of transformation being $by$, where $y$ is the number of grammes of $B$ existing at time $t$.

Write down the relations between $x$, $y$ and their differential coefficients with respect to the time, and show that these equations are satisfied by putting
\[
x=Pe^{-at}, \quad y=Qe^{-at}+Re^{-bt},
\]
provided $Q$ and $R$ are properly determined in terms of $P$, $y$ being zero when $t=0$.

Also show that the quantity of $B$ existing will be greatest at
$\dfrac{\log\text{Nap}\,a - \log\text{Nap}\,b}{a-b}$ seconds after the zero of time. \source{I. C. S., 1908.}

\item Prove that the integral of the sum of an infinite series taken over a range within which the series is absolutely convergent, is equal to the sum of the integrals of the terms.

Employ this theorem to find an expansion of $\log(1+z)$ in ascending powers of $z$, pointing out the range for which the expansion is valid.

Having given $\log_e 2 = 0.69315$, prove that $\log_e 61 = 4.1109$ to five significant figures. \source{I. C. S., 1904.}

\item $OK$ is the diameter bounding a semicircle of radius $r$, $P$ any point on $OK$, and $PQ$ an ordinate to the diameter $OK$. If $x$ denote the length $OP$ and $z$ the area which $PQ$ cuts off the semicircle, interpret $\dfrac{dz}{dx}$ and $\dfrac{d^2z}{dx^2}$.

Find a curve for which the area bounded by the curve, the axes of $x$ and $y$ and the ordinate at a distance $x$ from the axis of $y$, is $a^2\tan\dfrac{x}{a}$. \source{I. C. S., 1902.}

\item From the equation
\[
\dfrac{1}{y}\left(ah+\int_h^x y\, dx\right)=h,
\]
where $a$ and $h$ are constants, find $y$ in terms of $x$.

The value of $a$ being 2 feet, and of $h$ 10 feet, evaluate $y$ when $x$ is 30 feet. \source{I. C. S., 1910.}

\item Denoting by $A$ the area between the curve $y=f(x)$ and the axis of $x$, from the value zero to the value $a$ of $x$, show that, when $f(x)$ is a rational integral algebraic function of the third degree,
\[
A=\dfrac{a}{6}(y_0+4y_1+y_2),
\]
where $y_0=f(0)$, $y_1=f\left(\dfrac{a}{2}\right)$, $y_2=f(a)$.

Compare the result given by this rule with the true value, taken to three places of decimals, for the curve $y=\sin x$, between the values 0 and 0.5 of $x$ reckoned in radians. \source{I. C. S., 1912.}

\item Verify that the area of the curve
\[
y=A+Bx+Cx^2+Dx^3
\]
between the limits $x=h$ and $x=-h$ is equal to the product of $h$ and the sum of the ordinates at
\[
x=h/\sqrt3 \quad\text{and}\quad x=-h/\sqrt3.
\]
In the case of the curve
\[
y=A+Bx+Cx^2+Dx^3+Ex^4+Fx^5\equiv f(x),
\]
verify in like manner that the area between $x=h$ and $x=-h$ is equal to
\[
\{5f(h\sqrt{3/5})+8f(0)+5f(-h\sqrt{3/5})\}h/9.
\]
\source{I. C. S., 1913.}

\item Find the differential coefficient of
\[
e^{-x}(1+x+x^2/2!+x^3/3!+x^4/4!);
\]
and deduce that the sum of the first five terms of the exponential series is less than $e^x$ by the quantity
\[
e^x\int_0^x \dfrac{a^4e^{-a}}{4!}da.
\]
What would be the corresponding result if the series were taken to $n$ terms instead of to five terms? \source{I. C. S., 1913.}

\item Weddle's Rule for finding, approximately, the area bounded by a curve, two ordinates, and a base forming part of the axis of $x$, is: ``Divide the base into six equal parts and draw the ordinates at the points of division, making, with the extreme ordinates, seven in all. Of these ordinates add the first, third, fourth, fifth and seventh, and five times the second, fourth and sixth. Multiply the sum by one-twentieth of the base.''

Prove that the rule gives the true result when the limits are 0 and 1 and the curve has any of the forms
\[
y=a, \quad y=ax^2, \quad y=ax^4, \quad y=a(x-1/2)^n,
\]
where $a$ is a constant and $n$ is an odd positive integer.

Find, by the rule, the value of $6\times\displaystyle\int_0^{1/\sqrt3}\dfrac{dx}{1+x^2}$ to seven places of decimals. Check the result by integration. \source{I. C. S., 1911.}

\item Show that the work done by a gas in altering its volume from $v_1$ to $v_2$ according to the Adiabatic Law
\[
pv^\gamma = p_0v_0^\gamma
\]
is
\[
\dfrac{p_0v_0^\gamma}{\gamma-1}\left[\dfrac{1}{v_1^{\gamma-1}}-\dfrac{1}{v_2^{\gamma-1}}\right].
\]
If the law be Isothermal ($\gamma=1$), show that this becomes
\[
p_0v_0\log\dfrac{v_2}{v_1}.
\]
If a gas expands isothermally from state $p_1$, $v_1$ to state $p_2$, $v_2$ (Operation I.), then expands adiabatically from state $p_2$, $v_2$ to state $p_3$, $v_3$ (Operation II.), then contracts isothermally from state $p_3$, $v_3$ to state $p_4$, $v_4$ (Operation III.), then contracts adiabatically from state $p_4$, $v_4$ to state $p_1$, $v_1$ (Operation IV.),
\begin{enumerate}[label=(\arabic*)]
\item find the amounts of work done by or upon the gas during each of these four operations, drawing a graph of the whole cycle of changes;
\item show that the work done in the whole cycle of operations is measured by
\[
(p_1v_1-p_3v_3)\log\dfrac{v_2}{v_1} \text{ (the adiabatic portions cancelling)};
\]
\item that $v_1v_3 = v_2v_4$;
\item that, writing
\[
p_1v_1=p_2v_2=\alpha_1, \qquad p_3v_3=p_4v_4=\alpha_2,
\]
\[
p_2v_2^\gamma=p_3v_3^\gamma=\beta_1, \qquad p_4v_4^\gamma=p_1v_1^\gamma=\beta_2,
\]
the above expression for the work may be written
\[
\dfrac{\alpha_1-\alpha_2}{\gamma-1}\log\dfrac{\beta_1}{\beta_2}.
\]
\end{enumerate}
[This cycle of operations is known as a Carnot's cycle for a perfect heat engine.]

\item If $dQ$ be the whole heat absorbed by a body of uniform temperature whilst its temperature changes continuously from $\theta$ to $\theta+d\theta$, and if $\phi$ be a function of the independent variables which define the state of the body and such that
\[
d\phi = \dfrac{dQ}{\theta},
\]
$\phi$ is called the Entropy of the body (Clausius).

Show that if a graph be drawn to represent $\theta$ as a function of $\phi$, the area between the graph, the $\phi$-axis and the ordinates corresponding to the initial and final states represents on some scale the heat absorbed.

In the case of a perfect gas satisfying the law $\dfrac{pv}{\theta}=\text{const.}=R$, assume the Thermodynamic Equation
\[
dQ = C_v d\theta + p\, dv,
\]
where $C_v$ is the specific heat at constant volume, and show that in changing from state $\theta_1$, $v_1$, $\phi_1$ to state $\theta_2$, $v_2$, $\phi_2$,
\[
\left(\dfrac{\theta_2}{\theta_1}\right)^{C_v}\left(\dfrac{v_2}{v_1}\right)^R = \dfrac{e^{\phi_2}}{e^{\phi_1}}.
\]

Taking the temperature as a function of the entropy and simultaneous values of $\phi$ and $\theta$ as given in the following table:

\begin{center}
\begin{tabular}{@{}l|cccc@{}}
\toprule
$\phi$ & 1.60 & 1.70 & 1.75 & 1.85 \\
\midrule
$\theta$ & 450 & 400 & 370 & 340 \\
\bottomrule
\end{tabular}
\end{center}

and assuming that one unit of area indicates one unit of heat, what is the total heat received during these changes?

[There is a brief sketch of the fundamental formulae of Thermodynamics on pages 56 and 57 of \textit{Solutions of Senate House Problems for 1878} which may be found useful. Students may also read Tait's \textit{Thermodynamics} or Parker's \textit{Thermodynamics} for detailed accounts of the theory; other useful books are Zeuner, \textit{Th\'eorie M\'ecanique de la Chaleur}; Briot, \textit{Th\'eorie M\'ecanique de la Chaleur}.]

\item In the case of a saturated vapour, if $C'$ be the specific heat of the vapour, i.e.\ the heat imparted to one gramme of the saturated vapour to keep it constantly in the saturated state when slowly compressed till the temperature rises one degree Fahrenheit; $C$ that of the liquid from which it is derived at the same pressure and temperature, $L$ the latent heat, then it can be shown that
\[
C' = C+\dfrac{dL}{d\theta}-\dfrac{L}{\theta},
\]
where $\theta$ is the absolute temperature.

Let $C_p$ be the specific heat of the liquid at constant pressure, which, as liquids are practically incompressible, is so nearly equal to $C$ that no appreciable error results from regarding them as identical.

Then Regnault has shown experimentally that the sum of the free and latent heat, viz.\ $L+\displaystyle\int_{273}^\theta C_p\, d\theta$, is not a constant as had been supposed by Watt in his earlier experiments, but is a function of the temperature $\theta$, viz.\ putting $\theta = 273+\theta'$ and $J$ being the number of ergs in one calorie (41,539,739.8 ergs or about 3 foot-lbs.), he obtained the equations
\begin{enumerate}[label=(\arabic*)]
\item $L+\displaystyle\int_{273}^\theta C_p\, d\theta = J(\alpha+\beta\theta'+\gamma\theta'^2)$,
\item $C_p = J(\alpha'+\beta'\theta'+\gamma'\theta'^2)$,
\end{enumerate}
experimentally, determining the constants $\alpha$, $\beta$, $\gamma$; $\alpha'$, $\beta'$, $\gamma'$ for several vapours.

Using these data, prove that
\begin{enumerate}[label=(\arabic*)]
\item $\dfrac{dL}{d\theta}+C_p = J(\beta+2\gamma\theta')$,
\item $L=J\left[\alpha+(\beta-\alpha')\theta'+(2\gamma-\beta')\dfrac{\theta'^2}{2}-\gamma'\dfrac{\theta'^3}{3}\right]$,
\item $c'\equiv\dfrac{C'}{J}=\beta+2\gamma\theta'-\dfrac{\alpha+(\beta-\alpha')\theta'+(2\gamma-\beta')\frac{\theta'^2}{2}-\gamma'\frac{\theta'^3}{3}}{273+\theta'}$.
\end{enumerate}

\item Show that the integral equivalent of the equation
\[
\int_0^x \dfrac{dx}{1+x+x^2}+\int_0^y \dfrac{dy}{1+y+y^2}+\int_0^z \dfrac{dz}{1+z+z^2}=0
\]
is of the form
\[
xyz+a(yz+zx+xy)+b(x+y+z)+c=0,
\]
where $a$, $b$, $c$ are certain constants. \source{Oxf. I. P., 1913.}

\item If the variables $x$, $y$, $z$ be so related that
\[
yz=F(x), \quad zx=F(y), \quad xy=F(z),
\]
show that $\displaystyle\int_{x_1}^x F(x)\, dx+\int_{y_1}^y F(y)\, dy+\int_{z_1}^z F(z)\, dz = xyz-x_1y_1z_1$.

For example, if $x+y+z=0$,
and $yz+zx+xy=-\dfrac14 x^2y^2z^2$,
show that
\[
\int_{x_1}^x \dfrac{\sqrt{1+x^4}}{x^2}dx+\int_{y_1}^y \dfrac{\sqrt{1+y^4}}{y^2}dy+\int_{z_1}^z \dfrac{\sqrt{1+z^4}}{z^2}dz = \dfrac34(xyz-x_1y_1z_1).
\]
\source{Bertrand, Calc. Int., p. 383.}

\item If $y=\displaystyle\int_0^x e^{\sin x}dx$, expand $y$ in powers of $x$ as far as $x^5$. \source{Oxf. I. P., 1911.}

\item Prove that $\displaystyle\int_0^\infty \dfrac{x^2+3x+3}{(x+1)^3}e^{-x}\sin x\, dx = \dfrac12$. \source{Oxf. I. P., 1915.}

\item Integrate $\displaystyle\int x^{-2n-2}(1-x)^n(1-cx)^n dx$, where $n$ is a positive integer. \source{Oxf. I. P., 1917.}

\item Prove that the fifth differential coefficient of
\[
x^4\int \phi(x)\, dx - 4x^3\int x\phi(x)\, dx + 6x^2\int x^2\phi(x)\, dx - 4x\int x^3\phi(x)\, dx + \int x^4\phi(x)\, dx
\]
is
\[
24\phi(x). \source{Oxf. I. P., 1917.}
\]

\item Integrate $\displaystyle\int \dfrac{\cos 8\theta - \cos 7\theta}{1+2\cos 5\theta}d\theta$.

\item If $f(x)$ and $F(x)$ be two functions continuous and finite between $0$ and $x$, such that
\[
F(x) \equiv \int_0^x f(t)\, dt,
\]
\[
f(x) \equiv 1-\int_0^x F(t)\, dt,
\]
obtain their expansions in ascending powers of $x$. \source{Oxf. I. P., 1915.}

\item Prove that $\displaystyle\int_0^1 \dfrac{dx}{(1+x)(2+x)} = 0.288$ nearly. \source{Math. Trip. I., 1916.}

\item If $y^2=a^2x^2+c$, express in terms of $y$ the differential coefficients of the functions
\[
\log(ax+y), \quad xy
\]
with regard to $x$.

Hence evaluate $\displaystyle\int \dfrac{dx}{y}$ and $\displaystyle\int y\, dx$, and prove that
\[
\int_2^{\frac{5}{2}}\sqrt{x^2-4}\, dx = \dfrac{15}{8}-2\log 2. \source{Math. Trip. I., 1914.}
\]

\item Prove that
\[
\{\log(a+\theta h)-\log a\}-\theta\{\log(a+h)-\log a\}
\]
can be expressed in the form
\[
\theta(1-\theta)\int_0^h \dfrac{x\, dx}{(a+x)(a+\theta x)}.
\]
Deduce that in calculating a logarithm to base 10 by the method of proportional parts from tables which give the logarithms of all integers from $10^4$ to $10^5$, the error is one of defect and cannot amount to $\frac18 10^{-8}\mu$, where $\mu = \log_{10}e = .43429$. Is this negligible in seven-figure tables? \source{Math. Trip. Pt. II., 1919.}

\item Integrate $\displaystyle\int \sin\theta\sqrt{\dfrac{1+\cos 9\theta}{1+\cos\theta}}\, d\theta$, and show that
\[
\int \sqrt{\dfrac{1+\cos(4n+1)\theta}{1+\cos\theta}}\, d\theta = \theta - 2\left\{\dfrac{\sin\theta}{1}-\dfrac{\sin 2\theta}{2}+\dfrac{\sin3\theta}{3}-\dots-\dfrac{\sin 2n\theta}{2n}\right\},
\]
$n$ being a positive integer.

\end{enumerate}

\newpage

\begin{center}
    {\LARGE \bfseries \color{accent} Chapter III: Change of the Independent Variable --- Exercises}\\
    \vspace{0.2cm}
    {\color{accent}\rule{0.55\linewidth}{0.8pt}}
\end{center}
\vspace{0.35cm}

\section*{Exercises I --- Substitution Technique (Arts. 47--58)}

\begin{enumerate}[leftmargin=2em]

\item Evaluate
\begin{enumerate}[label=(\roman*)]
\item $\displaystyle\int \dfrac{3x^2}{1+x^3}dx$. Put $x^3=z$.
\item $\displaystyle\int_0^1 \dfrac{dx}{x^2+4x+5}$. Put $x+2=z$.
\item $\displaystyle\int_2^3 \dfrac{dx}{(x-1)\sqrt{x^2-2x}}$. Put $x-1=z$; $\therefore \dfrac{dx}{x-1}=\dfrac{dz}{z}$, etc.
\item $\displaystyle\int \dfrac{dx}{e^x+e^{-x}}$. Put $e^x=z$.
\item $\displaystyle\int \dfrac{e^x\,dx}{2e^{2x}+2e^x+1}$. Put $e^x=\dfrac{z}{1-z}$.
\item $\displaystyle\int \tan^2x\sec^2x\,dx$.
\item $\displaystyle\int \dfrac{dx}{\cosh^2 mx}$.
\end{enumerate}

\item Evaluate
\begin{enumerate}[label=(\roman*)]
\item $\displaystyle\int_0^a \sqrt{a^2-x^2}\,dx$. Put $x=a\sin\theta$.
\item $\displaystyle\int_0^{2a} \sqrt{2ax-x^2}\,dx$. Put $x=a(1-\cos\theta)$.
\end{enumerate}
Draw graphs to illustrate these two integrations.

\item Find the values of
\begin{enumerate}[label=(\roman*)]
\item $\displaystyle\int_0^a x\sqrt{a^2-x^2}\,dx$.
\item $\displaystyle\int_0^a x^2\sqrt{a^2-x^2}\,dx$.
\end{enumerate}
Interpret the meaning of these integrations.

\item Integrate $\displaystyle\int \dfrac{(ax^2-b)\,dx}{x\sqrt{c^2x^2-(ax^2+b)^2}}$. Put $ax+\dfrac{b}{x}=z$.

\item Integrate $\displaystyle\int \dfrac{x^{2n}\,dx}{(a^2+x^2)^{n+\frac32}}$. Put $x=a\tan\theta$. \source{St. John's, 1883.}

\item Integrate $\displaystyle\int \sec^{\frac23}\theta\,\text{cosec}^{\frac23}\theta\,d\theta$. Put $\tan\theta=z$. \source{St. John's, 1883.}

\item Integrate $\displaystyle\int \sqrt{\dfrac{\sin x}{\cos^5 x}}\,dx$. Put $\tan x=z$. \source{Trinity, 1883.}

\item Integrate
\begin{enumerate}[label=(\roman*)]
\item $\displaystyle\int \dfrac{dx}{x\sqrt{x^4-1}}$.
\item $\displaystyle\int \dfrac{dx}{x\sqrt{1-x^4}}$.
\item $\displaystyle\int \dfrac{dx}{x\sqrt{1+x^4}}$.
\end{enumerate}

\item Integrate
\begin{enumerate}[label=(\roman*)]
\item $\displaystyle\int \left(1-\dfrac{1}{x^2}\right)e^{x+\frac{1}{x}}dx$.
\item $\displaystyle\int \dfrac{ax^2-b}{x^3+(ax^2+b)^2}dx$.
\item $\displaystyle\int \dfrac{ax^2-b}{x^{n+2}}(ax^2+b)^n dx$.
\item $\displaystyle\int \dfrac{1}{(c+ex)^2}\cos\dfrac{a+bx}{c+ex}dx$.
\item $\displaystyle\int \dfrac{e^{a\tan^{-1}x}}{1+x^2}dx$.
\item $\displaystyle\int \dfrac{e^{a\sin^{-1}x}}{\sqrt{1-x^2}}dx$.
\item $\displaystyle\int \dfrac{\sin x\cos x\,dx}{a^2\cos^2x+b^2\sin^2x}$.
\end{enumerate}

\item Integrate
\begin{enumerate}[label=(\roman*)]
\item $\displaystyle\int \{\phi(x)\psi'(x)+\phi'(x)\psi(x)\}\,dx$.
\item $\displaystyle\int \dfrac{\phi(x)\psi'(x)-\phi'(x)\psi(x)}{[\phi(x)]^2}dx$.
\item $\displaystyle\int \dfrac{\phi'(x)\,dx}{1+[\phi(x)]^2}$.
\item $\displaystyle\int e^{\phi(x)}\phi'(x)\,dx$.
\item $\displaystyle\int e^{-\psi(x)}\dfrac{\phi'(x)-\phi(x)\psi'(x)\log\phi(x)}{\phi(x)}dx$.
\end{enumerate}

\item Show that
\[
\int (x-a)\sqrt{\dfrac{x-4a}{x}}\,dx = a^2(\sinh 2u - 2\sinh u)
\]
where $x=4a\cosh^2\dfrac{u}{2}$. \source{Ox. I. Pub., 1899.}

\end{enumerate}

\section*{Exercises II --- The Hyperbolic Functions (Arts. 59--73)}

Establish the following results:

\begin{enumerate}[leftmargin=2em]

\item $\displaystyle\int \cosh x\,dx=\sinh x$.

\item $\displaystyle\int \sinh x\,dx=\cosh x$.

\item $\displaystyle\int \text{sech}^2x\,dx=\tanh x$.

\item $\displaystyle\int \text{cosech}^2x\,dx=-\coth x$.

\item $\displaystyle\int \dfrac{\sinh x}{\cosh^2 x}dx=\int \text{sech}\,x\tanh x\,dx=-\text{sech}\,x$.

\item $\displaystyle\int \dfrac{\cosh x}{\sinh^2 x}dx=\int \text{cosech}\,x\coth x\,dx=-\text{cosech}\,x$.

\item $\displaystyle\int \text{cg}\,x\,dx=\text{gd}\,x$.

\item $\displaystyle\int \text{cg}^2x\,dx=\text{sg}\,x$.

\item $\displaystyle\int \dfrac{dx}{\text{cg}\,x}=\text{tg}\,x$.

\item
\begin{enumerate}[label=(\alph*)]
\item $\text{sg}(u+v)=\dfrac{\text{sg}\,u+\text{sg}\,v}{1+\text{sg}\,u\,\text{sg}\,v}$,
\item $\text{cg}(u+v)=\dfrac{\text{cg}\,u\,\text{cg}\,v}{1+\text{sg}\,u\,\text{sg}\,v}$,
\item $\text{sg}\,2u=\dfrac{2\,\text{sg}\,u}{1+\text{sg}^2u}$,
\item $\text{cg}\,2u=\dfrac{\text{cg}^2u}{1+\text{sg}^2u}$,
\item $\text{tg}\,2u=2\dfrac{\text{sg}\,u}{\text{cg}^2u}$.
\end{enumerate}

\end{enumerate}

\section*{Exercises III --- Cosec, Sec \& Radical Standard Forms (Arts. 74--89)}

Write down the integrals of

\begin{enumerate}[leftmargin=2em]

\item $\dfrac{1}{9-x^2}$, $\dfrac{1}{9-4x^2}$, $\dfrac{1}{x^2-4}$, $\dfrac{1}{9x^2-4}$, $\sqrt{16-9x^2}$, $\sqrt{3x^2-5}$, $\sqrt{2+3x^2}$.

\item $\dfrac{1}{\sqrt{x(x-4)}}$, $\dfrac{1}{\sqrt{x(4-x)}}$, $\dfrac{1}{\sqrt{x(4+x)}}$, $\dfrac{1}{\sqrt{2+2x-x^2}}$, $\dfrac{1}{\sqrt{x^2-2x+2}}$, $\sqrt{x^2+2ax}$.

\item $\dfrac{x}{\sqrt{9-x^2}}$, $\dfrac{x}{\sqrt{x^2-9}}$, $\dfrac{x}{\sqrt{9-4x^2}}$, $\dfrac{x^2}{\sqrt{1-x^2}}$, $\dfrac{x^2}{\sqrt{x^2+1}}$.

\item $x\sqrt{x^2+a^2}$, $(x+b)\sqrt{x^2+a^2}$, $\dfrac{ax+b}{\sqrt{x^2+c^2}}$.

\item $x(x^2+a^2)^{\frac{n}{2}}$, $(x+a)(x^2+2ax+b)^{\frac{n}{2}}$, $(ax-b)(ax^2-2bx+c)^{\frac{n}{2}}$.

\item $\dfrac{x^2+2x+3}{\sqrt{1-x^2}}$, $\dfrac{x^2+2x+3}{\sqrt{x^2+1}}$, $\dfrac{x^2+2x+3}{\sqrt{x^2+x+1}}$, $\dfrac{x^2+ax+b}{\sqrt{x^2+cx+d}}$.

\item $\sqrt{x^2+4x+5}$, $\sqrt{-x^2+4x+5}$, $\sqrt{4x^2+4x+5}$, $\sqrt{-4x^2+4x+5}$.

\item $\sqrt{\dfrac{x+a}{x-a}}$, $\sqrt{\dfrac{a+x}{a-x}}$, $x\sqrt{\dfrac{a+x}{a-x}}$, $(x+a)\sqrt{\dfrac{x+b}{x-b}}$, $\dfrac{(x+a)^{\frac32}}{(x-a)^{\frac12}}$.

\item $\text{cosec}\,nx$, $\text{cosec}(2x+b)$, $\dfrac{1}{4\cos^3x-3\cos x}$, $\dfrac{1+\tan^2x}{1-\tan^2x}$, $\dfrac12(\cot x+\tan x)$.

\item $\dfrac{1}{a\sin x+b\cos x}$, $\dfrac{1}{\sin 2x+\cos 2x}$, $\dfrac{a\sin x+b\cos x}{c\sin x+d\cos x}$.

\item Deduce $\displaystyle\int \text{cosec}\,x\,dx=\log\tan\dfrac{x}{2}$ by expressing $\text{cosec}\,x$ as
\[
\dfrac12\left(\cot\dfrac{x}{2}+\tan\dfrac{x}{2}\right).
\]

\item Deduce $\displaystyle\int \sec x\,dx=\log\tan\left(\dfrac{\pi}{4}+\dfrac{x}{2}\right)$ by
\begin{enumerate}[label=(\roman*)]
\item putting $\sin x=z$,
\item putting $\sec x+\tan x=z$.
\end{enumerate}
Show that $\displaystyle\int \sec x\,dx=\cosh^{-1}(\sec x)$.

\item Integrate $\displaystyle\int \dfrac{\cos\theta\,d\theta}{\sin\theta\sqrt{1-\sin^{2n}\theta}}$.

\end{enumerate}

\section*{General Examples}

\begin{enumerate}[leftmargin=2em]

\item If $APB$ be a semicircle, centre $O$, and $PN$ an ordinate to the diameter $AB$, and $P'N'$ a contiguous ordinate, show that
\[
Lt\sum \dfrac{NN'}{NP} = \text{circular measure of the angle } OPN,
\]
the summation being from the centre to any point $N$ between $O$ and $B$, and $NN'$ being indefinitely diminished.

\item Find the area in the first quadrant bounded by the axes of coordinates, the ordinate $x=x_1$ and the curve $x^2=a^2-\dfrac{b^4}{y^2}$.

If the range $x=0$ to $x=a$ on the $x$-axis be divided into $n$ equal portions of length $h$ and rectangles be inscribed in the Newtonian manner, examine the limit of the area of the last of these rectangles when $h$ is indefinitely diminished. Find the whole area from $x=0$ to $x=a$.

\item Find the value of $\displaystyle\int \sqrt{e^{2x}+ae^x}\,dx$. \source{R. P.}

\item Evaluate
\begin{enumerate}[label=(\roman*)]
\item $\displaystyle\int \dfrac{dx}{\sqrt{1-3x-x^2}}$. \source{I. C. S., 1884.}
\item $\displaystyle\int \dfrac{dx}{x\sqrt{x^2+x-6}}$. (Put $x=\dfrac{1}{y}$.) \source{Oxford Second Public Ex., 1880.}
\item $\displaystyle\int \dfrac{(1+x)\,dx}{(2x^2+3x+4)^{\frac32}}$. \source{Colleges $\beta$, 1891.}
\item $\displaystyle\int \dfrac{x-1}{(x^2+2x-1)^{\frac32}}dx$. \source{Trinity, 1892.}
\item $\displaystyle\int \dfrac{3x+4}{(x^2+2x+5)^{\frac32}}dx$. \source{Math. Trip., 1887.}
\end{enumerate}

\item Show that the result of integrating $\displaystyle\int \dfrac{dx}{\sqrt{a^2-x^2}}$ may be exhibited as
\begin{enumerate}[label=(\roman*)]
\item $2\cos^{-1}\dfrac{\sqrt{a+x}-\sqrt{a-x}}{2\sqrt a}$ or as
\item $2\sin^{-1}\dfrac{\sqrt{a+x}+\sqrt{a-x}}{2\sqrt a}$
\end{enumerate}
or as (iii) $2\tan^{-1}\sqrt{\dfrac{a+R}{a-R}}$, where $R=\sqrt{a^2-x^2}$.

\item If $R\equiv ax^2+2bx+c$ and $K=\dfrac{b^2-ac}{a}$, show that
\[
\int \dfrac{dx}{\sqrt R}=\dfrac{1}{\sqrt a}\tanh^{-1}\sqrt{\dfrac{R}{R+K}}
\]
or $=\dfrac{1}{\sqrt{-a}}\tan^{-1}\sqrt{-\dfrac{R}{R+K}}$,
according as $a$ is positive or negative.

\item Evaluate
\begin{enumerate}[label=(\roman*)]
\item $\displaystyle\int \dfrac{dx}{x\sqrt{1+x^6}}$. (Put $x^3=\tan\theta$.) \source{Oxf. I., 1888.}
\item $\displaystyle\int \dfrac{dx}{(x+1)\sqrt{x^2-1}}$. (Put $x=\sec\theta$.) \source{Oxf. I., 1888.}
\end{enumerate}

\item Show that
\begin{enumerate}[label=(\roman*)]
\item $\displaystyle\int_0^1 \dfrac{dx}{(1+x^2)(1-x^2)^{\frac12}}=\dfrac{\pi}{2\sqrt2}$. \source{Trinity, 1888.}
\item $\displaystyle\int \dfrac{dx}{x^4\sqrt{1+x^2}}=\dfrac{2x^2-1}{3x^3}\sqrt{1+x^2}$. \source{Trinity, 1882.}
\end{enumerate}

\item Integrate $\displaystyle\int \sqrt{1+e^x+e^{2x}}\,dx$.

\item A sphere of given radius $a$ consists of an infinite number of concentric shells of very small thickness, the density at the surface of any shell varying as the $n$th power of its radius. Find the mass of the sphere. \source{Ox. I. Pub., 1903.}

If any diameter $AB$ cut one of the shells at $P$ and the density of the shell varies inversely as (i) $OP\sqrt{AP.PB}$, (ii) $OP^2\sqrt{AP.PB}$, find the mass of the sphere in each case, $O$ being the centre.

\item A triangle $ABC$ is divided into strips by lines parallel to $BC$; a point is taken in each strip, and the square of the perpendicular from this point to $BC$ is multiplied by the area of the strip; the same is done with all the strips, and the sum of the products is formed. Express by a definite integral the limiting value of this sum when the breadths of all the strips are diminished indefinitely, and evaluate the integral in terms of the base $BC$ and the distance of $A$ therefrom. \source{Ox. I. Pub., 1901.}

\item Prove that if $u^2=x^2+2px+q$, an integral of the form
\[
\int f(x,u)\,dx
\]
can always be rationalized (provided $f$ is a rational algebraic function) by one of the substitutions
\[
\dfrac{u}{\sqrt{p^2-q}}=\dfrac{2y}{1-y^2} \quad\text{or}\quad \dfrac{u}{\sqrt{q-p^2}}=\dfrac{1+y^2}{1-y^2}. \source{Coll. $a$, 1890.}
\]

\item Find the relation connecting $x$ and $y$, being given
\[
x^2y^4\left(\dfrac{dy}{dx}\right)^2=a^2+b^2y^3. \source{I. C. S., 1889.}
\]

\item Show that $\displaystyle\int_{\sqrt3}^{2\sqrt3} z^3(z^2-3)^{\frac32}dz=\dfrac{16038}{35}$. \source{Ox. I., 1888.}

\item Integrate
\begin{enumerate}[label=(\roman*)]
\item $\displaystyle\int \dfrac{\sin\theta\,d\theta}{\sqrt{a\cos^2\theta+2b\cos\theta+c}}$,
\item $\displaystyle\int \dfrac{\cos\theta\,d\theta}{\sqrt{a\sin^2\theta+2b\sin\theta+c}}$,
\item $\displaystyle\int \dfrac{d\theta}{\cos\theta\sqrt{a\cos^2\theta+2b\sin\theta\cos\theta+c\sin^2\theta}}$,
\item $\displaystyle\int \dfrac{d\theta}{\sin\theta\sqrt{a\cos^2\theta+2b\sin\theta\cos\theta+c\sin^2\theta}}$,
\item $\displaystyle\int \dfrac{d\theta}{\sin\theta\sqrt{a\cos^2\theta+b\sin^2\theta+c}}$.
\end{enumerate}
\source{Trin., 1888.}

\item Integrate
\begin{enumerate}[label=(\roman*)]
\item $\displaystyle\int \dfrac{x^4}{(a^2-x^2)^{\frac32}}dx$,
\item $\displaystyle\int \dfrac{x^2\,dx}{(a+bx^2)\sqrt{c^2-x^2}}$.
\end{enumerate}
\source{Trin., 1888.}

\item
\begin{enumerate}[label=(\alph*)]
\item Evaluate $\displaystyle\int_1^7 (x^2-6x+13)\,dx$, first directly, second by putting $x^2-6x+13=y$. (Draw a graph and explain fully.)
\item Evaluate $\displaystyle\int_0^{\frac{2b}{a}} (ax^2-2bx+c)\,dx$,
and explain by a graph the result when $2b^2=3ac$.
Obtain the same result by substituting
$ax^2-2bx+c=y$, taking $b^2<ac$.
Also obtain $\displaystyle\int_0^{\frac{3b}{a}} (ax^2-2bx+c)\,dx$ by this substitution, explaining your limits for $y$ by means of the graph.
\end{enumerate}

\item Point out the fallacy in the following argument:
\[
\int_{-1}^1 \dfrac{dx}{1+x^2}=\left[\tan^{-1}x\right]_{-1}^1=\dfrac{\pi}{4}-\left(-\dfrac{\pi}{4}\right)=\dfrac{\pi}{2}.
\]
But putting $x=\dfrac{1}{y}$, $dx=-\dfrac{dy}{y^2}$. When $x=-1$, $y=-1$. When $x=1$, $y=1$.
\[
\therefore \int_{-1}^1 \dfrac{dx}{1+x^2}=-\int_{-1}^1 \dfrac{dy}{1+y^2}=-\int_{-1}^1 \dfrac{dx}{1+x^2},
\]
for, as the result is numerical, the letter used in integration cannot affect the result.

Hence $2\displaystyle\int_{-1}^1 \dfrac{dx}{1+x^2}=0$; but $\displaystyle\int_{-1}^1 \dfrac{dx}{1+x^2}=\dfrac{\pi}{2}$; $\therefore \pi=0$!

\item Point out the fallacy in the following reasoning:

We have, if we put $x=e^t$,
\[
\left(x\dfrac{d}{dx}\right)^a x^k=\left(\dfrac{d}{dt}\right)^a e^{kt}=k^a e^{kt}=k^a x^k.
\]
But when $a=-1$, we have
\[
\left(x\dfrac{d}{dx}\right)^{-1}x^k=\dfrac1x\int x^k dx=\dfrac{x^k}{k+1},
\]
and these two results do not agree. \source{R.P.}

\item Prove that $\iota\,\text{gd}\left(\dfrac1\iota \text{gd}\,u\right)=u$, \source{Cayley, E.F.}

and show that if $\text{gd}\,u=a_1u+a_3u^3+a_5u^5+\dots$,
then will $\text{gd}^{-1}u=a_1u-a_3u^3+a_5u^5-\dots$.

\item If $\sec x+\tan x=1+S_1\dfrac{x}{1!}+S_2\dfrac{x^2}{2!}+S_3\dfrac{x^3}{3!}+\dots$,

show that $\text{gd}^{-1}x=x+S_3\dfrac{x^3}{3!}+S_5\dfrac{x^5}{5!}+\dots$,

that $S_{p+1}=S_p+\dbinom{p}{2}S_{p-2}S_2+\dbinom{p}{4}S_{p-4}S_4+\dots$

and $S_n-\dbinom{n}{2}S_{n-2}+\dbinom{n}{4}S_{n-4}-\dots+\cos\dfrac{n\pi}{2}=\sin\dfrac{n\pi}{2}$,

and that $S_3=2$, $S_5=16$, $S_7=272$, $S_9=7936$, etc. \source{Diff. Calc., Art. 573, etc.}

\item Integrate $\displaystyle\int \dfrac{dx}{(a+x)(c+x)^{\frac12}}$
by putting $c+x=(a-c)z^2$ or $(c-a)z^2$,
according as $a>$ or $<c$.

Taking the case $a>c$, consider the same integral with $a+da$ replacing $a$, subtract the original integral, divide by $da$, and take the limit when $da$ is indefinitely diminished.

Hence obtain $\displaystyle\int \dfrac{dx}{(a+x)^2(c+x)^{\frac12}}$.

Deduce also $\displaystyle\int \dfrac{dx}{(a+x)(c+x)^{\frac32}}$.

\item Evaluate $\displaystyle\int \dfrac{dx}{(x^2-a^2)(x^2-c^2)^{\frac12}}$.
\begin{enumerate}[label=(\roman*)]
\item if $a>c$,
\item if $a<c$. (Put $x=c\sec\phi$).
\end{enumerate}
\source{Math. Trip., 1878.}

\item Show that $\displaystyle\int \dfrac{(x-p)^{2n+1}}{(ax^2+2bx+c)^{n+\frac32}}dx=\dfrac{1}{(b^2-ac)^{n+1}}\int (y^2-q^2)^n dy$,

where $q^2\equiv ap^2+2bp+c$

and $y\sqrt{ax^2+2bx+c}=(ap+b)x+bp+c$. \source{Colleges, 1901.}

\item If $F(x)=a_1f(x)+a_2f(2x)+a_3f(3x)+\dots$,

prove that $\dfrac{a_1}{1^x}+\dfrac{a_2}{2^x}+\dfrac{a_3}{3^x}+\dots=\dfrac{\int_0^\infty t^{x-1}F(t)\,dt}{\int_0^\infty t^{x-1}f(t)\,dt}$. \source{Ox. J. M. Sch., 1904.}

\item Integrate
\begin{enumerate}[label=(\roman*)]
\item $\displaystyle\int \dfrac{1+x^2}{1-x^2}\dfrac{dx}{\sqrt{1+x^4}}$.
\item $\displaystyle\int \dfrac{1-x^2}{1+x^2}\dfrac{dx}{\sqrt{1+x^4}}$.
\end{enumerate}
\source{Euler.}

\item Show that if $F(x,y)$ be a rational function of $x$ and $y$,
\[
\int F\left(x,\sqrt[n]{\dfrac{\alpha x+\beta}{\gamma x+\delta}}\right)dx
\]
can be thrown into rational form by the substitution
\[
\dfrac{\alpha x+\beta}{\gamma x+\delta}=z^n.
\]
Hence show that
\[
\int \left(\dfrac{1-2x}{1+2x}\right)^{\frac32}dx=3\tan^{-1}\sqrt{\dfrac{1-2x}{1+2x}}-2\sqrt{\dfrac{1-2x}{1+2x}}-\dfrac{\sqrt{1-4x^2}}{2}.
\]

\item Show that if $F(x,y)$ be any rational integral function of $x$ and $y$,
\[
\int F(x,\sqrt{ax^2+2bx+c})\,dx
\]
can be thrown into rational form by any of the substitutions
\begin{enumerate}[label=(\arabic*)]
\item $\sqrt{ax^2+2bx+c}=\sqrt a(x+z)$,
\item $\sqrt{ax^2+2bx+c}=xz+\sqrt c$,
\item $x-x_2=z^2(x-x_1)$,
\end{enumerate}
where $x_1$, $x_2$ are the roots of $ax^2+2bx+c=0$. \source{Bertrand, C.I., p. 39.}

Apply each of these methods to the integration of
\[
\int \dfrac{x\,dx}{\sqrt{x^2-6x+8}},
\]
showing that the result in each case reduces to
\[
\sqrt{x^2-6x+8}+3\cosh^{-1}(x-3),
\]
as derived by the method of Art. 85.

\item If $x^3-3a^2x=a^2y$,

show that $\dfrac{dx}{\sqrt{x^2-4a^2}}=\dfrac13\dfrac{dy}{\sqrt{y^2-4a^2}}$,

and hence obtain Cardan's formula for the solution of a cubic. \source{J. M. Sch. Ox.}

\item Evaluate $\displaystyle\int_0^{\pi/2} \dfrac{\cos x\,dx}{1-\sin^2a\cos^2x}$,

and deduce the expansion
\[
\dfrac{2a}{\sin 2a}=1+\dfrac23\sin^2a+\dfrac{2.4}{3.5}\sin^4a+\dots, \quad\text{where } \dfrac{\pi}{2}>a>0. \source{Oxf. I. P., 1915.}
\]

\item Show that
\[
\int \dfrac{dx}{(1+x^4)\{(1+x^4)^{\frac12}-x^2\}^{\frac12}}=\tan^{-1}\dfrac{x}{\{(1+x^4)^{\frac12}-x^2\}^{\frac12}}. \source{Euler, C.I., iv.}
\]
Integrate $\displaystyle\int \dfrac{dx}{(1+x^{2n})\{(1+x^{2n})^{\frac1n}-x^2\}^{\frac12}}$.

\item Show that the integrals
\[
\int \dfrac{dx}{(1-x^m)(2x^m-1)^{\frac{1}{2m}}} \quad\text{and}\quad \int \dfrac{x^{m-1}dx}{(1-x^m)(2x^m-1)^{\frac{1}{2m}}}
\]
are reduced to the integration of rational fractions by the respective substitutions $2x^m-1=u^{2m}x^{2m}$ and $2x^m-1=u^{2m}$. \source{Lexell, Actes de P\'etersbourg, 1781, ii.; Lacroix, C.D., ii., p. 65.}

\end{enumerate}

\newpage

\begin{center}
    {\LARGE \bfseries \color{accent} Chapter IV: Integration by Parts. Powers of Sines and Cosines --- Exercises}\\
    \vspace{0.2cm}
    {\color{accent}\rule{0.55\linewidth}{0.8pt}}
\end{center}
\vspace{0.35cm}

\section*{Exercises I --- Integration by Parts (Arts. 90--100)}

Integrate by parts

\begin{enumerate}[leftmargin=2em]

\item $xe^{3x}$, $x^2e^{ax}$, $x^5e^{-x}$, $x\cosh x$, $x^2\sinh x$.

\item $x\cos x$, $x^5\cos 2x$, $x^2\cos^2 x$, $x^2\cos 3x\sin x$, $x\sin x\sin 2x\sin 3x$.

\item $e^x\sin 2x$, $e^x\sin^2 x$, $e^{3x}\sin^3 x\cos x$, $e^{-5x}\cos x\sin^2 x\cos 3x$.

\item $x^3\log x$, $x^n\log x$, $x^n(\log x)^2$, $x^n(\log x)^3$.

\item $e^{ax}\sin px\sin qx\sin rx$, $e^{ax}\sin px\sin qx\cos rx$.

\item $e^{ax}\sin px\sin qx\cos^2 rx$, $e^{ax}\cos px\cos qx\cos^2(p+q)x$.

\item Evaluate
\[
\int_0^\pi x\sin x\,dx, \qquad \int_0^{\pi/2} x^2\cos x\,dx, \qquad \int_0^{\pi/2} x^2\cos 2x\,dx.
\]

\item Integrate
\[
\int \sin^{-1}x\,dx, \qquad \int x\sin^{-1}x\,dx, \qquad \int x^3\sin^{-1}x\,dx, \qquad \int x\tan^{-1}x\,dx.
\]

\end{enumerate}

\section*{Exercises II --- Reduction Formulae for $\int x^m\sin nx\,dx$, $\int x^m\cos nx\,dx$ (Arts. 101--102)}

Write down the integrals of

\begin{enumerate}[leftmargin=2em]

\item $\displaystyle\int x^6e^x\,dx$, $\displaystyle\int x^5\sinh x\,dx$, $\displaystyle\int x^5\cosh^2 x\,dx$.

\item $\displaystyle\int_0^{\pi/2} x^3\sin x\,dx$, $\displaystyle\int_0^{\pi/2} x^3\sin^2 x\,dx$, $\displaystyle\int_0^{\pi/2} x^4\sin x\cos x\,dx$.

\item $\displaystyle\int_0^\pi x^5\sin x\,dx$, $\displaystyle\int_0^\pi x^5\cos^2 x\,dx$, $\displaystyle\int_0^1 x^3\cosh x\,dx$.

\item $\displaystyle\int_0^{\pi/2} x^3(a^2\cos^2 x+b^2\sin^2 x)\,dx$, $\displaystyle\int_1^3 x^3\log x\,dx$, $\displaystyle\int_0^1 x\tan^{-1}x\,dx$.

\item $\displaystyle\int_0^{\pi/2} e^x\sin x\cos^2 x\,dx$, $\displaystyle\int_0^{\pi/2} x\sin x\sin 2x\sin 3x\,dx$.

\end{enumerate}

\section*{Exercises III --- Powers and Products of Sines, Cosines, Secants, Cosecants, Tangents and Cotangents (Arts. 103--126)}

\begin{enumerate}[leftmargin=2em]

\item Integrate $\sin^2x$, $\sin^3x$, $\sin^4x$, $\sin^5x$, $\sin^8x$, $\sin^9x$, $\sin^{2n}x$, $\sin^{2n+1}x$, doing those with odd indices in two ways.

\item Integrate $\sin^2x\cos^2x$, $\sin^3x\cos^3x$, $\sin^4x\cos^4x$, $\sin^3x\cos^6x$, $\sin^6x\cos^3x$, $\sin^6x\cos^4x$.

\item Integrate $\dfrac{\sin^2x}{\cos^4x}$, $\cos^2x\,\text{cosec}^4x$, $\sec^2x\,\text{cosec}^2x$, $\dfrac{1}{\sin^4x\cos^4x}$.

\item Evaluate $\displaystyle\int_0^{\pi/4}\sin^2x\,dx$, $\displaystyle\int_0^{\pi/4}\cos^5x\,dx$, $\displaystyle\int_0^{\pi/4}\cos^6x\,dx$.

\item Integrate $\sin ax\cos^2bx$, $\sin 3x\cos^3x$, $\sin nx\cos^2x$.

\item Show that
\[
\int \sin x\sin 2x\sin 3x\,dx = -\dfrac{1}{8}\cos 2x-\dfrac{1}{16}\cos 4x+\dfrac{1}{24}\cos 6x.
\]

\item Show that
\begin{enumerate}[label=(\roman*)]
\item $\displaystyle\int \sin mx\cos nx\,dx = -\dfrac{\cos(m+n)x}{2(m+n)}-\dfrac{\cos(m-n)x}{2(m-n)}$.
\item $\displaystyle\int \sin mx\sin nx\,dx = \dfrac{\sin(m-n)x}{2(m-n)}-\dfrac{\sin(m+n)x}{2(m+n)}$.
\item $\displaystyle\int \cos mx\cos nx\,dx = \dfrac{\sin(m-n)x}{2(m-n)}+\dfrac{\sin(m+n)x}{2(m+n)}$.
\end{enumerate}
Deduce from (ii) and (iii) the values of
\[
\int \sin^2mx\,dx \quad\text{and}\quad \int \cos^2mx\,dx,
\]
and verify the results by independent integration.

\item Prove that $\displaystyle\int_0^\pi \sin mx\sin nx\,dx$ and $\displaystyle\int_0^\pi \cos mx\cos nx\,dx$ are both zero so long as $m$ and $n$ are integral and unequal. But if $m$ and $n$ are equal integers their values are each equal to $\dfrac{\pi}{2}$.

\end{enumerate}

\section*{General Examples}

\begin{enumerate}[leftmargin=2em]

\item Prove that $\displaystyle\int u\dfrac{d^2v}{dx^2}dx = u\dfrac{dv}{dx}-v\dfrac{du}{dx}+\int v\dfrac{d^2u}{dx^2}dx$.

\item Perform the following integrations:
\begin{enumerate}[label=(\roman*)]
\item $\displaystyle\int \cos^{-1}x\,dx$.
\item $\displaystyle\int \cos^{-1}\dfrac{1}{x}\,dx$.
\item $\displaystyle\int x^3\tan^{-1}x\,dx$.
\item $\displaystyle\int x\sec^2x\,dx$.
\item $\displaystyle\int x\sec x\tan x\,dx$.
\item $\displaystyle\int (ax+b)\log(cx+d)\,dx$.
\item $\displaystyle\int \tan^{-1}\sqrt{1-x^2}\,dx$. \source{St. John's, 1886.}
\item $\displaystyle\int \dfrac{x^4}{1+x^2}\tan^{-1}x\,dx$. \source{Ox. II. P., 1889.}
\end{enumerate}

\item Integrate
\begin{enumerate}[label=(\roman*)]
\item $\displaystyle\int e^{a\sin^{-1}x}\,dx$.
\item $\displaystyle\int \dfrac{x\sin^{-1}x}{(1-x^2)^{\frac12}}\,dx$.
\item $\displaystyle\int \dfrac{x^3\sin^{-1}x}{(1-x^2)^{\frac32}}\,dx$.
\end{enumerate}

\item Integrate
\begin{enumerate}[label=(\roman*)]
\item $\displaystyle\int \dfrac{e^{m\tan^{-1}x}}{1+x^2}\,dx$.
\item $\displaystyle\int \dfrac{e^{m\tan^{-1}x}}{(1+x^2)^{\frac32}}\,dx$.
\item $\displaystyle\int \dfrac{e^{m\tan^{-1}x}}{(1+x^2)^2}\,dx$.
\item $\displaystyle\int \dfrac{e^{m\tan^{-1}x}}{(1+x^2)^{\frac52}}\,dx$.
\item $\displaystyle\int \dfrac{e^{m\tan^{-1}x}}{(1+x^2)^{\frac{n}{2}+1}}\,dx$ ($n\equiv$ a positive integer).
\end{enumerate}

\item Integrate
\begin{enumerate}[label=(\roman*)]
\item $\displaystyle\int xe^{bx}\cos ax\,dx$.
\item $\displaystyle\int x^2e^{ax}\sin bx\,dx$. \source{a, 1888.}
\item $\displaystyle\int xe^x\sin^2x\,dx$.
\end{enumerate}

\item Integrate
\begin{enumerate}[label=(\roman*)]
\item $\displaystyle\int e^{ax}(\sin bx+\cos bx)\,dx$.
\item $\displaystyle\int e^{ax}(\sinh bx+\cosh bx)\,dx$.
\item $\displaystyle\int e^{ax}\sinh bx\cosh ax\,dx$.
\item $\displaystyle\int e^{ax}\cosh ax\sin bx\,dx$.
\item $\displaystyle\int x^2 3^x\sin 4x\,dx$.
\item $\displaystyle\int \cos\left(b\log\dfrac{x}{a}\right)dx$.
\item $\displaystyle\int \cosh\left(b\log\dfrac{x}{a}\right)dx$.
\item $\displaystyle\int_0^\pi \theta\sin\theta\cosh(\cos\theta)\,d\theta$. \source{a, 1891.}
\end{enumerate}

\item Integrate
\begin{enumerate}[label=(\roman*)]
\item $\displaystyle\int \dfrac{xe^x}{(x+1)^2}\,dx$.
\item $\displaystyle\int e^x\dfrac{1+\sin x}{1+\cos x}\,dx$.
\item $\displaystyle\int e^x\dfrac{1-\sin x}{1-\cos x}\,dx$.
\item $\displaystyle\int \dfrac{\cosh x+\sinh x\sin x}{1+\cos x}\,dx$.
\item $\displaystyle\int \dfrac{dx}{1+e^x}$. \source{Mech. Sc. Trip.}
\item $\displaystyle\int \sqrt{1+e^{nx}}\,dx$.
\item $\displaystyle\int e^x\dfrac{1+x^2}{(1+x)^2}\,dx$. \source{Ox. I. P., 1890.}
\end{enumerate}

\item Integrate
\begin{enumerate}[label=(\roman*)]
\item $\displaystyle\int (\log x)^2\,dx$. \source{Ox. I. P., 1888.}
\item $\displaystyle\int \left(x+\dfrac{1}{x}+\dfrac{1}{x^2}\right)\log x\,dx$. \source{Ox. I. P., 1889.}
\item $\displaystyle\int x^{-2}\tan^{-1}x\,dx$. \source{Ox. II. P., 1887.}
\item $\displaystyle\int \log\left(x+\sqrt{a^2+x^2}\right)dx$. \source{Math. Trip., 1882.}
\item $\displaystyle\int x\log\left(x+\sqrt{a^2+x^2}\right)dx$. \source{St. John's, 1884.}
\item $\displaystyle\int (a+x)\sqrt{a^2+x^2}\,dx$. \source{St. John's, 1888.}
\item $\displaystyle\int (a^2+x^2)\sqrt{a+x}\,dx$. \source{St. John's, 1888.}
\item $\displaystyle\int e^{ax}x^2\sin(bx+c)\,dx$. \source{Coll., 1892.}
\item $\displaystyle\int x^3(1-x^{\frac23})^{\frac32}\,dx$. \source{Ox. I. P., 1890.}
\end{enumerate}

\item Integrate
\begin{enumerate}[label=(\roman*)]
\item $\displaystyle\int xe^{ax}\sin bx\sin cx\,dx$.
\item $\displaystyle\int xe^{ax}\sin bx\sin^2cx\,dx$.
\end{enumerate}

\item Show that if $u$ be a rational integral function of $x$,
\[
\int e^{x/a}u\,dx = ae^{x/a}\left\{u-a\dfrac{du}{dx}+a^2\dfrac{d^2u}{dx^2}-a^3\dfrac{d^3u}{dx^3}+\dots\right\},
\]
where the series within the brackets is necessarily finite. \source{Trin. Coll., 1881.}

\item If $u\equiv\displaystyle\int e^{ax}\cos bx\,dx$, $v\equiv\displaystyle\int e^{ax}\sin bx\,dx$, prove that
\[
\tan^{-1}\dfrac{v}{u}+\tan^{-1}\dfrac{b}{a}=bx,
\]
and that $(a^2+b^2)(u^2+v^2)=e^{2ax}$.

\item Evaluate $\displaystyle\int x^2\log(1-x^2)\,dx$, and deduce that
\[
\dfrac{1}{1.5}+\dfrac{1}{2.7}+\dfrac{1}{3.9}+\dots = \dfrac{8}{9}-\dfrac{2}{3}\log_e 2. \source{a, 1889.}
\]

\item Integrate $\displaystyle\int \sec^{\frac53}\theta\,\text{cosec}^{\frac53}\theta\,d\theta$, $\displaystyle\int \sec^{\frac53}\theta\sin\theta\,d\theta$.

\item Find the value of
\[
\int \left\{u\dfrac{d^nv}{dx^n}-(-1)^nv\dfrac{d^nu}{dx^n}\right\}dx. \source{$\gamma$, 1890.}
\]

\item Evaluate
\[
\int\left[\dfrac{d^3u}{dx^3}\left(\dfrac{dv}{dx}-\dfrac{dw}{dx}\right)+\dfrac{d^3v}{dx^3}\left(\dfrac{dw}{dx}-\dfrac{du}{dx}\right)+\dfrac{d^3w}{dx^3}\left(\dfrac{du}{dx}-\dfrac{dv}{dx}\right)\right]dx. \source{$\gamma$, 1890.}
\]

\item Establish the following formulae for integration by parts, $u$ and $v$ being functions of $x$, and accents denoting differentiations and suffixes integrations with respect to $x$, and $u^{(n)}$ denoting $u$ with $n$ accents:
\begin{enumerate}[label=(\roman*)]
\item $\displaystyle\int uv\,dx = uv_1-u'v_2+u''v_3-u'''v_4+\dots+(-1)^{n-1}u^{(n-1)}v_n$
\[
+(-1)^n\int u^{(n)}\,dv_{n+1}.
\]
\item $\displaystyle\iint(uv)(dx)^2 = uv_2-2u'v_3+3u''v_4-4u'''v_5+\dots+(-1)^{n-1}nu^{(n-1)}v_{n+1}$
\[
+(-1)^n n\int u^{(n)}v_{n+1}\,dx+(-1)^n\int dx\int u^{(n)}v_n\,dx. \source{a, 1888.}
\]
\end{enumerate}

\item If $u$ be a function of $x$, and differentiations and integrations are respectively denoted by accents and suffixes, and $(n)$ means $n$ accents, show that
\[
\log u = 1-\left(\dfrac{1}{u}\right)'u_1+\left(\dfrac{1}{u}\right)''u_2-\left(\dfrac{1}{u}\right)'''u_3+\dots+(-1)^n\int\left(\dfrac{1}{u}\right)^{(n)}du_n. \source{St. John's, 1889.}
\]

\item If $u,v,w$ be functions of $x$, and accents and suffixes denote differentiations and integrations respectively, show that
\[
2uvw = (vw)'u_1-(vw)''u_2+(vw)'''u_3-\dots+(-1)^{m-1}\int(vw)^{(m)}du_m
\]
\[
+(wu)'v_1-(wu)''v_2+(wu)'''v_3-\dots+(-1)^{n-1}\int(wu)^{(n)}dv_n
\]
\[
+(uv)'w_1-(uv)''w_2+(uv)'''w_3-\dots+(-1)^{p-1}\int(uv)^{(p)}dw_p. \source{St. John's, 1889.}
\]

\item Prove that
\[
\int_0^1 v^{xv}\,dv = 1-\dfrac{x}{2^2}+\dfrac{x^2}{3^3}-\dfrac{x^3}{4^4}+\dfrac{x^4}{5^5}-\dots\ \text{etc.} \source{Math. Trip., 1878.}
\]

\item Find the value of $\displaystyle\int_0^1 x^x\,dx$ correct to five decimal places. \source{J. M. Sch. Ox., 1904.}

\item Prove that
\[
e^{a^2x^2}\int_0^x e^{-a^2x^2}\,dx = x+\dfrac{2}{3}a^2x^3+\dfrac{2^2}{3.5}a^4x^5+\dots
\]
\[
+\dfrac{(2a^2)^n}{3.5\dots(2n-1)}e^{a^2x^2}\int_0^x x^{2n}e^{-a^2x^2}\,dx. \source{Ox. I. Pub., 1899.}
\]

\item Find the sum of the series, supposed convergent,
\[
\dfrac{x^5}{1.3.5}-\dfrac{x^7}{3.5.7}+\dfrac{x^9}{5.7.9}-\text{etc. to }\infty. \source{Coll., 1881.}
\]

\item If $y$ and $z$ be functions of $x$, and $u=yz'-zy'$, prove the following:
\begin{enumerate}[label=(\roman*)]
\item $\displaystyle\int zu^{-2}(y'z''-z'y'')\,dx = -y^{-1}(1+y'zu^{-1})$,
\item the integration of $zy^{-1}u^{-2}(yz''-zy'')$ can be reduced to that of $y^{-2}$.
\end{enumerate}
\source{St. John's, 1886.}

\item Show how the method of integration by parts may be applied to find
\[
\int f(x)\dfrac{d^{n+1}V}{dx^{n+1}}\,dx,
\]
where $f(x)$ is a rational algebraical expression of the $n$th degree.

Prove that $\displaystyle\int_{-1}^1 f(x)\dfrac{d^{n+1}(x^2-1)^{n+1}}{dx^{n+1}}\,dx = 0$. \source{Coll., 1876.}

\item Prove that $\displaystyle\int (\cos x)^n\,dx$ may be expressed by the series,
\[
\sin x - N_1\dfrac{\sin^3x}{3}+N_2\dfrac{\sin^5x}{5}-N_3\dfrac{\sin^7x}{7}+\dots,\ \text{etc.},
\]
$N_1$, $N_2$, $N_3$, $\dots$ being the coefficients of the expansion $(1+a)^{\frac{n-2}{2}}$, and $n$ having any real value positive or negative. \source{Smith's Prize, 1876.}

\item Prove that
\[
\int x^ne^x\sin x\,dx = e^x\sum_{r=0}^{r=n}(-1)^r\dfrac{n!}{(n-r)!}x^{n-r}2^{-\frac{r+1}{2}}\sin\left\{x-\dfrac{(r+1)\pi}{4}\right\}.
\]

\item Express the infinite series
\[
\dfrac{1}{2}+\dfrac{1.3}{2.4}\cdot\dfrac{1}{2}+\dfrac{1.3.5}{2.4.6}\cdot\dfrac{1}{3}+\dots
\]
as a definite integral, and find its value. \source{Ox. II. P., 1902.}

\item Show that
\[
2^m\int \cos mx\cos^mx\,dx = A+x+m\dfrac{\sin 2x}{2}+\dfrac{m(m-1)}{1.2}\dfrac{\sin 4x}{4}+\dots+\dfrac{\sin 2mx}{2m},
\]
where $m$ is an integer and $A$ is independent of $x$. \source{Coll. a, 1885.}

\item Evaluate the integral
\[
\int_0^T a_1\sin\dfrac{2\pi t}{T}\cdot a_2\sin\dfrac{2\pi}{T}(t+\lambda)\,dt,
\]
and draw curves showing how its value depends on that of $\lambda$. \source{Mech. Sc. Trip., 1899.}

\item Prove that if $y=f(x)$ and $x=\phi(y)$ are equivalent relations, then, between any corresponding limits,
\[
\int\sqrt{f'(x)}\,dx = \int\sqrt{\phi'(y)}\,dy.
\]
Hence, or otherwise, prove that if $\tan\beta=\sqrt{1-c}\tan\alpha$,
\[
\int_0^\alpha \dfrac{dx}{\sqrt{1-c\sin^2x}} = \int_0^\beta \dfrac{dx}{\sqrt{1-c\cos^2x}}. \source{Ox. II. P., 1886.}
\]

\item Prove that the remainder $R$ in the series
\[
\theta = \tan\theta-\dfrac13\tan^3\theta+\dots+\dfrac{(-1)^n}{2n+1}\tan^{2n+1}\theta+R
\]
may be written as a definite integral,
\[
(-1)^{n+1}\int_0^\theta \tan^{2n+2}\theta\,d\theta. \source{Coll., 1881.}
\]

\item Show that the integrals $\displaystyle\int_0^x f(z)\,dz$, $\displaystyle\int_0^x z^nf^{(n)}(z)\,dz$ are connected thus:
\[
\int_0^x f(z)\,dz = xf(x)-\dfrac{x^2}{2!}f'(x)+\dfrac{x^3}{3!}f''(x)-\dots
\]
\[
+(-1)^{n-1}\dfrac{x^n}{n!}f^{(n-1)}(x)+(-1)^n\int_0^x z^nf^{(n)}(z)\,dz,
\]
and that if one can be integrated the other can also be integrated. \source{Bernoulli.}

\item Integrate
\[
\int\{(2n+1)\cos(2n+\tfrac12)\theta+(2n-2)\cos(2n-\tfrac32)\theta\}(\cos\theta)^{\frac12}\,d\theta,
\]
and prove that when $n$ is a positive integer,
\[
\int_0^{\pi/2}\cos(2n+\tfrac12)\theta\,(\cos\theta)^{\frac12}\,d\theta = 0. \source{Oxford II. Pub., 1913.}
\]

\item Find the sum of the areas included between the axis of $x$ and the arc of the curve $y=x\sin(x/a)$ from the ordinate $x=0$ to the ordinate $x=n\pi a$, $n$ being any positive integer, odd or even. \source{Oxf. I. P., 1911.}

\item Evaluate $\displaystyle\int_0^{2a}\dfrac{x^n}{\sqrt{2ax-x^2}}\,dx$ when $n$ is any positive integer. \source{Oxf. I. P., 1916.}

\item Show that $\displaystyle\int_0^1 x\log\left(1+\tfrac12x\right)dx = \tfrac34\left(1-2\log\tfrac32\right)$, and prove that this is less than $\displaystyle\int_0^1 \tfrac12x^2\,dx$. \source{Math. Trip., Part I., 1913.}

\item If $T_n=\displaystyle\int_0^x \tan^nx\,dx$, show that $(n-1)(T_n+T_{n-2})=\tan^{n-1}x$.

Given that $\pi=3.141592\dots$, $\log_e2=0.693147\dots$, show that
\[
\int_0^{\pi/4}\tan^5x\,dx = 0.09657\dots,\qquad \int_0^{\pi/4}\tan^4x\,dx = 0.11873\dots. \source{Math. Trip. I., 1915.}
\]

\item Prove that $\displaystyle\int_0^{1/\sqrt2} \dfrac{\sin^{-1}x}{(1-x^2)^{\frac32}}\,dx = \dfrac{\pi}{4}-\dfrac12\log_e2$. \source{Math Trip. I., 1917.}

\item Find the area $A$ between the curve
\[
y=a\left(\sin x+\dfrac13\sin 3x+\dfrac15\sin 5x\right)
\]
and the axis of $x$ between the limits $0$ and $\pi$; and the volume $V$ obtained by rotating this area about the axis of $x$.

Prove that $4V=\pi^2aA$. \source{Math. Trip. I., 1913.}

\item Show that
\[
\int_0^1 x^{2p-1}\log(1+x)\,dx = \dfrac{1}{2p}\left\{\dfrac{1}{1.2}+\dfrac{1}{3.4}+\dots+\dfrac{1}{(2p-1)2p}\right\}. \source{Math. Trip., Pt. I., 1916.}
\]

\end{enumerate}

\newpage

\begin{center}
    {\LARGE \bfseries \color{accent} Chapter V: Rational Algebraic Fractional Forms --- Exercises}\\
    \vspace{0.2cm}
    {\color{accent}\rule{0.55\linewidth}{0.8pt}}
\end{center}
\vspace{0.35cm}

\section*{Exercises --- Partial Fractions and Rational Algebraic Fractional Forms (Arts. 127--169)}

Integrate with regard to $x$ the expressions in the following seven groups:

\begin{enumerate}[leftmargin=2em]

\item Linear unrepeated factors
\begin{enumerate}[label=(\roman*)]
\item $\dfrac{1}{x(x^2-1)}$.
\item $\dfrac{1}{(x-1)(x-3)(x-5)}$.
\item $\dfrac{1-3x^2}{3x-x^3}$.
\item $\dfrac{x^2+x+1}{(x^2-1)(x^2-4)}$.
\item $\dfrac{(x-1)(x-2)}{(x^2-9)(x^2-16)}$.
\item $\dfrac{(x-a)(x-b)(x-c)}{(x-a_1)(x-b_1)(x-c_1)}$.
\item $\dfrac{(x-a)(x-b)(x-c)}{(x^2-a_1^2)(x^2-b_1^2)(x^2-c_1^2)}$.
\item $\dfrac{x+1}{x^2+10x-75}$.
\item $\dfrac{x+1}{x^2+10x-119}$.
\item $\dfrac{x+1}{x^3-31x^2+311x-1001}$.
\end{enumerate}

\item Linear repeated factors:
\begin{enumerate}[label=(\roman*)]
\item $\dfrac{1}{(x-1)^3(x+1)}$.
\item $\dfrac{1}{(x-1)^4(x+1)^4}$.
\item $\dfrac{x+1}{x^4(x-1)^4}$.
\item $(ax^2+bx^3)^{-1}$.
\item $(x^2-7x+12)^{-2}$.
\item $\dfrac{x^3}{(x-a)^2(x-b)}$.
\item $\dfrac{x^2-3x+3}{x^3-7x^2+16x-12}$. \source{I. C. S., 1900.}
\end{enumerate}

\item Quasi-linear occurrence of factors. Powers of $x$ all even:
\begin{enumerate}[label=(\roman*)]
\item $\displaystyle\int \dfrac{dx}{(x^2+a^2)(x^2+b^2)}$.
\item $\displaystyle\int \dfrac{(x^2+a^2)(x^2+b^2)}{(x^2+c^2)(x^2+d^2)}dx$.
\item $\displaystyle\int \dfrac{x^2(x^2+a^2)}{(x^2+c^2)}dx$.
\item $\displaystyle\int \dfrac{x^2\,dx}{(x^2+1)(2x^2+1)}$.
\item $\displaystyle\int \dfrac{ax^2+b}{(cx^2+d)(ex^2+f)}dx$.
\item $\displaystyle\int \dfrac{ax^2+b}{x^2(cx^2+d)(ex^2+f)(gx^2+h)}dx$.
\end{enumerate}
In the last two $c$, $d$, $e$, $f$, $g$, $h$ may be considered positive.

\item Quasi-linear factors. Numerator an odd function, Denominator even, or Numerator even, Denominator odd:
\begin{enumerate}[label=(\roman*)]
\item $\displaystyle\int \dfrac{dx}{x(x^2+1)}$.
\item $\displaystyle\int \dfrac{x^2+2}{x(x^4-1)}dx$.
\item $\displaystyle\int \dfrac{dx}{x^7-6x^5+11x^3-6x}$.
\item $\displaystyle\int \dfrac{x\,dx}{(ax^2+bx+c)^2+(ax^2-bx+c)^2}$.
\end{enumerate}

\item Quadratic factors not repeated:
\begin{enumerate}[label=(\roman*)]
\item $\displaystyle\int \dfrac{dx}{x^4+x^2+1}$.
\item $\displaystyle\int \dfrac{(x+1)^2}{x^4+x^2+1}dx$.
\item $\displaystyle\int \dfrac{x^2+1}{x^4+1}dx$.
\item $\displaystyle\int \dfrac{x^2+1}{x^4-x^2+1}dx$.
\item $\displaystyle\int (x^2+a^2)(x^4+a^2x^2+a^4)^{-1}dx$.
\item $\displaystyle\int (x^2-a^2)(x^4+a^2x^2+a^4)^{-1}dx$.
\item $\displaystyle\int \dfrac{x^2+3x+1}{x^4+x^2+1}dx$.
\item $\displaystyle\int \dfrac{dx}{x^4+1}$.
\end{enumerate}

\item Linear factors repeated. Quadratic factors not repeated.
\begin{enumerate}[label=(\roman*)]
\item $\dfrac{x^2\,dx}{(x-1)^2(x^2-2x+4)}$.
\item $\dfrac{dx}{(1+x)^2(1+2x+4x^2)}$.
\item $\dfrac{x^4\,dx}{(x-1)^2(x^2+4)}$.
\item $\dfrac{dx}{(x+1)^2(x^2+1)}$.
\item $\dfrac{dx}{(x-1)^2(x^2+1)}$.
\item $\dfrac{dx}{x(x-1)^2(x^2+1)}$.
\item $\dfrac{dx}{x^3(a^2+x^2)}$.
\item $\dfrac{dx}{x^3(a^2+x^2)(b^2+x^2)}$.
\item $\dfrac{dx}{(x-1)^3(x^2+x+1)}$.
\item $\dfrac{dx}{(2x-3)^2(4x^2+5)}$.
\end{enumerate}

\item Repeated quadratic factors:
\begin{enumerate}[label=(\roman*)]
\item $\dfrac{dx}{x(x^2+1)^2}$.
\item $\dfrac{dx}{(x-1)^2(x^2+1)^2}$.
\item $\dfrac{(x+1)\,dx}{(x^2+1)^2}$.
\item $\dfrac{(x+a)(x+b)}{(x^2+c^2)^3}dx$.
\end{enumerate}

\item Evaluate $\displaystyle\int_0^{\pi/4}\sqrt{\tan\theta}\,d\theta$ and $\displaystyle\int_0^{\pi/4}\sqrt{\cot\theta}\,d\theta$.

\item Evaluate
\begin{enumerate}[label=(\roman*)]
\item $\displaystyle\int_0^{\pi/4} \dfrac{dx}{\cos^4x-\cos^2x\sin^2x+\sin^4x}$,
\item $\displaystyle\int_0^{\pi/4} \dfrac{dx}{\cos^4x+\cos^2x\sin^2x+\sin^4x}$.
\end{enumerate}

\item Evaluate $\displaystyle\int_0^{\pi/2} \dfrac{\cos x\,dx}{(1+\sin x)(2+\sin x)}$.

\item Show that $\displaystyle\int_0^\infty \dfrac{x^2\,dx}{(x^2+a^2)(x^2+b^2)(x^2+c^2)}=\dfrac{\pi}{2(a+b)(b+c)(c+a)}$.

\item Show that $\displaystyle\int_{-\infty}^{+\infty} \dfrac{dx}{(x^2\pm ax+a^2)(x^2\pm bx+b^2)}=\dfrac{2\pi}{\sqrt3}\dfrac{a+b}{ab(a^2+ab+b^2)}$. \source{$\gamma$, 1891.}

\item Show that the sum of the infinite series
\[
\dfrac{1}{a}-\dfrac{1}{a+b}+\dfrac{1}{a+2b}-\dfrac{1}{a+3b}+\dots \quad (a>0,\ b>0)
\]
can be expressed as a definite integral, viz.
\[
\int_0^1 \dfrac{t^{a-1}}{1+t^b}dt.
\]
And hence prove that
\[
\dfrac11-\dfrac14+\dfrac17-\dfrac{1}{10}+\dfrac{1}{13}-\dfrac{1}{16}+\dots = \dfrac13\left(\pi 3^{-\frac12}+\log_e 2\right). \source{Oxford, 1887.}
\]

\item Integrate:
\begin{enumerate}[label=(\roman*)]
\item $\displaystyle\int \dfrac{25\,dx}{2x^4+3x^3+3x-2}$. \source{Colleges, 1882.}
\item $\displaystyle\int \dfrac{x^5+2}{x^5-x}dx$. \source{St. John's, 1881.}
\item $\displaystyle\int \dfrac{(1+x^2)\,dx}{1-2x^2\cos\alpha+x^4}$. \source{Colleges, 1882.}
\item $\displaystyle\int \dfrac{dx}{1+x^5}$. \source{Colleges a, 1891.}
\item $\displaystyle\int_0^1 \dfrac{1+x^2}{(1-x^2)^2+a^2x^2}dx$.
\end{enumerate}

\item Prove that $\displaystyle\int_0^\infty \dfrac{dx}{1+x^6}=\dfrac{\pi}{3}$. \source{St. John's, 1881.}

\item Prove that
\[
\int \dfrac{dx}{(x-a)^p(x-b)^q}
\]
\[
=\sum_{r=0}^{r=p-2} \dfrac{Q_r}{(a-b)^{q+r}}\dfrac{(x-a)^{-p+r+1}}{-p+r+1}+\sum_{r=0}^{r=q-2} \dfrac{P_r}{(b-a)^{p+r}}\dfrac{(x-b)^{-q+r+1}}{-q+r+1}
\]
\[
+\dfrac{1}{(a-b)^{p+q-1}}Q_{p-1}\log(x-a)+\dfrac{1}{(b-a)^{p+q-1}}P_{q-1}\log(x-b),
\]
where $P_r$ and $Q_r$ are the coefficients of $z^r$ in $(1+z)^{-p}$ and $(1+z)^{-q}$ respectively.

\item Integrate
\begin{enumerate}[label=(\roman*)]
\item $\displaystyle\int \dfrac{dx}{(5x^3-3x)^2(x^2-1)}$. \source{Math. Trip., 1878.}
\item $\displaystyle\int \dfrac{dx}{(5x^3-3x^2)^2(x^2-1)}$.
\item $\displaystyle\int \dfrac{\sqrt{x}\,dx}{(1+x)(2+x)(3+x)}$. \source{Oxford I., 1888.}
\end{enumerate}

\item Prove
\begin{enumerate}[label=(\roman*)]
\item $\displaystyle\int_{-\infty}^{\infty} \dfrac{2x^2-x+1}{(x^2+x+1)^3}dx=\dfrac{10}{9}\pi\sqrt3$. \source{Colleges $\beta$, 1891.}
\item $\displaystyle\int_0^\infty \dfrac{\sqrt{x}}{1+x^2}dx=\dfrac{\pi}{\sqrt2}$. \source{Trinity, 1882.}
\item $\displaystyle\int_x^1 \dfrac{dx}{1+x^4}dx=\dfrac{1}{2\sqrt2}\left\{\pi-2\tan^{-1}\dfrac{\sqrt2}{x^{-1}-x}\right\}$. \source{Trinity, 1895.}
\end{enumerate}

\item Integrate
\[
\int \dfrac{x\,dx}{x^3+1}.
\]
Prove that $\dfrac{1}{2.5}+\dfrac{1}{8.11}+\dfrac{1}{14.17}+\dots$ to $\infty = \dfrac19\left[\dfrac{\pi}{\sqrt3}-\log 2\right]$. \source{Colleges, 1896.}

\item Integrate $\displaystyle\int \dfrac{(\sqrt{\cot x}-\sqrt{\tan x})\,dx}{1+3\sin 2x}$. \source{Colleges $\beta$, 1890.}

\item Integrate
\begin{enumerate}[label=(\roman*)]
\item $\displaystyle\int \tan^{-1}\sqrt{\dfrac{a^2x-1}{b^2x-1}}\,dx$.
\item $\displaystyle\int \sqrt{a^2+\sqrt{b^2+c/x}}\,dx$. \source{Math. Trip., 1898.}
\end{enumerate}

\item Integrate $\displaystyle\int \dfrac{(\sqrt a-\sqrt x)^2\,dx}{(a^2+ax+x^2)\sqrt{x}}$. \source{Colleges, 1896.}

\item Integrate $\displaystyle\int \dfrac{5x^3+3x-1}{(x^3+3x+1)^3}dx$. \source{J. M. Sch., Ox., 1904.}

\item Evaluate $\displaystyle\int_0^{\pi/4} \dfrac{\sin^4x}{\cos^5x}dx$. \source{St. John's, 1892.}

\item Integrate $\displaystyle\int \dfrac{dx}{x(x-a)^n}$, $n$ being a positive integer. \source{St. John's, 1882.}

\item Integrate $\displaystyle\int \left(\dfrac{x-b}{x-a}\right)^{\frac{n}{2}}dx$. \source{Colleges a, 1885.}

\item Integrate $\displaystyle\int \dfrac{dx}{1+3e^x+2e^{2x}}$. \source{Math. Trip., 1895.}

\item Sum the series
\[
\dfrac{x^3}{1.3}+\dfrac{x^5}{3.5}+\dfrac{x^7}{5.7}+\dots\ ad\ inf.,
\]
assuming it to be convergent.

Deduce that
\[
\dfrac{1}{1.3}\cdot\dfrac{1}{2^3}+\dfrac{1}{3.5}\cdot\dfrac{1}{2^5}+\dfrac{1}{5.7}\cdot\dfrac{1}{2^7}+\dots\ ad\ inf. = \dfrac14-\dfrac{3}{16}\log 3. \source{I. C. S., 1899.}
\]

\item Prove that
\[
1-\dfrac17+\dfrac19-\dfrac{1}{15}+\dfrac{1}{17}-\dots\ ad\ inf. = \dfrac{\pi}{8}(1+\sqrt2). \source{Colleges $\beta$, 1888.}
\]

\item Evaluate $\displaystyle\int x^2\log(1-x^2)\,dx$,

and deduce that
\[
\dfrac{1}{1.5}+\dfrac{1}{2.7}+\dfrac{1}{3.9}+\dots = \dfrac89-\dfrac23\log 2. \source{Colleges a, 1889.}
\]

\item Integrate $\displaystyle\int (a^4+x^4)^{-1}dx$.

Prove that
\[
\dfrac{1}{1.5}+\dfrac{1}{9.13}+\dfrac{1}{17.21}+\dfrac{1}{25.29}+\dots = \dfrac{\sqrt2}{32}\left\{\pi+\log(3+2\sqrt2)\right\}. \source{Math. Trip., 1896.}
\]

\item Show that
\[
\int \dfrac{dx}{x(x+1)(x+2)(x+3)\dots(x+n)} = \dfrac{1}{n!}\sum_{r=0}^{r=n}(-1)^r\, {}^nC_r\log(x+r).
\]

\item Show that
\[
\int \dfrac{(1+x)^n}{(1-2x)^3}dx = \dfrac{3^n}{2^{n+2}}\dfrac{1}{(1-2x)^2}-n\dfrac{3^{n-1}}{2^{n+1}}\dfrac{1}{1-2x}-\dfrac{n(n-1)3^{n-2}}{2^{n+2}}\log(1-2x)
\]
+ a rational integral algebraic expression of a finite number of terms.

\item Show that if $c<1$,
\[
\int \dfrac{dx}{(1-x)(1-cx)(1-c^2x)\dots\text{to }\infty}
\]
\[
=x+\dfrac12\dfrac{x^2}{1-c}+\dfrac13\dfrac{x^3}{(1-c)(1-c^2)}+\dfrac14\dfrac{x^4}{(1-c)(1-c^2)(1-c^3)}+\dots\text{to }\infty.
\]

\item Show that
\[
\int \dfrac{x^{n+2}\,dx}{(x-a_1)(x-a_2)(x-a_3)\dots(x-a_n)}
\]
\[
=\dfrac{x^3}{3}+H_1\dfrac{x^2}{2}+H_2x+\sum \dfrac{a_1^{n+2}}{(a_1-a_2)(a_1-a_3)\dots(a_1-a_n)}\log(x-a_1),
\]
where $H_r$ is the sum of the homogeneous products $r$ at a time of $a_1$, $a_2$, $\dots$, $a_n$.

\item Show that the part of the indefinite integral
\[
\int \dfrac{1}{x^3}\dfrac{f(x)}{\phi(x)}dx
\]
which becomes infinite when $x=0$, $f$ and $\phi$ being rational integral functions of $x$ which do not vanish when $x=0$, is
\[
-\dfrac{1}{2x^2}\dfrac{f(0)}{\phi(0)}-\dfrac1x\dfrac{f'(0)\phi(0)-\phi'(0)f(0)}{[\phi(0)]^2}
\]
\[
+\dfrac12\log x\dfrac{f''(0)\{\phi(0)\}^2-2f'(0)\phi'(0)\phi(0)-f(0)[\phi(0)\phi''(0)-2\{\phi'(0)\}^2]}{[\phi(0)]^3}. \source{Ox. I. P., 1901.}
\]

\item Show that when a rational fraction is decomposed into its simple or ``partial'' fractions, the decomposition is unique.

\item If $F(x)$ be a function of the $(n-1)$th degree which assumes the values $u_1$, $u_2$, $u_3$, $\dots$, $u_n$ when $x=x_0$, $x_1$, $x_2$, $\dots$, $x_n$ respectively, show that
\[
F(x)=u_1\dfrac{(x-x_2)(x-x_3)\dots(x-x_n)}{(x_1-x_2)(x_1-x_3)\dots(x_1-x_n)}
\]
\[
+u_2\dfrac{(x-x_1)(x-x_3)\dots(x-x_n)}{(x_2-x_1)(x_2-x_3)\dots(x_2-x_n)}
\]
\[
+\dots\dots\dots\dots\dots\dots\dots\dots\dots
\]
\[
+u_n\dfrac{(x-x_1)(x-x_2)\dots(x-x_{n-1})}{(x_n-x_1)(x_n-x_2)\dots(x_n-x_{n-1})}.
\]

\item Prove that if $p<n+1$,
\[
na^{n-p}\int \dfrac{x^{p-1}\,dx}{x^n-a^n} = \log(x-a) + \sum_{r=1}^{r=\frac{n-1}{2}}\cos\dfrac{2rp\pi}{n}\log\left(x^2-2ax\cos\dfrac{2r\pi}{n}+a^2\right)
\]
\[
-2\sum_{r=1}^{r=\frac{n-1}{2}}\sin\dfrac{2rp\pi}{n}\tan^{-1}\dfrac{x-a\cos\dfrac{2r\pi}{n}}{a\sin\dfrac{2r\pi}{n}}\qquad\text{if $n$ be odd,}
\]
and
\[
=\log(x-a)+(-1)^p\log(x+a)
\]
\[
+\sum_{r=1}^{r=\frac{n-2}{2}}\cos\dfrac{2rp\pi}{n}\log\left(x^2-2ax\cos\dfrac{2r\pi}{n}+a^2\right)
\]
\[
-2\sum_{r=1}^{r=\frac{n-2}{2}}\sin\dfrac{2rp\pi}{n}\tan^{-1}\dfrac{x-a\cos\dfrac{2r\pi}{n}}{a\sin\dfrac{2r\pi}{n}}\qquad\text{if $n$ be even.}
\]

\item Prove that if $p<n+1$,
\[
na^{n-p}\int \dfrac{x^{p-1}\,dx}{x^n+a^n} = (-1)^{p-1}\log(x+a)
\]
\[
-\sum_{r=1}^{r=\frac{n-1}{2}}\cos(2r-1)\dfrac{p\pi}{n}\log\left\{x^2-2ax\cos(2r-1)\dfrac{\pi}{n}+a^2\right\}
\]
\[
+2\sum_{r=1}^{r=\frac{n-1}{2}}\sin(2r-1)\dfrac{p\pi}{n}\tan^{-1}\dfrac{x-a\cos(2r-1)\dfrac{\pi}{n}}{a\sin(2r-1)\dfrac{\pi}{n}}\qquad\text{if $n$ be odd,}
\]
and
\[
=-\sum_{r=1}^{r=\frac{n}{2}}\cos(2r-1)\dfrac{p\pi}{n}\log\left\{x^2-2ax\cos(2r-1)\dfrac{\pi}{n}+a^2\right\}
\]
\[
+2\sum_{r=1}^{r=\frac{n}{2}}\sin(2r-1)\dfrac{p\pi}{n}\tan^{-1}\dfrac{x-a\cos(2r-1)\dfrac{\pi}{n}}{a\sin(2r-1)\dfrac{\pi}{n}}\qquad\text{if $n$ be even.}
\]

\item Prove that
\[
\int_0^x \dfrac{dx}{1-x^{2n}} = \dfrac{1}{2n}\sum_{r=0}^{r=n-1}\left(\cos\dfrac{r\pi}{n}\tanh^{-1}\dfrac{2x\cos\dfrac{r\pi}{n}}{1+x^2}+\sin\dfrac{r\pi}{n}\tan^{-1}\dfrac{2x\sin\dfrac{r\pi}{n}}{1-x^2}\right). \source{Math. Trip., 1884.}
\]

\item Show that $\displaystyle\int_0^\infty \dfrac{t}{t^5+1}dt = \dfrac{4\pi}{5\sqrt{10+2\sqrt5}}$.

\item
\begin{enumerate}[label=(\roman*)]
\item Show that the remainder left after dividing the rational integral function $f(x)$ by $(x-c)^2+b^2$ is
\[
\left[f(c)-\dfrac{b^2}{2!}f^{(2)}(c)+\dfrac{b^4}{4!}f^{(4)}(c)-\dots+(-1)^r\dfrac{b^{2r}}{(2r)!}f^{(2r)}(c)+\dots\right]
\]
\[
+(x-c)\left[f'(c)-\dfrac{b^3}{3!}f^{(3)}(c)+\dfrac{b^5}{5!}f^{(5)}(c)-\dots+(-1)^r\dfrac{b^{2r+1}}{(2r+1)!}f^{(2r+1)}(c)+\dots\right],
\]
where $f^{(s)}(c)$ denotes $\dfrac{d^sf(c)}{dc^s}$.
\item If $f(x)$ and $\phi(x)$ are rational integral functions of $x$, and $\phi(x)$ does not contain $(x-c)^2+b^2$ as a factor, show that it is possible to determine finite values for the constants $P$ and $Q$ in such a manner that
\[
f(x)-[P(x-c)+Q]\phi(x)
\]
is divisible, without remainder, by $(x-c)^2+b^2$.
\item Apply the last result to show (or prove in any manner) that
\[
\dfrac{f(x)}{[(x-c)^2+b^2]^r\phi(x)}
\]
can be expressed in the form
\[
\dfrac{\chi(x)}{\phi(x)}+\sum_{n=1}^{n=r}\dfrac{P_n(x-c)+Q_n}{[(x-c)^2+b^2]^n},
\]
$r$ being a positive integer, $\chi(x)$ a rational integral function of $x$, and $P_n$ and $Q_n$ constants. \source{I.C.S., 1892.}
\end{enumerate}

\item If
\[
\dfrac{F(x)}{f(x)}\equiv \phi(x)+\dfrac{\psi(x)}{f(x)},
\]
where $F$, $f$, $\phi$, $\psi$ are rational polynomials of degrees $m+n$, $n$, $m$, $n-1$ respectively, show that if $a_1$, $a_2$, $a_3$, $\dots$, $a_n$ be the roots of $f(x)=0$, considered all different, $\psi(x)$ will be determinable from
\[
\begin{vmatrix}
\psi(x), & x^{n-1}, & x^{n-2}, & \dots, & x, & 1 \\
F(a_1), & a_1^{n-1}, & a_1^{n-2}, & \dots, & a_1, & 1 \\
F(a_2), & a_2^{n-1}, & a_2^{n-2}, & \dots, & a_2, & 1 \\
\dots & \dots & \dots & \dots & \dots & \dots \\
F(a_n), & a_n^{n-1}, & a_n^{n-2}, & \dots, & a_n, & 1
\end{vmatrix} = 0.
\]
Also determine $\psi(x)$ when $f(x)=0$ has equal roots. \source{Oxford I. P., 1913.}

\item Integrate $\displaystyle\iint \left\{\dfrac{ax^2+bx+c}{(x-\alpha)(x-\beta)(x-\gamma)}\right\}^2 dx$. \source{Oxford I. P., 1917.}

\item Prove that if $n$ be a positive integer,
\begin{enumerate}[label=(\roman*)]
\item $\displaystyle\int \dfrac{\cos(n-2p)\theta}{\sin n\theta}d\theta = \dfrac1n\sum_{r=1}^{r=n}\cos\dfrac{2pr\pi}{n}\log\sin\left(\theta-\dfrac{r\pi}{n}\right)$,
\item $\displaystyle\int \dfrac{\sin(n-2p)\theta}{\sin n\theta}d\theta = \dfrac1n\sum_{r=1}^{r=(n-1)}\sin\dfrac{2pr\pi}{n}\log\,\text{cosec}\left(\theta-\dfrac{r\pi}{n}\right)$.
\end{enumerate}

\item Integrate
\begin{enumerate}[label=(\roman*)]
\item $\displaystyle\int \dfrac{dx}{1+e^x}$,
\item $\displaystyle\int \dfrac{dx}{\sqrt{\sin^3x\sin(x+a)}}$,
\end{enumerate}
and prove that
\begin{enumerate}[label=(\roman*)]
\setcounter{enumii}{2}
\item $\displaystyle\int_0^\infty \dfrac{x\,dx}{(1+x)(1+x^2)}=\dfrac{\pi}{4}$,
\item $\displaystyle\int_1^\infty \dfrac{(x^2+3)\,dx}{x^6(x^2+1)}=\dfrac{1}{30}(58-15\pi)$. \source{Math. Trip. I., 1917.}
\end{enumerate}

\item Obtain the rational part of $\displaystyle\int \dfrac{dx}{(x^5+1)^3}$. \source{Math. Trip. II., 1915.}

\item Prove that
\[
\dfrac{2n}{(1+x)^{2n}+(1-x)^{2n}} = \sum_{r=0}^{r=n}\dfrac{\sin a_r\cos^{2n-2}a_r}{\sin(2n-1)a_r}\dfrac{1}{x^2+\tan^2a_r},
\]
where $a_r=(2r+1)\pi/4n$. \source{Oxf. II. P., 1899.}

Write down the values of the integrals
\[
\int \dfrac{dx}{(1+x)^{2n}+(1-x)^{2n}}, \qquad \int \dfrac{x\,dx}{(1+x)^{2n}+(1-x)^{2n}}.
\]

\item Show that
\[
\int_0^\infty \dfrac{x\,dx}{(a+x)^n-(a-x)^n} = \dfrac{\pi}{2na^{n-2}}\Sigma(-1)^{\lambda-1}\sin\dfrac{\lambda\pi}{n}\cos^{n-3}\dfrac{\lambda\pi}{n},
\]
the summation extending from $\lambda=1$ to $\lambda=\dfrac{n-1}{2}$ or to $\lambda=\dfrac{n-2}{2}$, according as $n$ is odd or even. \source{Cf. Wolstenholme's Problems, No. 1912.}

Write down the value of the integral
\[
\int \dfrac{x^2\,dx}{(a+x)^n-(a-x)^n}.
\]

\item Show that if $n$ be even and $x+y=1$,
\[
\int \dfrac{dx}{x^ny^n} = \dfrac{1}{n-1}\left[\dfrac{1}{y^{n-1}}-\dfrac{1}{x^{n-1}}\right]+\dfrac{n}{1}\cdot\dfrac{1}{n-2}\left[\dfrac{1}{y^{n-2}}-\dfrac{1}{x^{n-2}}\right]
\]
\[
+\dfrac{n(n+1)}{1.2}\cdot\dfrac{1}{n-3}\left[\dfrac{1}{y^{n-3}}-\dfrac{1}{x^{n-3}}\right]+\dots+\dfrac{n(n+1)\dots(2n-2)}{1.2\dots(n-1)}\log\dfrac{x}{y}. \source{Murphy, Camb. Tr., vi.}
\]

\item Show that if $p<q$,
\[
Lt_{n=\infty}\dfrac{\displaystyle\prod_{r=1}^{r=n}\left(1-\dfrac{p^2x^2}{r^2\pi^2}\right)}{\displaystyle\prod_{r=1}^{r=n}\left(1-\dfrac{q^2x^2}{r^2\pi^2}\right)} = \sum_{r=1}^{r=\infty}\dfrac{2r\pi}{pq}\dfrac{\sin\left(\dfrac{p}{q}r\pi\right)}{\cos r\pi}\dfrac{1}{x^2-\dfrac{r^2\pi^2}{q^2}}. \source{Todhunter, I.C., p. 38.}
\]
Deduce that if $p<q$,
\[
\int \dfrac{\sin px}{\sin qx}dx = -2\sum_{r=1}^{r=\infty}\dfrac1q\dfrac{\sin\left(\dfrac{p}{q}r\pi\right)}{\cos r\pi}\tanh^{-1}\dfrac{qx}{r\pi} \qquad (-\pi<qx<\pi).
\]

\end{enumerate}

\newpage

\begin{center}
    {\LARGE \bfseries \color{accent} Chapter VI: Integrals of Forms $\displaystyle\int \dfrac{dx}{(a+b\cos x+c\sin x)^n}$, etc. --- Exercises}\\
    \vspace{0.2cm}
    {\color{accent}\rule{0.55\linewidth}{0.8pt}}
\end{center}
\vspace{0.35cm}

\section*{Exercises --- Various Standard Methods (Arts. 170--208)}

\begin{enumerate}[leftmargin=2em]

\item Integrate
\begin{enumerate}[label=(\roman*)]
\item $\displaystyle\int \dfrac{a\sin\theta+b\cos\theta}{c\sin\theta+e\cos\theta}d\theta$. \source{a, 1883.}
\item $\displaystyle\int \dfrac{d\theta}{\cos\theta-\sin\theta}$. \source{I.C.S., 1880.}
\item $\displaystyle\int \dfrac{a+\beta\sin\theta}{a+b\cos\theta}d\theta$. \source{Trin. H. and Magd.}
\item $\displaystyle\int \dfrac{dx}{\cos a+\cos x}$. \source{I.C.S., 1889.}
\item $\displaystyle\int \dfrac{dx}{a\cos x+b\sin x}$. \source{Coll., 1876.}
\item $\displaystyle\int \dfrac{1-\tan\theta}{1+\tan\theta}d\theta$. \source{Trin., 1884.}
\item $\displaystyle\int \dfrac{\sec\theta}{1+\text{cosec}\,\theta}d\theta$. \source{Ox. I. P., 1889.}
\item $\displaystyle\int \dfrac{dx}{3(1-\sin x)-\cos x}$. \source{a, 1881.}
\item $\displaystyle\int \dfrac{\sqrt2\,dx}{2\sqrt2+\cos x+\sin x}$. \source{Ox. I. P., 1888.}
\item $\displaystyle\int \dfrac{dx}{a+b\tan x}$. \source{St. John's, 1888.}
\item Apply the transformation $t=\tan\tfrac12 x$ to the integrals
\[
\int \dfrac{4\,dx}{5+3\cos x}, \qquad \int \dfrac{4\,dx}{3+5\cos x}.
\]
Hence or otherwise, evaluate these integrals to the nearest hundredth, when the limits are $x=0$ and $\tfrac14\pi$. Prove in any way that the second is the greater of the two integrals, when taken between $0$ and $\tfrac14\pi$. \source{Math. Trip. I., 1913.}
\item Prove that
\[
\int_0^a \dfrac{dx}{x+(a^2-x^2)^{\frac12}} = \tfrac14\pi, \qquad \int_0^a \dfrac{a\,dx}{\{x+(a^2-x^2)^{\frac12}\}^2} = \dfrac{1}{\sqrt2}\log_e(1+\sqrt2),
\]
the positive sign being taken for the radical in each of the subjects of integration. \source{Math. Trip. II., 1913.}
\end{enumerate}

\item Evaluate
\begin{enumerate}[label=(\roman*)]
\item $\displaystyle\int_0^\pi \dfrac{dx}{4+5\sin x}$. \source{I. C. S., 1889.}
\item $\displaystyle\int_0^\pi \dfrac{d\theta}{a+c\cos\theta}$ \ $(c<a)$. \source{I. C. S., 1879.}
\item $\displaystyle\int_0^{\frac\pi2} \dfrac{dx}{2+\cos x}$. \source{St. John's, 1882.}
\item $\displaystyle\int_0^\pi \dfrac{dx}{1-2a\cos x+a^2}$. \source{I. C. S., 1888.}
\end{enumerate}

\item Show that $\displaystyle\int_0^\pi \dfrac{dx}{1-\cos a\cos x} = \dfrac{\pi}{2}\text{cosec}\,a$, \source{Ox. II. P., 1889; Trin., 1887.}
and integrate
\[
\int \dfrac{\cos a\cos x+1}{\cos a+\cos x}dx.
\]

\item Evaluate
\begin{enumerate}[label=(\roman*)]
\item $\displaystyle\int \dfrac{dx}{a^2-\beta^2\cos^2x}$ \ $(a>\beta)$. \source{Ox. II. P., 1887.}
\item $\displaystyle\int \dfrac{dx}{\sin x\cos x+\sin x+\cos x-1}$. \source{$\gamma$, 1899.}
\item $\displaystyle\int \dfrac{dx}{\sin x+\sin 2x}$. \source{$\beta$, 1891.}
\end{enumerate}

\item Evaluate
\[
\int \dfrac{dx}{(e^x+e^{-x})^2}. \source{Ox. I. P., 1889.}
\]

\item Integrate
\begin{enumerate}[label=(\roman*)]
\item $\displaystyle\int \dfrac{dx}{(4+5\cos x)^2}$. \source{a, 1883.}
\item $\displaystyle\int \dfrac{dx}{(a+b\cos x)^2}$. \source{St. John's, 1884.}
\item $\displaystyle\int \dfrac{d\theta}{(a\cos\theta+b\sin\theta)^2}$. \source{a, 1881.}
\item $\displaystyle\int \dfrac{dx}{(a+b\cos x+c\sin x)^2}$. \source{Coll., 1892.}
\item Employ the substitution $\tan\dfrac\theta2=\sqrt{\dfrac{1+e}{1-e}}\tan\dfrac u2$ to evaluate the integral $\displaystyle\int_0^\phi \dfrac{d\theta}{(1+e\cos\theta)^2}$. \source{Math. Trip. I., 1909.}
\end{enumerate}

\item Prove that
\begin{enumerate}[label=(\roman*)]
\item $\displaystyle\int_0^\pi \dfrac{d\theta}{(1+\cos a\cos\theta)^2}=\pi\,\text{cosec}^3a$, \source{Coll., 1886; St. John's, 1886.}
\item $\displaystyle\int_0^\pi \dfrac{\sin x\,dx}{(1+\cos a\sin x)^2}=\dfrac{2(\sin a-a\cos a)}{\sin^3a}$, \source{$\beta$, 1887.}
\end{enumerate}
and evaluate
\[
\int_0^\pi \dfrac{x\sin x\,dx}{(1+\cos a\cos x)^3} \quad\text{and}\quad \int_0^\pi \dfrac{x\sin x\,dx}{(1+\cos a\cos x)^4},
\]
if $a$ be less than $\dfrac\pi2$.

\item Evaluate the integral
\[
\int \cos2\theta\log(1+\tan\theta)\,d\theta. \source{$\gamma$, 1882.}
\]

\item Find the values of the following integrals: \quad $\left(a<\dfrac\pi2\right)$
\begin{enumerate}[label=(\roman*)]
\item $\displaystyle\int_0^\infty \dfrac{dx}{2\cos a+e^x+e^{-x}}$.
\item $\displaystyle\int_0^\pi \dfrac{dx}{2+\cos a(e^x+e^{-x})}$. \source{Trin., 1882.}
\end{enumerate}

\item Evaluate
\begin{enumerate}[label=(\roman*)]
\item $\displaystyle\int_0^{\frac\pi2} \dfrac{d\theta}{a^2\cos^2\theta+b^2\sin^2\theta}$. \source{Ox. I. P., 1889.}
\item $\displaystyle\int_0^{\frac\pi2} \dfrac{d\theta}{4+5\sin^2\theta}$. \source{I. C. S., 1885.}
\item $\displaystyle\int_0^{\frac\pi2} \dfrac{a\sin^2\theta+b\cos^2\theta}{c\sin^2\theta+d\cos^2\theta}d\theta$ \ (c and d positive). \source{St. John's.}
\item $\displaystyle\int_0^{\frac\pi2} \dfrac{d\theta}{(a^2\cos^2\theta+\beta^2\sin^2\theta)^2}$. \source{Ox. II. P., 1887.}
\item $\displaystyle\int_0^{\frac\pi4} \dfrac{\sin2\theta\,d\theta}{\sin^4\theta+\cos^4\theta}$, \source{I. C. S., 1891.}
\end{enumerate}
and
\begin{enumerate}[label=(\roman*)]
\setcounter{enumii}{5}
\item Shew that if $c>a>0$,
\[
\int_0^a \dfrac{\sqrt{a^2-x^2}}{c^2-x^2}dx = \pi(c-\sqrt{c^2-a^2})/2c. \source{Math. Trip. I., 1908.}
\]
\end{enumerate}

\item Prove that $\displaystyle\int_0^\pi \dfrac{d\theta}{a+b\cos\theta}=\dfrac{\pi}{\sqrt{a^2-b^2}}$,

where $a>b$, and deduce or otherwise obtain the value of
\[
\int_0^\pi \dfrac{d\theta}{(a+b\cos\theta)^3}. \source{$\gamma$, 1899.}
\]

\item Prove that if $a>b$,
\[
\int_0^\pi \dfrac{dx}{(a+b\cos x)^4} = \dfrac\pi2\left\{\dfrac{5a^3}{(a^2-b^2)^{\frac72}}-\dfrac{3a}{(a^2-b^2)^{\frac52}}\right\}. \source{$\gamma$, 1888.}
\]
Evaluate $\displaystyle\int_0^\pi \dfrac{dx}{(1+e\cos x)^4}$, where $e<1$. \source{St. John's, 1892.}

\item Evaluate $\displaystyle\int \dfrac{dx}{a+be^x+ce^{-x}}$, where $a^2<4bc$. \source{I. C. S., 1897.}

\item Prove that if $a<\dfrac\pi6$,
\[
\int_0^a \dfrac{\sin x}{\cos 3x}dx = \dfrac16\log\dfrac{\cos^3a}{\cos3a}. \source{C. S., 1896.}
\]

\item Prove that $\displaystyle\int_0^{2\pi} \dfrac{d\phi}{a+b\cos\phi+c\sin\phi} = \dfrac{2\pi}{\sqrt{a^2-r^2}}$,

where $r^2=b^2+c^2$ and $r<a$. \source{C. S., 1900.}

\item Integrate
\begin{enumerate}[label=(\roman*)]
\item $\displaystyle\int \dfrac{\sqrt{\tan x}}{\sin x\cos x}dx$.
\item $\displaystyle\int \dfrac{\sec x\,dx}{a+b\tan x}$.
\item $\displaystyle\int \dfrac{\sec^2x\,dx}{a+b\tan x}$.
\end{enumerate}

\item Integrate
\begin{enumerate}[label=(\roman*)]
\item $\displaystyle\int \dfrac{d\theta}{15\sin^2\theta-16\cos\theta}$.
\item $\displaystyle\int_0^\pi \dfrac{x}{1+\sin x}dx$.
\item $\displaystyle\int \dfrac{\cot\theta-3\cot3\theta}{3\tan3\theta-\tan\theta}d\theta$. \source{Ox. I. P., 1888.}
\item $\displaystyle\int \dfrac{\sin2x\,dx}{(a+b\cos x)^2}$. \source{a, 1889.}
\end{enumerate}

\item Integrate
\begin{enumerate}[label=(\roman*)]
\item $\displaystyle\int \cos2\theta\log\dfrac{\cos\theta+\sin\theta}{\cos\theta-\sin\theta}\,d\theta$.
\item $\displaystyle\int \dfrac{\sin\theta-\cos\theta}{\sqrt{\sin2\theta}}d\theta$.
\item $\displaystyle\int \sqrt{\dfrac{1-\cos\theta}{\cos\theta(1+\cos\theta)(2+\cos\theta)}}\,d\theta$.
\end{enumerate}

\item Integrate $\displaystyle\int \sqrt{\dfrac{1+\sin x}{1-\sin x}\cdot\dfrac{2+\sin x}{2-\sin x}}\,dx$.

\item Integrate
\[
\int \dfrac{\sin\theta-\cos\theta}{(\sin\theta+\cos\theta)\sqrt{\sin\theta\cos\theta+\sin^2\theta\cos^2\theta}}d\theta.
\]

\item Integrate $\displaystyle\int \dfrac{\sin^3\theta\,d\theta}{(1+\cos^2\theta)\sqrt{1+\cos^2\theta+\cos^4\theta}}$.

\item Integrate
\begin{enumerate}[label=(\roman*)]
\item $\displaystyle\int \sqrt{1+\sin x}\,dx$.
\item $\displaystyle\int \dfrac{\sin x}{\sqrt{1+\sin x}}dx$.
\item $\displaystyle\int \dfrac{\tan x}{\sqrt{a+b\tan^2x}}dx$.
\end{enumerate}

\item Integrate $\displaystyle\int \dfrac{\sinh x\sin x-\cosh x}{1-\cos x}dx$.

\item Integrate $\displaystyle\int \dfrac{x^2\,dx}{(x\sin x+\cos x)^2}$.

\item Integrate $\displaystyle\int \dfrac{\sec x\,\text{cosec}\,x}{\log\tan x}dx$.

\item Integrate
\begin{enumerate}[label=(\roman*)]
\item $\displaystyle\int \sin^{-1}\dfrac{2x}{1+x^2}dx$.
\item $\displaystyle\int \tan^{-1}\dfrac{3x-x^3}{1-3x^2}dx$.
\item $\displaystyle\int \tan^{-1}\dfrac{-1+\sqrt{1+x^2}}{x}dx$.
\end{enumerate}

\item Integrate $\displaystyle\int \dfrac{\sqrt{1+x^2}}{1-x^2}dx$.

\item Integrate $\displaystyle\int \dfrac{\sin x}{\sin 2x}dx$, $\displaystyle\int \dfrac{\sin x}{\sin 3x}dx$, $\displaystyle\int \dfrac{\sin x}{\sin 4x}dx$,

and prove that
\[
5\int \dfrac{\sin x}{\sin 5x}dx = \sin\dfrac{2\pi}{5}\log\left\{\dfrac{\sin\left(x-\dfrac{2\pi}{5}\right)}{\sin\left(x+\dfrac{2\pi}{5}\right)}\right\} - \sin\dfrac\pi5\log\left\{\dfrac{\sin\left(x-\dfrac\pi5\right)}{\sin\left(x+\dfrac\pi5\right)}\right\}. \source{Trin. Coll., 1892.}
\]

\item Integrate $\displaystyle\int \dfrac{d\theta}{\sin^2\theta-\sin^2a}$ and $\displaystyle\int \dfrac{\cos\theta}{\sin^2\theta-\sin^2a}d\theta$.

Show that $\dfrac{\sin\theta}{\sin n\theta}$ can be expressed in partial fractions of type
\[
\dfrac{1}{\sin^2\theta-\sin^2a} \quad\text{or}\quad \dfrac{\cos\theta}{\sin^2\theta-\sin^2a},
\]
according as $n$ is an odd or an even integer and can thereby be integrated.

\item Integrate
\begin{enumerate}[label=(\roman*)]
\item $\displaystyle\int_0^{\frac\pi2} \dfrac{\sin3x\,dx}{(a+b\cos x)^2}$.
\item $\displaystyle\int_0^{\frac\pi2} \dfrac{\sin^3x\,dx}{a+b\cos x}$.
\item $\displaystyle\int_0^{\frac\pi2} \dfrac{\sin^3x\,dx}{(a+b\cos x)^2}$.
\item $\displaystyle\int_0^{\frac\pi2} \dfrac{\sin3x\,dx}{(a+b\cos x)^n}$ \ $(n>3)$.
\end{enumerate}

\item Show how to effect the integration of
\[
\int \dfrac{\cos^px}{\sin2nx}dx, \qquad \int \dfrac{\cos^px}{\cos nx}dx,
\]
$p$ and $n$ being integers. \source{$\epsilon$, 1883, and Coll., 1879.}

\item Integrate $\displaystyle\int \cot(x-a)\cot(x-\beta)\,dx$, \source{$\gamma$, 1891.}

and show that
\[
\int \cot(x-a)\cot(x-b)\cot(x-c)\,dx = \Sigma\cot(a-b)\cot(a-c)\log\sin(x-a). \source{Trinity, 1891.}
\]

\item Show that
\[
\int \sin x\sec(x-a)\sec(x-\beta)\,dx = \dfrac{1}{\sin(\beta-a)}\left[\cos a\cosh^{-1}\sec(x-a)-\cos\beta\cosh^{-1}\sec(x-\beta)\right]. \source{Trinity, 1889.}
\]

\item Prove that
\[
\int_0^\beta x\sec x\sec(\beta-x)\,dx = \beta\,\text{cosec}\,\beta\log\sec\beta. \source{Oxf. II. P., 1901.}
\]

\item Prove that, if $a$ and $\beta$ be positive quantities,
\[
\int_0^{\frac\pi2} \dfrac{dx}{(a\cos^2x+\beta\sin^2x)^{n+1}} = \dfrac\pi2\dfrac{(-1)^n}{n!}\left(\dfrac{d}{da}+\dfrac{d}{d\beta}\right)^n(a\beta)^{-\frac12}. \source{a, 1884.}
\]

\item Prove that, if
\[
f(x)=(x-a_1)(x-a_2)\dots(x-a_n)
\]
and
\[
P=\prod_{r=1}^{r=n}(1-2a_r\cos\theta+a_r^2),
\]
where $a_1,a_2,\dots a_n$ denote real quantities, then
\[
\int_0^\pi \dfrac{d\theta}{P} = \pi\sum_1^n \epsilon_r/A_r,
\]
where $A_r=a_rf(1/a_r)f'(a_r)$, and $\epsilon=-1$, or $+1$, according as $a_r$ is numerically greater or less than $1$. \source{St. John's, 1886.}

\item If $c$ be less than $a\sin\theta$, show that the coefficient of $c^m$ in the expansion of
\[
\dfrac{2c}{\sqrt{c^2-a^2}}\tan^{-1}\left\{\sqrt{\dfrac{c-a}{c+a}}\tan\left(\dfrac\pi4-\dfrac\theta2\right)\right\}
\]
is
\[
(-1)^m\dfrac{1}{a^m}\int \dfrac{d\theta}{\sin^m\theta}+A_m,
\]
where $A_m$ is independent of $\theta$. \source{Coll., 1892.}

\item Show that if $n$ be a positive integer,
\[
\int \dfrac{\cos n\theta-\cos na}{\cos\theta-\cos a}d\theta = \theta\dfrac{\sin na}{\sin a}+\dfrac{2}{\sin a}\sum_{r=1}^{r=n-1}\dfrac{\sin r\theta\sin(n-r)a}{r}. \source{Hermite.}
\]

\item Prove that
\[
\int_0^x (1+\cos x)^n dx = \dfrac{1}{2^n}\cdot\binom{2n}{n}x+\dfrac{1}{2^{n-1}}\sum_1^n \binom{2n}{n-r}\dfrac{\sin rx}{r}.
\]

\item Show that
\[
\int \prod_1^n \cot(\theta-a_r)\,d\theta = \theta\cos\dfrac{n\pi}{2}+\sum_1^n A_r\log\sin(\theta-a_r),
\]
where $A_r=\cot(a_r-a_1)\cot(a_r-a_2)\dots\cot(a_r-a_n)$, the factor $\cot(a_r-a_r)$ being omitted. \source{Hermite.}

\item (1) Show that
\[
\int \dfrac{dx}{a-\sin x} = \dfrac{2}{\sqrt{a^2-1}}\tan^{-1}\left(\dfrac{a\tan\dfrac x2-1}{\sqrt{a^2-1}}\right) \qquad (a>1).
\]

(2) Differentiate with regard to $x$,
\[
\tan^{-1}\sqrt{\dfrac{a+1}{a-1}\cdot\dfrac{1-\sin x}{1+\sin x}}.
\]

Deduce from (1) and (2) that
\[
\tan^{-1}\sqrt{\dfrac{a+1}{a-1}\cdot\dfrac{1-\sin x}{1+\sin x}}+\tan^{-1}\dfrac{a\tan\dfrac x2-1}{\sqrt{a^2-1}}
\]
is independent of $x$, and verify your conclusion. \source{C. S., 1898.}

\item Integrate $\displaystyle\int \dfrac{dx}{1+ab-a\cos x-b\sec x}$,

where $a<1$, $b>1$, $ab\neq1$. \source{Oxf. I. P., 1917.}

\item Integrate
\begin{enumerate}[label=(\roman*)]
\item $e^x(x\cos x+\sin x)$.
\item $(x^2+1)^{-\frac32}$.
\item $(x^4+2x^3+2x+1)/(x^2+x+1)^3$. \source{Oxf. I. P., 1918.}
\end{enumerate}

\item Integrate $\displaystyle\int \dfrac{\cos^3x\,dx}{\sin(x-a)\sin(x-\beta)\sin(x-\gamma)}$.

\item Deduce from the identity
\[
\int_0^{\frac\pi2}\cos n\theta\,d\theta = \int_0^{\frac\pi2}d\theta\left\{1-\dfrac{n^2}{2!}\sin^2\theta-\dfrac{n^2(2^2-n^2)}{4!}\sin^4\theta-\dfrac{n^2(2^2-n^2)(4^2-n^2)}{6!}\sin^6\theta-\dots\right\}
\]
the expression for $\sin x$ as an infinite product. \source{Oxf. II. P., 1887.}

\item Evaluate the integrals
\begin{enumerate}[label=(\roman*)]
\item $\displaystyle\int (x+a\log x)^2\,dx$.
\item $\displaystyle\int \dfrac{(a+x)\,dx}{x^2+ax\log x}$.
\item $\displaystyle\int \dfrac{(1-\log x)\,dx}{(x+a\log x)^2}$. \source{Math. Tripos, 1885.}
\end{enumerate}

\item Show that
\[
\int_0^x \dfrac{x^2+a(a-1)}{(x\sin x+a\cos x)^2}dx = \dfrac{a\sin x-x\cos x}{x\sin x+a\cos x}. \source{Trin. Coll., 1891.}
\]

\item Evaluate the indefinite integrals
\begin{enumerate}[label=(\roman*)]
\item $\displaystyle\int \dfrac{(\sin x+\cos x)^2}{\{(x-1)\cos x-(x+1)\sin x\}^2}dx$.
\item $\displaystyle\int \dfrac{x^2}{\{(x-1)\cos x-(x+1)\sin x\}^2}dx$. \source{Colleges, 1886.}
\end{enumerate}

\item If $T\equiv x+\sqrt{1+x^2}$, show that
\[
\int T^r\,dx = \tfrac12\left\{\dfrac{T^{r-1}}{r-1}+\dfrac{T^{r+1}}{r+1}\right\}.
\]

\item Integrate
\begin{enumerate}[label=(\roman*)]
\item $\displaystyle\int e^x\dfrac{x^3-x+2}{(1+x^2)^2}dx$. \source{Ox. II. P., 1899.}
\item $\displaystyle\int \dfrac{\log(\cos\theta+\sqrt{\cos2\theta})}{1-\cos^2\theta}d\theta$. \source{Coll. a, 1891.}
\end{enumerate}

\item Prove that, if $n$ be an integer,
\[
\int_0^\pi \dfrac{\cos nx}{1+\cos a\cos x}dx = \pi\,\text{cosec}\,a(\tan a-\sec a)^n,
\]
and deduce the value of
\[
\int_0^\pi \dfrac{\cos nx}{(1+\cos a\cos x)^2}dx. \source{Colleges $\gamma$, 1891.}
\]

\item From considering the integral
\[
\int_0^\pi \dfrac{\cos n\theta}{a-\cos\theta}d\theta,
\]
show that
\[
1+\dfrac{n+2}{1}\dfrac{\cos^2a}{2^2}+\dfrac{(n+3)(n+4)}{1.2}\dfrac{\cos^4a}{2^4}+\dots = 2^n(\sec a-\tan a)^n\sec^na\,\text{cosec}\,a.
\]

\item Prove that, if $0<a<\dfrac\pi2$ and $n$ be a positive integer,
\[
\int_0^\pi \sin n\phi\tan^{-1}\left(\dfrac{\tan\dfrac\phi2}{\tan\dfrac a2}\right)d\phi = \dfrac{\pi}{2n}\left[(\sec a-\tan a)^n-(-1)^n\right].
\]

\item Show that
\[
\int \dfrac{\sin n\theta}{\sin\theta}\dfrac{d\theta}{\cos n\theta-\cos na} = \dfrac1n\sum_{r=0}^{r=n-1}\text{cosec}\,\beta_r\log\dfrac{\sin\tfrac12(\beta_r+\theta)}{\sin\tfrac12(\beta_r-\theta)},
\]
where $\beta_r=a+\dfrac{2r\pi}{n}$.

\item Discuss the integration of
\[
(a)\ \int \dfrac{\sin p\theta}{\sin q\theta}d\theta, \qquad (b)\ \int \sin\theta\dfrac{\sin p\theta}{\sin q\theta}d\theta,
\]
where $p$ and $q$ are positive integers.

\item With the help of the substitution $x^{-1}=\sqrt{t^2-1}$, or otherwise prove that
\[
\int_0^\infty \dfrac{dx}{(9+25x^2)\sqrt{1+x^2}} = \dfrac{1}{12}\tan^{-1}\dfrac43. \source{Math. Trip., Pt. II., 1920.}
\]

\end{enumerate}

\begin{center}
    {\LARGE \bfseries \color{accent} Chapter VII: Further Reduction Formulae --- Exercises}\\
    \vspace{0.2cm}
    {\color{accent}\rule{0.55\linewidth}{0.8pt}}
\end{center}
\vspace{0.35cm}

\section*{Exercises --- Integrals of Form $\displaystyle\int x^{m-1}(a+bx^n)^p\,dx$ (Arts. 211--223)}

\begin{enumerate}[leftmargin=2em]

\item Prove that

\[
\int x^{m-1}(a+bx)^p\,dx=\dfrac{x^m(a+bx)^p}{m+p}+\dfrac{ap}{m+p}\int x^{m-1}(a+bx)^{p-1}\,dx.
\]

\item
\[
\int \dfrac{(a+bx)^p}{x^{m+1}}\,dx=-\dfrac{(a+bx)^{p+1}}{max^m}+\dfrac{(p-m+1)}{m}\dfrac{b}{a}\int \dfrac{(a+bx)^p}{x^m}\,dx.
\]

\item
\[
\int \dfrac{dx}{x^m(a+bx)}=-\dfrac{1}{m-1}\cdot\dfrac1a\cdot\dfrac{1}{x^{m-1}}+\dfrac{1}{m-2}\cdot\dfrac{b}{a^2}\cdot\dfrac{1}{x^{m-2}}-\dfrac{1}{m-3}\cdot\dfrac{b^2}{a^3}\cdot\dfrac{1}{x^{m-3}}
+\dots+(-1)^{m-1}\dfrac{1}{1}\cdot\dfrac{b^{m-2}}{a^{m-1}}\cdot\dfrac1x+(-1)^m\dfrac{b^{m-1}}{a^m}\log\dfrac{a+bx}{x}.
\]

\item
\[
\int \dfrac{(a+bx)^p}{x}\,dx=\dfrac{(a+bx)^p}{p}+a\int \dfrac{(a+bx)^{p-1}}{x}\,dx. \source{Bertrand.}
\]

\item
\[
\int \dfrac{dx}{(a+bx^2)^{p+1}}=\dfrac{x}{2ap(a+bx^2)^p}+\dfrac{2p-1}{2ap}\int \dfrac{dx}{(a+bx^2)^p}. \source{Bertrand.}
\]

\item
\[
\int \dfrac{x^n\,dx}{(a+bx^2)^{p+1}}=\dfrac{x^{n+1}}{3ap(a+bx^2)^p}-\dfrac{n-3p+1}{3ap}\int \dfrac{x^n}{(a+bx^2)^p}\,dx, \source{Bertrand.}
\]

and evaluate
\[
\int \dfrac{x^7}{a+bx^3}\,dx, \qquad \int \dfrac{x^3}{(a+bx^3)^2}\,dx, \qquad \int \dfrac{dx}{x^2(a+bx^3)}.
\]

\item
\[
\int x^n(a+bx^4)^p\,dx=\dfrac{x^{n+1}(a+bx^4)^p}{n+1}-\dfrac{4bp}{n+1}\int x^{n+4}(a+bx^4)^{p-1}\,dx,
\]
\[
\int \dfrac{x^n}{a+bx^4}\,dx=\dfrac{x^{n-3}}{(n-3)b}-\dfrac{a}{b}\int \dfrac{x^{n-4}}{a+bx^4}\,dx,
\]
\[
\int \dfrac{dx}{(a+bx^4)^{p+1}}=\dfrac{x}{4ap(a+bx^4)^p}+\dfrac{4p-1}{4ap}\int \dfrac{dx}{(a+bx^4)^p}, \source{Bertrand.}
\]

and evaluate
\[
\int \dfrac{dx}{(a+bx^4)^4}, \qquad \int \dfrac{dx}{x^2(a+bx^4)^3}.
\]

\end{enumerate}

\section*{Exercises --- Reduction Formulae for $\displaystyle\int \sin^px\cos^qx\,dx$ (Arts. 224--231)}

\begin{enumerate}[leftmargin=2em]

\item Prove that
\[
\int \sin^p\theta\cos^q\theta\,d\theta=\dfrac{\sin^{p+1}\theta\cos^{q-1}\theta}{p+q}-\dfrac{q-1}{(p+q)(p+q-2)}\sin^{p-1}\theta\cos^{q-1}\theta+\dfrac{(p-1)(q-1)}{(p+q)(p+q-2)}\int \sin^{p-2}\theta\cos^{q-2}\theta\,d\theta
\]
the indices being both diminished.

\item Prove that
\[
\int \dfrac{\sin^p\theta}{\cos^q\theta}\,d\theta=\dfrac{\sin^{p-1}\theta}{(q-1)\cos^{q-1}\theta}-\dfrac{p-1}{q-1}\int \dfrac{\sin^{p-2}\theta}{\cos^{q-2}\theta}\,d\theta.
\]

\item Prove that
\[
\int \dfrac{\cos^p\theta}{\sin^q\theta}\,d\theta=-\dfrac{\cos^{p-1}\theta}{(q-1)\sin^{q-1}\theta}-\dfrac{p-1}{q-1}\int \dfrac{\cos^{p-2}\theta}{\sin^{q-2}\theta}\,d\theta.
\]

\item Prove that
\[
\int \dfrac{d\theta}{\sin^p\theta\cos^q\theta}=\dfrac{1}{(q-1)\sin^{p-1}\theta\cos^{q-1}\theta}+\dfrac{p+q-2}{q-1}\int \dfrac{d\theta}{\sin^p\theta\cos^{q-2}\theta}
\]
\[
=-\dfrac{1}{(p-1)\sin^{p-1}\theta\cos^{q-1}\theta}+\dfrac{p+q-2}{p-1}\int \dfrac{d\theta}{\sin^{p-2}\theta\cos^{q}\theta}.
\]

\item
\[
\int \dfrac{\sin^p\theta}{\cos\theta}\,d\theta=-\dfrac{\sin^{p-1}\theta}{p-1}+\int \dfrac{\sin^{p-2}\theta}{\cos\theta}\,d\theta.
\]

\item
\[
\int \dfrac{\cos^p\theta}{\sin\theta}\,d\theta=\dfrac{\cos^{p-1}\theta}{p-1}+\int \dfrac{\cos^{p-2}\theta}{\sin\theta}\,d\theta.
\]

\item (a)
\[
\int \sin^{2n}\theta\,d\theta=-\dfrac{c}{2n}\left[s^{2n-1}+\dfrac{2n-1}{2n-2}s^{2n-3}+\dfrac{(2n-1)(2n-3)}{(2n-2)(2n-4)}s^{2n-5}
+\dots+\dfrac{1.3\dots(2n-1)}{2.4\dots(2n-2)}s\right]+\dfrac{1.3\dots(2n-1)}{2.4\dots2n}\theta.
\]

(b)
\[
\int \sin^{2n+1}\theta\,d\theta=-\dfrac{c}{2n+1}\left[s^{2n}+\dfrac{2n}{2n-1}s^{2n-2}+\dfrac{2n(2n-2)}{(2n-1)(2n-3)}s^{2n-4}
+\dots+\dfrac{2.4\dots2n}{1.3\dots(2n-1)}\right],
\]
where $c$ and $s$ stand respectively for $\cos\theta$ and $\sin\theta$. \source{Bertrand.}

\item (a)
\[
\int \cos^{2n}\theta\,d\theta=\dfrac{s}{2n}\left[c^{2n-1}+\dfrac{2n-1}{2n-2}c^{2n-3}+\dfrac{(2n-1)(2n-3)}{(2n-2)(2n-4)}c^{2n-5}
+\dots+\dfrac{1.3\dots(2n-1)}{2.4\dots(2n-2)}c\right]+\dfrac{1.3\dots(2n-1)}{2.4\dots2n}\theta.
\]

(b)
\[
\int \cos^{2n+1}\theta\,d\theta=\dfrac{s}{2n+1}\left[c^{2n}+\dfrac{2n}{2n-1}c^{2n-2}+\dots+\dfrac{2.4\dots2n}{1.3\dots(2n-1)}\right],
\]
$c$ and $s$ being respectively $\cos\theta$ and $\sin\theta$. \source{Bertrand.}

\item Prove

(a)
\[
\int \text{cosec}^5\theta\,d\theta=-\dfrac14\dfrac{c}{s^4}-\dfrac{1.3}{2.4}\dfrac{c}{s^2}+\dfrac{1.3}{2.4}\log\tan\dfrac{\theta}{2}.
\]

(b)
\[
\int \sec^5\theta\,d\theta=\dfrac14\dfrac{s}{c^4}+\dfrac{1.2}{2.4}\dfrac{s}{c^2}+\dfrac38\log\tan\left(\dfrac{\theta}{2}+\dfrac{\pi}{4}\right),
\]
where $c\equiv\cos\theta$, $s\equiv\sin\theta$.

\item Prove that

(a)
\[
\int \cos^{2n}\theta\,d\theta=\dfrac{1}{2^{2n-1}}\left[\dfrac{\sin2n\theta}{2n}+2n\dfrac{\sin(2n-2)\theta}{2n-2}+\dfrac{2n(2n-1)}{1.2}\dfrac{\sin(2n-4)\theta}{2n-4}+\dots+\dfrac{2n(2n-1)\dots(n+1)}{1.2\dots n}\theta\right].
\]

(b)
\[
\int \cos^{2n+1}\theta\,d\theta=\dfrac{1}{2^{2n}}\left[\dfrac{\sin(2n+1)\theta}{2n+1}+(2n+1)\dfrac{\sin(2n-1)\theta}{2n-1}
+\dfrac{(2n+1)2n}{1.2}\dfrac{\sin(2n-3)\theta}{2n-3}+\dots+\dfrac{(2n+1)2n\dots(n+2)}{1.2\dots n}\sin\theta\right].
\]

(c)
\[
\int \sin^{2n}\theta\,d\theta=\dfrac{(-1)^n}{2^{2n-1}}\left[\dfrac{\sin2n\theta}{2n}-2n\dfrac{\sin(2n-2)\theta}{2n-2}
+\dfrac{2n(2n-1)}{1.2}\dfrac{\sin(2n-4)\theta}{2n-4}-\dots
+(-1)^{n-1}\dfrac{2n(2n-1)\dots(n+2)}{1.2\dots(n-1)}\dfrac{\sin2\theta}{2}
+(-1)^n\dfrac{2n(2n-1)\dots(n+1)}{1.2\dots n}\dfrac{\theta}{2}\right].
\]

(d)
\[
\int \sin^{2n+1}\theta\,d\theta=\dfrac{(-1)^{n+1}}{2^{2n}}\left[\dfrac{\cos(2n+1)\theta}{2n+1}-(2n+1)\dfrac{\cos(2n-1)\theta}{2n-1}+\dots
+(-1)^n\dfrac{(2n+1)\dots(n+2)}{1.2\dots n}\cos\theta\right]. \source{Bertrand.}
\]

\end{enumerate}

\section*{Exercises --- Integrals of Form $\displaystyle\int x^m(a+bx+cx^2)^p\,dx$ (Art. 247)}

\begin{enumerate}[leftmargin=2em]

\item Taking
\[
\int \dfrac{dx}{X^{p+1}}=\dfrac{b+2cx}{pkX^p}+\dfrac{2(2p-1)c}{pk}\int \dfrac{dx}{X^p},
\]
where $X\equiv a+bx+cx^2$ and $k\equiv4ac-b^2$,

prove
\[
\int \dfrac{dx}{X^2}=\dfrac{b+2cx}{kX}+\dfrac{2c}{k}\int \dfrac{dx}{X},
\]
\[
\int \dfrac{dx}{X^3}=\dfrac{(b+2cx)}{k}\left(\dfrac{1}{2X^2}+\dfrac{3c}{kX}\right)+\dfrac{6c^2}{k^2}\int \dfrac{dx}{X},
\]
\[
\int \dfrac{dx}{X^4}=\dfrac{b+2cx}{k}\left(\dfrac{1}{3X^3}+\dfrac{5c}{3kX^2}+\dfrac{10c^2}{k^2X}\right)+\dfrac{20c^3}{k^3}\int \dfrac{dx}{X}. \source{Bertrand.}
\]

\item Show that if $I_n\equiv\displaystyle\int \dfrac{x^n}{X}\,dx$, then $cI_n+bI_{n-1}+aI_{n-2}=\dfrac{x^{n-1}}{n-1}$, and prove
\[
\int \dfrac{x\,dx}{X}=\dfrac{1}{2c}\log X-\dfrac{b}{2c}\int \dfrac{dx}{X}.
\]
Deduce
\[
\int \dfrac{x^2\,dx}{X}=\dfrac{x}{c}-\dfrac{b}{2c^2}\log X+\dfrac{b^2-2ac}{2c^2}\int \dfrac{dx}{X},
\]
\[
\int \dfrac{x^3\,dx}{X}=\dfrac{x^2}{2c}-\dfrac{bx}{c^2}+\dfrac{b^2-ac}{2c^3}\log X+\dfrac{3ac-b^2}{2c^3}b\int \dfrac{dx}{X}. \source{Bertrand.}
\]

\item Prove
\[
\int \dfrac{dx}{xX}=\dfrac{1}{2a}\log\dfrac{x^2}{X}-\dfrac{b}{2a}\int \dfrac{dx}{X},
\]
and deduce
\[
\int \dfrac{dx}{x^3X}=\dfrac{b}{a^2x}-\dfrac{1}{2ax^2}+\dfrac{b^2-ac}{2a^3}\log\dfrac{x^2}{X}+\dfrac{b(3ac-b^2)}{2a^3}\int \dfrac{dx}{X}. \source{Bertrand.}
\]

(The value of $\displaystyle\int \dfrac{dx}{X}$ occurring in each of these results is given in Art. 80.)

\item If $X\equiv a+bx+cx^2$, prove that
\[
\int \dfrac{x^{m+1}}{X^{p+1}}\,dx=-\dfrac{x^m}{2cpX^p}+\dfrac{m}{2cp}\int \dfrac{x^{m-1}\,dx}{X^p}-\dfrac{b}{2c}\int \dfrac{x^m}{X^{p+1}}\,dx. \source{Bertrand.}
\]

\item Prove that if $X\equiv x^2+ax+a^2$,
\[
\int X^{\frac{n}{2}}\,dx=\dfrac{2x+a}{2(n+1)}X^{\frac{n}{2}}+\dfrac{3na^2}{4(n+1)}\int X^{\frac{n}{2}-1}\,dx. \source{St. John's, 1889.}
\]

\item Prove that if $X\equiv x^2+x+1$,

(a)
\[
\int \dfrac{dx}{x^4X}=\dfrac12\log\dfrac{x^2}{X}-\dfrac{1}{3x^3}+\dfrac{1}{2x^2}-\dfrac{1}{\sqrt3}\tan^{-1}\dfrac{2x+1}{\sqrt3}.
\]

(b)
\[
\int \dfrac{dx}{X^{p+1}}=\dfrac{1+2x}{3pX^p}+\dfrac{2(2p-1)}{3p}\int \dfrac{dx}{X^p}.
\]

\item Show that if $p$ be a positive integer and $X\equiv x^2+x+1$,

(a)
\[
\int \dfrac{dx}{(x^2+x+1)^{p+1}}
=\dfrac{(1+2x)}{3}\left[\dfrac{1}{pX^p}+\dfrac{2p-1}{p(p-1)}\left(\dfrac23\right)\dfrac{1}{X^{p-1}}+\dfrac{(2p-1)(2p-3)}{p(p-1)(p-2)}\left(\dfrac23\right)^2\dfrac{1}{X^{p-2}}
+\dots+\dfrac{(2p-1)(2p-3)\dots3.1}{p(p-1)\dots2.1}\left(\dfrac23\right)^{p-1}\dfrac{1}{X}\right]
\]
\[
+\dfrac{(2p-1)(2p-3)\dots3.1}{p(p-1)(p-2)\dots2.1}\dfrac{2}{\sqrt3}\left(\dfrac23\right)^p\tan^{-1}\dfrac{2x+1}{\sqrt3}.
\]

(b) $\displaystyle\int \dfrac{dx}{(x^2+x+1)^{\frac{2n+1}{2}}}$, $n$ being a positive integer,
\[
=\dfrac{1+2x}{3}\left[\dfrac{1}{2n-1}\dfrac{2}{X^{\frac{2n-1}{2}}}+\dfrac{n-1}{(2n-1)(2n-3)}\left(\dfrac23\right)\dfrac{2^3}{X^{\frac{2n-3}{2}}}
+\dfrac{(n-1)(n-2)}{(2n-1)(2n-3)(2n-5)}\left(\dfrac23\right)^2\dfrac{2^5}{X^{\frac{2n-5}{2}}}
+\dots+\dfrac{(n-1)!}{(2n-1)(2n-3)\dots1}\left(\dfrac23\right)^{n-1}\dfrac{2^{2n-1}}{X^{\frac12}}\right].
\]

\end{enumerate}

\section*{Exercises --- Miscellaneous Reduction Formulae (Arts. 248--274)}

\begin{enumerate}[leftmargin=2em]

\item If $X\equiv ax^n+b$, obtain reduction formulae for the integral $u_{p,q}\equiv\displaystyle\int \dfrac{x^p}{X^q}\,dx$ of the forms,
\begin{enumerate}[label=(\roman*)]
\item $Au_{p,q}+Bu_{p-n,q}+R=0$,
\item $A'u_{p,q}+B'u_{p,q-1}+R'=0$,
\end{enumerate}
where $A,B,A',B'$ are constants and $R,R'$ are algebraic functions of $x$. \source{Math. Trip., 1896.}

\item Prove that
\begin{enumerate}[label=(\alph*)]
\item $\displaystyle\int \cos^{2n}\phi\,d\phi=\dfrac{1}{2n}\tan\phi\cos^{2n}\phi+\left(1-\dfrac{1}{2n}\right)\int \cos^{2n-2}\phi\,d\phi$, \source{Trinity, 1891.}
\item $\displaystyle\int \sec^{2n+1}\phi\,d\phi=\dfrac{1}{2n}\tan\phi\sec^{2n-1}\phi+\left(1-\dfrac{1}{2n}\right)\int \sec^{2n-1}\phi\,d\phi$. \source{I.C.S., 1886.}
\end{enumerate}

\item Prove that
\[
\int (a^2+x^2)^{\frac{2n+1}{2}}\,dx=\dfrac{x}{2n+2}(a^2+x^2)^{\frac{2n+1}{2}}+\dfrac{2n+1}{2n+2}a^2\int (a^2+x^2)^{\frac{2n-1}{2}}\,dx. \source{I.C.S., 1886.}
\]

\item Investigate a formula of reduction applicable to
\[
\int x^m(1+x^2)^{\frac{n}{2}}\,dx,
\]
where $m$ and $n$ are positive integers, and complete the integration if $m=5$, $n=7$. \source{St. John's Coll., 1881.}

\item If $\phi(n)\equiv a^{2n-1}\displaystyle\int_0^\infty \dfrac{dx}{(a^2+x^2)^n}$, prove that $\phi(n)=\dfrac{2n-3}{2n-2}\phi(n-1)$. \source{R.P.}

\item Investigate formulae of reduction for
\begin{enumerate}[label=(\alph*)]
\item $\displaystyle\int \dfrac{dx}{(x^2+a^2)^{\frac{n}{2}}}$.
\item $\displaystyle\int x^n(a+bx)^{p+\frac12}\,dx$.
\item $\displaystyle\int \dfrac{x^m}{(a^2+x^2)^{\frac12}}\,dx$.
\item $\displaystyle\int \dfrac{x^n}{(a^3+x^3)^{\frac{n}{3}}}\,dx$.
\item $\displaystyle\int \dfrac{x^m}{(x^3-1)^{\frac13}}\,dx$.
\item $\displaystyle\int x^{2n}(x^2+a^2)^{\frac{2p+1}{2}}\,dx$,
\end{enumerate}
and obtain the value of $\displaystyle\int x^8(x^3-1)^{-\frac13}\,dx$. \source{Colleges, Camb.}

\item Investigate a formula of reduction for
\[
\int \dfrac{x^{2n+1}}{(1-x^2)^{\frac12}}\,dx,
\]
and by means of this integral show that
\[
\dfrac{1}{2n+2}+\dfrac12\cdot\dfrac{1}{2n+4}+\dfrac{1.3}{2.4}\cdot\dfrac{1}{2n+6}+\dfrac{1.3.5}{2.4.6}\cdot\dfrac{1}{2n+8}+\dots\ \text{ad inf.}
=\dfrac{2.4.6\dots2n}{3.5.7\dots(2n+1)}.
\]
Sum also the series
\[
\dfrac{1}{2n+1}+\dfrac12\cdot\dfrac{1}{2n+3}+\dfrac{1.3}{2.4}\cdot\dfrac{1}{2n+5}+\dfrac{1.3.5}{2.4.6}\cdot\dfrac{1}{2n+7}+\dots\ \text{ad inf.} \source{Math. Trip., 1897.}
\]

\item Find a reduction formula for
\[
\int_1^x \dfrac{x^n\,dx}{\sqrt{x-1}}.
\]
Show that
\[
\dfrac{2.4.6\dots2n}{3.5.7\dots(2n+1)}\left[1+\dfrac12x+\dfrac{1.3}{2.4}x^2+\dots+\dfrac{1.3\dots(2n-1)}{2.4\dots2n}x^n\right]
=1+\dfrac{a_1}{3}(x-1)+\dfrac{a_2}{5}(x-1)^2+\dots+\dfrac{1}{2n+1}(x-1)^n,
\]
where $a_1,a_2,\dots$ are the binomial coefficients. \source{St. John's, 1886.}

\item Prove that if $u_n\equiv\displaystyle\int_0^{\pi/2}\sin^{2n}x\,dx$, then
\[
u_n=\left(1-\dfrac{1}{2n}\right)u_{n-1}-\dfrac{1}{n2^{n+1}},
\]
and deduce
\[
\int_0^{\pi/2}\sin^{2n}x\,dx=-\dfrac{1}{2^{n+1}}\left\{\dfrac1n+\dfrac{2n-1}{n(n-1)}+\dfrac{(2n-1)(2n-3)}{n(n-1)(n-2)}+\dots\right\}
+\dfrac{(2n-1)(2n-3)\dots3}{2n(2n-2)\dots4.2}\cdot\dfrac{\pi}{4}. \source{Math. Trip., 1878.}
\]

\item Prove that
\[
\int_0^1 x^{4m+1}\sqrt{\dfrac{1-x^2}{1+x^2}}\,dx=\dfrac{1.3.5\dots(2m-1)}{2.4.6\dots2m}\cdot\dfrac{\pi}{4}-\dfrac{2.4.6\dots2m}{3.5.7\dots2m+1}\cdot\dfrac12.
\]

\item Find a reduction formula for
\[
\int e^{ax}\cos^nx\,dx,
\]
where $n$ is a positive integer, and evaluate
\[
\int e^{ax}\cos^4x\,dx. \source{Oxford, 1889.}
\]

\item Find formulae of reduction for
\begin{enumerate}[label=(\arabic*)]
\item $\displaystyle\int x^n\sin x\,dx$,
\item $\displaystyle\int e^{ax}\sin^nx\,dx$.
\end{enumerate}
Deduce from the latter a formula of reduction for
\[
\int \cos ax\sin^nx\,dx. \source{Colleges $\gamma$, 1890.}
\]

\item Show that
\[
(m+n)(m+n-2)\int \sin^m\theta\cos^n\theta\,d\theta
=(m-1)\sin^{m+1}\theta\cos^{n-1}\theta-(n-1)\sin^{m-1}\theta\cos^{n+1}\theta
\]
\[
+(m-1)(n-1)\int \sin^{m-2}\theta\cos^{n-2}\theta\,d\theta. \source{Trin. Coll. Camb., 1889.}
\]

\item Show that
\[
2^m\int \cos mx\cos^mx\,dx
=C+x+m\cdot\dfrac{\sin2x}{2}+\dfrac{m(m-1)}{1.2}\cdot\dfrac{\sin4x}{4}+\dots+\dfrac{\sin2mx}{2m},
\]
where $m$ is a positive integer. \source{Colleges $a$, 1885.}

\item Show that
\[
\int_0^{\pi/2}\sin^{2m}\theta\cos^{2m-1}\theta\,d\theta=\dfrac{(2m-2)(2m-4)\dots4.2}{(4m-1)(4m-3)\dots(2m+1)},
\]
$m$ being a positive integer. \source{Oxford, 1889.}

\item Evaluate the integral $\displaystyle\int_{-\pi/2}^{\pi/2}e^{-nx}\cos^mx\,dx$, $m$ being a positive integer. \source{Coll., 1886.}

\item Prove that if
\[
I_{m,n}\equiv\int \cos^mx\sin nx\,dx,
\]
\[
(m+n)I_{m,n}=-\cos^mx\cos nx+mI_{m-1,n-1}
\]
and
\[
\left[I_{m,m}\right]_0^{\pi/2}=\dfrac{1}{2^{m+1}}\left(\dfrac21+\dfrac{2^2}{2}+\dfrac{2^3}{3}+\dots+\dfrac{2^m}{m}\right). \source{Bertrand.}
\]

\item If $u_{m,n}\equiv\displaystyle\int_0^{\pi/2}\cos^mx\sin nx\,dx$, prove that
\[
u_{m,n}=\dfrac{1}{m+n}+\dfrac{m}{m+n}u_{m-1,n-1}.
\]
Hence find the value, when $m$ is a positive integer, of
\[
\int_0^{\pi/2}\cos^mx\sin2mx\,dx. \source{$\gamma$, 1887.}
\]

\item If $I_{m,n}\equiv\displaystyle\int \cos^mx\cos nx\,dx$, prove that
\[
I_{m,n}=-\dfrac{\cos^2nx}{m^2-n^2}\dfrac{d}{dx}\left(\dfrac{\cos^mx}{\cos nx}\right)+\dfrac{m(m-1)}{m^2-n^2}I_{m-2,n},
\]
and show that
\[
\left[I_{m,n}\right]_0^{\pi/2}=\dfrac{m(m-1)}{m^2-n^2}\left[I_{m-2,n}\right]_0^{\pi/2}.
\]

\item Prove that $\displaystyle\int_0^{\pi/2}\cos^nx\cos nx\,dx=\dfrac{\pi}{2^{n+1}}$, $n$ being a positive integer. \source{Bertrand.}

\item If $m$ and $n$ be positive integers, and if $m+n$ be even, prove that
\[
\int_0^{\pi/2}\cos^m\theta\cos n\theta\,d\theta=\dfrac{\pi}{2^{m+1}}\dfrac{m!}{\dfrac{m+n}{2}!\,\dfrac{m-n}{2}!}. \source{Colleges, 1882.}
\]

\item If $\displaystyle\int_0^{\pi/2}\cos^mx\cos nx\,dx$ be denoted by $f(m,n)$, show that
\[
f(m,n)=\dfrac{m}{m-n}f(m-1,n+1)=\dfrac{m}{m+n}f(m-1,n-1). \source{Oxford, 1890.}
\]

\item Prove that if $n$ be a positive integer, greater than unity,
\[
\int_0^{\pi/2}\cos^{n-2}x\sin nx\,dx=\dfrac{1}{n-1}. \source{Oxford I. P., 1889.}
\]

\item If $u_{m,n}\equiv\displaystyle\int x^m\text{cosec}^nx\,dx$, prove that
\[
(n-1)(n-2)u_{m,n}=(n-2)^2u_{m,n-2}+m(m-1)u_{m-2,n-2}
-x^{m-1}\{m\sin x+(n-2)x\cos x\}\text{cosec}^{n-1}x. \source{Math. Trip., 1896.}
\]

\item If $\displaystyle\int_0^\infty e^{-x}x^{n-1}\log x\,dx\equiv\phi(n)$,

prove that $\phi(n+2)-(2n+1)\phi(n+1)+n^2\phi(n)=0$. \source{R.P., St. John's Coll., 1881.}

\item Show that if $U_n\equiv\displaystyle\int_0^1 \dfrac{x^ne^x\,dx}{\sqrt{1-x}}$,
\[
2U_{n+1}+U_n(2n-1)-2nU_{n-1}=0. \source{Colleges $\beta$, 1887.}
\]

\item Prove that if $\phi(m)\equiv\displaystyle\int x^m(x^3+3ax+c)^{-\frac12}\,dx$,

then
\[
(2m-1)\phi(m)+3a(2m-3)\phi(m-2)+(2m-4)c\phi(m-3)
=2x^{m-2}(x^3+3ax+c)^{\frac12}. \source{Trinity, 1886.}
\]

\item If $u_m\equiv\displaystyle\int_0^{\pi/2}e^{-nx}\cos^mx\,dx$,

prove that $(m^2+n^2)u_m=m(m-1)u_{m-2}+n$. \source{Oxford I. P., 1900.}

\item Prove that if $I_m\equiv\displaystyle\int_0^{\pi/2}\dfrac{\sin^mx\,dx}{(1-k^2\sin^2x)^{\frac12}}$,

then $(m-1)k^2I_m-(m-2)(1+k^2)I_{m-2}+(m-3)I_{m-4}=0$. \source{Trinity, 1890.}

\item Obtain a reduction formula for the integral
\[
I_n\equiv\int (a\cos^2\theta+2h\sin\theta\cos\theta+b\sin^2\theta)^{-n}\,d\theta
\]
in the form
\[
2(n+1)(ab-h^2)I_{n+2}-(2n+1)(a+b)I_{n+1}+2nI_n=-\dfrac{1}{2n}\dfrac{d^2I_n}{d\theta^2}. \source{Math. Trip., 1898.}
\]

\item Show that
\[
\int_0^{\pi/2}\dfrac{d\phi}{(1-e^2\sin^2\phi)^3}=\dfrac{8-8e^2+3e^4}{(1-e^2)^{\frac52}}\dfrac{\pi}{16},
\]
$e$ being less than unity. \source{St. John's Coll., 1885.}

\item If $I_n\equiv\displaystyle\int \dfrac{\sin nx}{\sin x}\,dx$, prove that $(n-1)(I_n-I_{n-2})=2\sin(n-1)x$,

and hence that
\[
\int_0^{\pi/2}\dfrac{\sin nx}{\sin x}=\dfrac{\pi}{2}, \text{ if $n$ be odd,}
\]
\[
=2\left(1-\dfrac13+\dfrac15-\dfrac17+\dots+(-1)^{\frac n2+1}\dfrac{1}{n-1}\right), \text{ if $n$ be even.}
\]

\item If $X\equiv a+bx^n+cx^{2n}$ and $I_{m,p}\equiv\displaystyle\int x^mX^p\,dx$,

prove the existence of reduction formulae of the nature of
\begin{align*}
\text{(i)}\quad & x^{m+1}X^{p+1}=A_1I_{m,p}+B_1I_{m+n,p}+C_1I_{m+2n,p};\\
\text{(ii)}\quad & x^{m-2n+1}X^{p+1}=A_2I_{m,p}+B_2I_{m-n,p}+C_2I_{m-2n,p};\\
\text{(iii)}\quad & x^{m+1}X^p=A_3I_{m,p}+B_3I_{m+n,p-1}+C_3I_{m+2n,p-1};\\
\text{(iv)}\quad & x^{m+1}X^p=A_4I_{m,p}+B_4I_{m,p-1}+C_4I_{m+n,p-1};\\
\text{(v)}\quad & x^{m-n+1}X^{p+1}=A_5I_{m,p}+B_5I_{m-n,p}+C_5I_{m-n,p+1};
\end{align*}
and find the values of the fifteen constants.

\item Show that
\[
\int \dfrac{dx}{(a+bx^2+cx^4)^p}
\]
can be reduced to the integration of
\[
\int \dfrac{dx}{a+bx^2+cx^4} \quad\text{and}\quad \int \dfrac{x^2\,dx}{a+bx^2+cx^4} \quad (b>0,\ b^2>4ac),
\]
and integrate these expressions; $p$ being integral. \source{Bertrand.}

\item Show that, if $u_{m,n}\equiv\displaystyle\int \dfrac{x^m}{(\log x)^n}\,dx$,
\[
(n-1)u_{m,n}=-\dfrac{x^{m+1}}{(\log x)^{n-1}}+(m+1)u_{m,n-1}. \source{Oxford, I. P., 1889.}
\]

\item Find reduction formulae for
\begin{enumerate}[label=(\alph*)]
\item $\displaystyle\int \tanh^nx\,dx$;
\item $\displaystyle\int \dfrac{x}{\sin^nx}\,dx$;
\item $\displaystyle\int \dfrac{dx}{(a+be^x+ce^{-x})^n}$.
\end{enumerate}

\item If $I_m\equiv\displaystyle\int \dfrac{x^m\,dx}{\sqrt X}$, where $X\equiv ax^2+2bx+c$, show that
\[
amI_m+(2m-1)bI_{m-1}+(m-1)cI_{m-2}=x^{m-1}\sqrt X. \source{$\beta$, 1891.}
\]

\item Establish a reduction formula for
\[
I_n\equiv\int \dfrac{dx}{(Ax^2+B)^n\sqrt X},
\]
where $X\equiv ax^4+bx^2+c$, in the form
\[
\dfrac{x\sqrt X}{(Ax^2+B)^{n-1}}=\lambda I_n+\mu I_{n-1}+\nu I_{n-2}+\xi I_{n-3},
\]
showing that
\begin{align*}
\lambda&=(2n-2)\dfrac{B}{A^2}(A^2c-ABb+B^2a),\\
\mu&=-(2n-3)\dfrac{1}{A^2}(A^2c-2ABb+3B^2a),\\
\nu&=-(2n-4)\dfrac{1}{A^2}(Ab-3Ba),\\
\xi&=-(2n-5)\dfrac{1}{A^2}a.
\end{align*}

\item Show that, if
\[
I_{m,n}\equiv\int \sin^m\theta\cos n\theta\,d\theta, \qquad J_{m,n}\equiv\int \sin^m\theta\sin n\theta\,d\theta,
\]
then
\[
J_{m-1,n+1}=\left(1-\dfrac{n}{m}\right)I_{m,n}, \qquad I_{m-1,n+1}=-\left(1-\dfrac{n}{m}\right)J_{m,n},
\]
where $m$ is a positive integer; and point out how these results can be used to find the values of $I_{m,n}$ and $J_{m,n}$. \source{C. S., 1896.}

\item If $T$ be a function of $x$ such that
\[
\left(\dfrac{dT}{dx}\right)^2=A+3BT+3CT^2+DT^3,
\]
prove that
\[
\dfrac{d}{dx}\left(\dfrac{1}{T^{n-1}}\dfrac{dT}{dx}\right)=-\dfrac{(n-1)A}{T^n}-\dfrac32(2n-3)\dfrac{B}{T^{n-1}}-3\dfrac{(n-2)C}{T^{n-2}}-\dfrac{(2n-5)D}{2T^{n-3}}.
\]
Apply the result to investigate a reduction formula for
\[
\int \dfrac{dx}{T^n}.
\]
By a consideration of the case where $C=0$, $D=0$ (or in any other way), obtain a reduction formula for
\[
\int \dfrac{dx}{(a+2bx+cx^2)^n}. \source{I. C. S., 1897.}
\]

\item Prove that
\[
\int_0^\infty e^{-x^2}x^{2n}\,dx=\sqrt\pi\dfrac{1.3.5\dots(2n-1)}{2.4.6\dots2n}\int_0^\infty e^{-x^2}x^{2n+1}\,dx,
\]
where $n$ is a positive integer. \source{Colleges $a$, 1890.}

\item If $u_n\equiv\displaystyle\int_a^b x^n\{(b-x)(x-a)\}^{-\frac12}\,dx$,

show that $2nu_n=(2n-1)(a+b)u_{n-1}-2(n-1)abu_{n-2}$. \source{Oxford I. Pub., 1912.}

\item By applying the substitution $x=a\cos^2\theta+b\sin^2\theta$ (or otherwise), prove that the definite integral
\[
\int_a^b \dfrac{x^n\,dx}{\sqrt{(x-a)(b-x)}}
\]
is a rational integral function of $a$ and $b$ when $n$ is an integer; and evaluate it when $n=3$. \source{Oxf. I. P., 1913.}

\item If $u_n\equiv\displaystyle\int \dfrac{\cos2nx}{\sin^2x}\,dx$,

obtain a formula of reduction connecting $u_n$ and $u_{n-1}$.

Hence, or otherwise, evaluate
\[
\int_x^{\pi/2}\dfrac{\cos2nx}{\sin^2x}\,dx,
\]
where $n$ is a positive integer and $\pi/2>x>0$.

Consider the case when the lower limit is negative. \source{Oxf. I. P., 1915.}

\item By multiplying the inequality $1\geqslant2\sin x-\sin^2x$ by $\sin^{2n-1}x$ and by $\sin^{2n}x$, and integrating between $0$ and $\tfrac12\pi$, show that
\[
\left\{\dfrac{(4n+3)(2n+1)}{4n+4}\cdot\dfrac\pi2\right\}^{\frac12}>\dfrac{2.4\dots2n}{1.3\dots(2n-1)}>\left\{\dfrac{2n(2n+1)}{4n+1}\cdot\pi\right\}^{\frac12}. \source{Math. Trip. I., 1915.}
\]

\item The expression
\[
\dfrac{1-a}{(1-a\sin^2\theta)^{\frac12}}-(1-a\sin^2\theta)^{\frac12},
\]
where $1>a>0$, is expanded in ascending powers of $a$, and the coefficient of $a^n$ is denoted by $u_n$. Prove that
\[
\int_0^{\pi/2}u_n\,d\theta=0. \source{Math. Trip. I., 1916.}
\]

\item If $s_n\equiv\displaystyle\int \dfrac{\sin(2n-1)x}{\sin x}\,dx$, $v_n\equiv\displaystyle\int \dfrac{\sin^2nx}{\sin^2x}\,dx$,

prove the reduction formulae
\[
n(s_{n+1}-s_n)=\sin2nx, \qquad v_{n+1}-v_n=s_{n+1};
\]
and show that if $v_n$ be taken between the limits $0$ and $\tfrac12\pi$, its value is $\tfrac12n\pi$, where $n$ is an integer. \source{Math. Trip. I., 1914.}

\item If $\lambda^2\equiv\cos^2\phi/a^2+\sin^2\phi/b^2$, find $\displaystyle\int_0^{\pi/2}\dfrac{d\phi}{\lambda^2}$,

and prove that
\[
\int_0^{\pi/2}\dfrac{d\phi}{\lambda^6}=\pi ab\{3(a^4+b^4)+2a^2b^2\}/16,
\]
and that
\[
\int_0^1 \dfrac{3\mu^2-1}{2}\dfrac{d\mu}{(1-e^2\mu^2)^{\frac32}}=\dfrac{e^2}{3(1-e^2)^{\frac32}}. \source{$\epsilon$, 1883.}
\]

\item If $U_n\equiv\displaystyle\int \sin^mx(a+b\cos x)^{-n}\,dx$, prove that $U_n$ can be calculated by means of a reduction formula of the nature
\[
AU_n+BU_{n-1}+CU_{n-2}=\sin^{m+1}x(a+b\cos x)^{-n+1};
\]
and determine the constants $A$, $B$, $C$.

\item Prove that
\[
\int_c^\infty \dfrac{dx}{(a^2-2cx+x^2)^{n+1}}=\dfrac{2n!}{n!\,n!}\dfrac{\pi}{\lambda^{n+\frac12}},
\]
where $\lambda$ denotes $4(a^2-c^2)$ and is supposed positive. \source{Trin., 1887.}

\end{enumerate}

\newpage

\begin{center}
    {\LARGE \bfseries \color{accent} Chapter VIII: Form $\displaystyle\int F(x,\sqrt{R})\,dx$, where $R$ is Quadratic --- Exercises}\\
    \vspace{0.2cm}
    {\color{accent}\rule{0.55\linewidth}{0.8pt}}
\end{center}
\vspace{0.35cm}

\section*{Exercises --- Cases I--III: $X$ and $Y$ Linear or Quadratic (Arts. 275--290)}

Integrate the following expressions:

\begin{enumerate}[leftmargin=2em]

\item $\dfrac{1}{x\sqrt{x+1}}$, $\dfrac{1}{(x-1)\sqrt{x+2}}$, $\dfrac{x+1}{(x-1)\sqrt{x+2}}$, $\dfrac{x^2+x+1}{(x+2)\sqrt{x-1}}$.

\item $\dfrac{1}{(x^2+1)\sqrt{x}}$, $\dfrac{1}{(x^2+2x+2)\sqrt{x+1}}$, $\dfrac{x}{(x^2+2x+2)\sqrt{x+1}}$, $\dfrac{x^2+1}{(x^2+2x+2)\sqrt{x+1}}$.

\item $\dfrac{1}{x\sqrt{x^2+1}}$, $\dfrac{1}{(x+1)\sqrt{x^2+1}}$, $\dfrac{x}{(x+1)\sqrt{x^2+1}}$, $\dfrac{x^2+x+1}{(x+1)\sqrt{x^2+2x+3}}$.

\item Prove that
\[
\int \dfrac{dx}{(x+c)\sqrt{x}}=\dfrac{2}{\sqrt{c}}\tan^{-1}\sqrt{\dfrac{x}{c}} \quad\text{or}\quad \dfrac{1}{\sqrt{-c}}\log_e\dfrac{\sqrt{x}-\sqrt{-c}}{\sqrt{x}+\sqrt{-c}},
\]
according as $c$ is positive or negative. \source{C. S., 1904.}

\item Integrate $\dfrac{\sqrt{\sin\theta}}{\cos\theta(\cos\theta+\sin\theta)\sqrt{2\cos\theta+3\sin\theta}}$.

\item Integrate $\dfrac{\sqrt{\sin\theta}+\sqrt{\cos\theta}}{(a^2\cos^2\theta-b^2\sin^2\theta)\sqrt{\cos\left(\theta-\dfrac{\pi}{4}\right)}}$ \quad $(a>b>0)$.

\item Integrate $\dfrac{\tan2\theta}{\sqrt{\cos^6\theta+\sin^6\theta}}$. \qquad Integrate $\dfrac{(x^3+1)^2}{(x-1)\sqrt{x^2+1}}$.

\item Integrate
\[
\dfrac{1}{\sqrt{(a-x)(x-b)}}, \qquad \dfrac{1}{(a-x)\sqrt{x-b}}, \qquad \dfrac{1}{(x-b)\sqrt{a-x}}, \qquad \dfrac{1}{(1\pm x)\sqrt{1\pm x^2}}.
\]
\source{St. John's, 1890.}

\item Integrate $\dfrac{1}{x\sqrt{(x-a)(x-b)}}$. \source{Colleges, 1876.}

\item Show that
\[
\int \dfrac{x-3}{x\sqrt{x^2-3x+2}}\,dx=\cosh^{-1}(2x-3)+\dfrac{3}{\sqrt2}\cosh^{-1}\left(\dfrac{4-3x}{x}\right).
\]
\source{St. John's, 1883.}

\item Integrate $\displaystyle\int \dfrac{x-\alpha}{x-\beta}\dfrac{dx}{\sqrt{x^2+2px+q}}$ \quad $(p^2>q)$. \source{$a$, 1887.}

\item Integrate
\begin{enumerate}[label=(\roman*)]
\item $\displaystyle\int \dfrac{dx}{x\sqrt{x^2+x+1}}$, \source{Coll., 1879.}
\item $\displaystyle\int \dfrac{dx}{x\sqrt{x^2-a^2}}$, \source{$\epsilon$, 1883.}
\item $\displaystyle\int \dfrac{dx}{(1+x^2)\sqrt{1-x^2}}$. \source{I. C. S., 1888.}
\end{enumerate}

\item Integrate
\begin{enumerate}[label=(\roman*)]
\item $\displaystyle\int \dfrac{dx}{(1-x)\sqrt{1-x^2}}$ (in two ways). \source{Trinity, 1890.}
\item $\displaystyle\int \dfrac{dx}{(1-2x)\sqrt{1+4x}}$. \source{Coll., 1892.}
\end{enumerate}

\item Integrate
\[
\int \dfrac{dx}{(x-\lambda)\sqrt{x-\mu}}, \qquad \int \dfrac{x^3\,dx}{x^2+\lambda^2}, \qquad \int \dfrac{dx}{(x+1)(x+2)(x+3)}.
\]
\source{Math. Trip., Pt. I., 1920.}

\end{enumerate}

\section*{General Examples}

\begin{enumerate}[leftmargin=2em]

\item Obtain the following integrals:
\begin{enumerate}[label=(\roman*)]
\item $\displaystyle\int (1+x)^{-1}x^{-\frac12}\,dx$.
\item $\displaystyle\int (1+x)^{-1}(1+2x)^{-\frac12}\,dx$.
\item $\displaystyle\int x^{-1}(2-3x+x^2)^{-\frac12}\,dx$.
\item $\displaystyle\int (1+x)^{-1}(1+x+x^2)^{-\frac12}\,dx$.
\item $\displaystyle\int \dfrac{\sqrt{1+x+x^2}}{1+x}\,dx$.
\item $\displaystyle\int \dfrac{x^2+x-1}{(x+1)\sqrt{x^2-1}}\,dx$.
\item $\displaystyle\int \dfrac{dx}{x\sqrt{a^n+x^n}}$.
\item $\displaystyle\int \dfrac{1-x}{1+x}\dfrac{1}{\sqrt{x+x^2+x^3}}\,dx$.
\end{enumerate}

\item Integrate
\begin{enumerate}[label=(\roman*)]
\item $\displaystyle\int \dfrac{dx}{(x+2)\sqrt{x^2-1}}$.
\item $\displaystyle\int \dfrac{\sqrt{x^2-1}}{x+2}\,dx$. \source{Barnes Schol., 1887.}
\end{enumerate}

\item Show that
\[
\int \dfrac{dx}{(x-p)\sqrt{a+2bx+cx^2}}=\dfrac{1}{\{-(a+2bp+cp^2)\}^{\frac12}}\sin^{-1}\left\{\dfrac{(a+bx)+p(b+cx)}{(x-p)\sqrt{b^2-ac}}\right\}
\]
where $p$ lies between the roots of $a+2bx+cx^2=0$, supposed real. \source{Trinity, 1886 and 1891.}

\item Show that
\[
\int \dfrac{dx}{(x^2+a^2)\sqrt{x^2+b^2}}=\dfrac{1}{a\sqrt{b^2-a^2}}\cos^{-1}\dfrac{a}{b}\sqrt{\dfrac{x^2+b^2}{x^2+a^2}}, \quad \text{if } a<b,
\]
\[
\text{and} \quad =\dfrac{1}{a\sqrt{a^2-b^2}}\cosh^{-1}\dfrac{a}{b}\sqrt{\dfrac{x^2+b^2}{x^2+a^2}}, \quad \text{if } a>b.
\]

\item Prove that
\[
\int \dfrac{(x+1)\,dx}{(2x^2-2x+1)\sqrt{3x^2-2x+1}}=\cosh^{-1}\sqrt{\dfrac{3x^2-2x+1}{2x^2-2x+1}}
+2\cos^{-1}\dfrac{1}{\sqrt2}\sqrt{\dfrac{3x^2-2x+1}{2x^2-2x+1}}.
\]

\item Integrate
\begin{enumerate}[label=(\roman*)]
\item $\displaystyle\int \dfrac{(3x+4)\,dx}{(5x^2+8x)\sqrt{4x^2-2x+1}}$,
\item $\displaystyle\int \dfrac{dx}{(x^2+2ax+b^2)\sqrt{x^2+2ax+c^2}}$,
\end{enumerate}
where $a<b<c$.

\item Integrate
\begin{enumerate}[label=(\roman*)]
\item $\displaystyle\int \dfrac{x\,dx}{(a^2+b^2-x^2)\sqrt{(a^2-x^2)(x^2-b^2)}}$. \source{St. John's, 1888.}
\item $\displaystyle\int \dfrac{(x+b)\,dx}{(x^2+a^2)\sqrt{x^2+c^2}}$ \quad $(a>c)$. \source{St. John's, 1889.}
\item $\displaystyle\int \dfrac{d\theta}{\sin\theta\sqrt{a\cos^2\theta+b\sin^2\theta+c}}$. \source{Trinity, 1888.}
\end{enumerate}

\item Find the values of
\begin{enumerate}[label=(\roman*)]
\item $\displaystyle\int \dfrac{\sin x\,dx}{(\cos x+\cos\alpha)\sqrt{(\cos x+\cos\beta)(\cos x+\cos\gamma)}}$. \source{$\gamma$, 1890.}
\item $\displaystyle\int \dfrac{dx}{\cos(x+\alpha)\sqrt{\cos(x+\beta)\cos(x+\gamma)}}$. \source{$\gamma$, 1890.}
\end{enumerate}

\item Integrate
\[
\int \dfrac{a^2-x^2}{(a^2-ax+x^2)(a^4+a^2x^2+x^4)^{\frac12}}\,dx,
\]
transforming by the substitution
\[
x^2+ax+a^2=y^2(x^2-ax+a^2). \source{$a$, 1884.}
\]

\item Integrate
\begin{enumerate}[label=(\roman*)]
\item $\displaystyle\int \dfrac{(8x-13)\,dx}{(3x^2-10x+9)\sqrt{-x^2+10x-13}}$.
\item $\displaystyle\int \dfrac{dx}{(x-1)(x-2)\sqrt{(x-3)(x-4)}}$. \source{Coll., 1892.}
\item $\displaystyle\int \dfrac{(x+9)\,dx}{(x^2-5x+4)\sqrt{x^2-2x+2}}$.
\item $\displaystyle\int \dfrac{(x-a)\,dx}{(x-b)(x-c)(x-d)\sqrt{x-e}}$.
\item $\displaystyle\int \dfrac{(x+3)\,dx}{(x^2+x+1)\sqrt{x^2+x+2}}$.
\end{enumerate}

\item Integrate
\begin{enumerate}[label=(\roman*)]
\item $\displaystyle\int \dfrac{x^4-1}{x^2\sqrt{x^4+x^2+1}}\,dx$; \source{Coll., 1901.}
\item $\displaystyle\int \dfrac{x^{-1}(x^2-1)}{\sqrt{x^4+x^2+1}}\,dx$.
\end{enumerate}

\item Show that
\[
\int \sqrt{\dfrac{\sin(x-a)}{\sin(x+a)}}\,dx=\cos a\cos^{-1}\left(\dfrac{\cos x}{\cos a}\right)-\sin a\cosh^{-1}\left(\dfrac{\sin x}{\sin a}\right).
\]
\source{Coll., 1901.}

\item Integrate
\begin{enumerate}[label=(\roman*)]
\item $\displaystyle\int \dfrac{1-2x^2}{1+2x^2}\sqrt{\dfrac{1+x^2}{1-x^2}}\,x\,dx$. \source{R. P.}
\item $\displaystyle\int \dfrac{dx}{(1+x^4)\sqrt{\left(\sqrt{1+x^4}-x^2\right)}}$. \source{R. P.; Euler, C.I., vol. iv.}
\item $\displaystyle\int \dfrac{dx}{(a^2-ax-x^2)\sqrt{ax+x^2}}$. \source{Oxf. I. P., 1900.}
\end{enumerate}

\item Evaluate the integral
\[
\int \dfrac{dx}{(a^2-\tan^2x)(b^2-\tan^2x)^{\frac12}}. \source{Math. Trip., 1886.}
\]

\item Prove that
\[
\int \dfrac{dx}{\sin x\sin^{\frac12}(2x+a)}=-\dfrac{1}{\sqrt{\sin a}}\cosh^{-1}\dfrac{\sin(x+a)}{\sin x}. \source{Coll., 1892.}
\]

\item Show how to integrate
\[
\int (fx+g)(ax^2+bx+c)^{\frac{n}{2}}\,dx,
\]
where $n$ is any positive or negative integer. \source{$a$, 1890.}

\item Show that $\displaystyle\int \dfrac{(a+bx)\,dx}{(a+2bx+cx^2)^{\frac32}}=\dfrac{x}{(a+2bx+cx^2)^{\frac12}}$. \source{Trinity, 1889.}

\item Prove by the substitution
\[
y^2=(ax^2+2bx+c)/(Ax^2+2Bx+C),
\]
where $A$ and $AC-B^2$ are positive, that the integral
\[
\int \dfrac{(Mx+N)\,dx}{(Ax^2+2Bx+C)\sqrt{ax^2+2bx+c}}
\]
becomes of the form
\[
P_1\int \dfrac{dy}{\sqrt{\lambda_1-y^2}}+P_2\int \dfrac{dy}{\sqrt{y^2-\lambda_2}},
\]
where $\lambda_1$ and $\lambda_2$ are the roots of the quadratic
\[
(a-\lambda A)(c-\lambda C)-(b-\lambda B)^2=0,
\]
and $P_1$, $P_2$ are definite constants.

Integrate completely the function
\[
\dfrac{x+3}{(2x^2-10x+17)\sqrt{4x^2-26x+49}}. \source{Math. Trip., 1891.}
\]

\item Prove $\displaystyle\int_0^a \sqrt{\tanh^2a-\tanh^2x}\,dx=\dfrac{\pi}{2}(1-\operatorname{sech}a)$. \source{$\beta$, 1889.}

\item Integrate $\displaystyle\int (x^2+b^2)^{-1}(x^2+a^2+b^2)^{-\frac12}\,dx$. \source{Ox. II. P., 1902.}

\item Integrate
\[
\int \dfrac{\sin^3x\,dx}{(1+\cos^2x)\sqrt{1+\cos^2x+\cos^4x}}, \source{St. John's, 1882.}
\]
and evaluate
\[
\int_{-1}^{1} \dfrac{dx}{(a^2+c^2x^2)\sqrt{1-x^2}}.
\]

\item Show that $\displaystyle\int \dfrac{dx}{(x-a)\sqrt{Ax^2+2Bx+C}}$ is transcendental unless $Aa^2+2Ba+C=0$. \source{J. M. Sch. Ox., 1904.}

Establish the results
\begin{enumerate}[label=(\roman*)]
\item $\displaystyle\int \dfrac{dx}{(x-1)\sqrt{x^2-4x+5}}=\dfrac{1}{\sqrt2}\cosh^{-1}\dfrac{\sqrt{2(x^2-4x+5)}}{x-1}$
\end{enumerate}
and
\begin{enumerate}[label=(\roman*)]
\setcounter{enumii}{1}
\item $\displaystyle\int \dfrac{dx}{(x-2)\sqrt{3+2x-2x^2}}=\sin^{-1}\dfrac{\sqrt{3+2x-2x^2}}{\sqrt7(x-2)}$.
\end{enumerate}
\source{Coll. $a$, 1890.}

\item Show that
\[
\int_{-1}^{+1} \dfrac{(1-ax)(1-bx)}{(1-2ax+a^2)(1-2bx+b^2)}\dfrac{dx}{\sqrt{1-x^2}}=\dfrac{\pi}{2}\dfrac{2-ab}{1-ab}. \source{$\delta$, 1884.}
\]

\item Describe the steps whereby the integral of a rational function of a single variable, $x$, can be obtained.

Prove that if the sign of summation refer to the suffixes 1, 2, 3 in cyclical order, the integral
\[
\int dx\sqrt{(x-a)(x-b)}\,\Sigma(c_2-c_3)\left[\dfrac{2(c_1-a)(c_1-b)}{(x-c_1)^2}-\dfrac{2c_1-a-b}{x-c_1}\right]
\]
is a certain constant multiple of
\[
(x-a)^{\frac32}(x-b)^{\frac32}(x-c_1)^{-1}(x-c_2)^{-1}(x-c_3)^{-1}. \source{Math. Trip., 1896.}
\]

\item Determine the degenerate form of the elliptic integral
\[
\int \dfrac{ds}{\sqrt{4(s-s_1)(s-s_2)(s-s_3)}}, \quad s_1>s_2>s_3,
\]
when $s_2$ is made to coincide with $s_1$ or with $s_3$. \source{Int. Arts, London.}

\item Prove, by means of the substitution $\dfrac{x-\beta}{x-\alpha}=y^2$, that
\[
\int \dfrac{dx}{(x-\gamma)\sqrt{(x-\alpha)(x-\beta)}}=\dfrac{-2}{\sqrt{(\alpha-\gamma)(\gamma-\beta)}}\tan^{-1}\sqrt{\dfrac{\alpha-\gamma}{\gamma-\beta}\cdot\dfrac{x-\beta}{x-\alpha}}
\]
\[
\text{or} \quad =\dfrac{2}{\sqrt{(\alpha-\gamma)(\beta-\gamma)}}\tanh^{-1}\sqrt{\dfrac{\alpha-\gamma}{\beta-\gamma}\cdot\dfrac{x-\beta}{x-\alpha}}. \source{Int. Arts, London.}
\]

\item Prove that
\[
\int_0^1 \dfrac{dx}{(1+x)(2+x)\sqrt{x(1-x)}}=\pi\left(\dfrac{1}{\sqrt2}-\dfrac{1}{\sqrt6}\right). \source{Math. Trip. I., 1912.}
\]

\item Show that the integral
\[
\int \dfrac{dx}{x\sqrt{3x^2+2x-1}}
\]
is rationalized by the assumption $x=(1+y^2)/(3-y^2)$, and hence, or otherwise, find its value.

Prove that if $m$ be a positive proper fraction, the value of the above integral when taken between limits $\dfrac13$ and $\dfrac{1}{2+m}$ is the same as when taken between limits $\dfrac{1}{2+m}$ and $\dfrac{1}{m(2+m)}$. \source{Math. Trip. I., 1910.}

\item Prove by means of the substitution
\[
\dfrac{a-x}{x-b}=\dfrac{a-d}{b-c}\cdot\dfrac{c-y}{y-d},
\]
that, if $m$ be any positive quantity, and $a>b>c>d$,
\[
\int_b^a \dfrac{\{(a-x)(x-c)\}^{m-1}}{\left\{\dfrac{(a-x)(x-d)}{a-d}+\dfrac{(x-b)(x-c)}{b-c}\right\}^m}\,dx
=\int_d^c \dfrac{\{(a-x)(c-x)\}^{m-1}}{\left\{\dfrac{(a-x)(x-d)}{a-d}+\dfrac{(b-x)(c-x)}{b-c}\right\}^m}\,dx.
\]
\source{Math. Trip., 1878.}

[See Wolstenholme's \textit{Mathematical Problems}, Numbers 1900--1903 for a group of similar examples.]

\item By the transformation $px+\dfrac{q}{x}=\sqrt{2pq}/z$, integrate
\[
\int \dfrac{px^2-q}{px^2+q}\dfrac{dx}{\sqrt{p^2x^4+q^2}}. \source{Cf. Euler, C.I., iv., p. 22.}
\]

\item Apply the transformation $x^2+\dfrac{1}{x^2}=\dfrac{2}{z^2}$ to integrate
\begin{enumerate}[label=(\roman*)]
\item $\displaystyle\int \dfrac{\sqrt{1+x^4}}{1-x^4}\,dx$,
\item $\displaystyle\int \dfrac{x^2\,dx}{(1-x^4)\sqrt{1+x^4}}$.
\end{enumerate}
\source{Euler, C.I., iv.}

\item Show that
\[
\int \dfrac{2\,d\theta}{\sin\theta\sqrt[4]{\cos2\theta}}=\tanh^{-1}\dfrac{\cos\theta}{\sqrt[4]{\cos2\theta}}+\tan^{-1}\dfrac{\cos\theta}{\sqrt[4]{\cos2\theta}}.
\]

\item Show that the transformation
\[
\{(a+bx^n)^\lambda-b^\lambda x^{\lambda n}\}=\left(\dfrac{x}{u}\right)^{\lambda n}
\]
will reduce the integration
\[
\int \dfrac{x^{m-1}\,dx}{(a+bx^n)\{(a+bx^n)^\lambda-b^\lambda x^{\lambda n}\}^{\frac{m}{\lambda n}}} \quad\text{to the form}\quad \dfrac{1}{a}\int \dfrac{u^{m-1}\,du}{1+b^\lambda u^{\lambda n}}.
\]
\source{Euler, C.I., iv., 53 and 56; Peacock, p. 305.}

\item[34.] (i) Show that
\[
\int \dfrac{e^x(2-x^2)}{(1-x)\sqrt{1-x^2}}\,dx=e^x\sqrt{\dfrac{1+x}{1-x}}, \source{Peacock, p. 309.}
\]
and (ii) integrate
\[
\int e^x\dfrac{1+nx^{n-1}-x^{2n}}{(1-x^n)\sqrt{1-x^{2n}}}\,dx.
\]

\item Integrate
\begin{enumerate}[label=(\roman*)]
\item $\displaystyle\int \dfrac{2-3x}{2+3x}\sqrt{\dfrac{1+x}{1-x}}\,dx$.
\item $\displaystyle\int \dfrac{\sin^5\theta+2\cos^5\theta}{\cos\theta\sin4\theta}\,d\theta$.
\end{enumerate}
\source{St. John's, 1881.}

\item Show that
\begin{enumerate}[label=(\roman*)]
\item $\displaystyle\int \sqrt{\dfrac{a^2-c^2x^2}{a^2-x^2}}\,\dfrac{dx}{x}=\dfrac12\log\left\{\dfrac{y-1}{y+1}\dfrac{(y+c)^c}{(y-c)^c}\right\}$,\\[0.3em]
where $y=\sqrt{\dfrac{a^2-c^2x^2}{a^2-x^2}}$.
\item $\displaystyle\int \dfrac{x^2+1}{x^2-1}\dfrac{dx}{\sqrt{1-ax^2+x^4}}=\dfrac{1}{\sqrt{a-2}}\cos^{-1}\dfrac{x\sqrt{a-2}}{x^2-1}$.
\end{enumerate}
\source{Hall, I.C., p. 325.}

\item If $F(x,y)$ be a rational algebraic function of $x$ and $y$, show that
\[
\int F\!\left(x,\sqrt{1+x^2}\right)\left(x+\sqrt{1+x^2}\right)^{\frac{p}{q}}dx
\]
may be integrated by the transformation $x=\sinh(q\log z)$.

\item Show that
\begin{enumerate}[label=(\roman*)]
\item $\displaystyle\int (\cos2\theta)^{\frac32}\cos\theta\,d\theta=\dfrac18\sin\theta(3+2\cos2\theta)\sqrt{\cos2\theta}+\dfrac{3}{8\sqrt2}\sin^{-1}(\sqrt2\sin\theta)$.
\item $\displaystyle\int_\alpha^\theta \dfrac{\sin\theta}{\sqrt{\sin^2\alpha-\sin^2\theta}}\,d\theta=\log\left(\dfrac{\cos\alpha}{\cos\theta+\sqrt{\sin^2\alpha-\sin^2\theta}}\right)$.
\item $\displaystyle\int_0^\pi e^{2x\cos\theta}\,d\theta=\pi\left[1+\dfrac{x^2}{(1!)^2}+\dfrac{x^4}{(2!)^2}+\dfrac{x^6}{(3!)^2}+\dots\right]$.
\end{enumerate}

\item Show that
\begin{enumerate}[label=(\roman*)]
\item $\displaystyle\int_0^1 x^{a\pm cx}\,dx=\dfrac{1}{a+1}\mp\dfrac{c}{(a+2)^2}+\dfrac{c^2}{(a+3)^3}\mp\dfrac{c^3}{(a+4)^4}+\dots$.
\item $\displaystyle\int_0^1 x^{\pm x^2}\,dx=1\pm\dfrac{1}{3^2}+\dfrac{1}{5^3}\mp\dfrac{1}{7^4}+\dots$.
\end{enumerate}
\source{Anglin.}

\item If $\phi(x)=a_0+\dfrac{1}{2}a_2x^2+\dfrac{1\cdot3}{2\cdot4}a_4x^4+\dots$, show that
\begin{enumerate}[label=(\roman*)]
\item $\displaystyle\dfrac{4}{\pi}\int_0^{\pi/2}\cos^2\theta\,\phi(\sin\theta)\,d\theta=1^2\cdot a_0+\dfrac13\left(\dfrac12\right)^2a_2+\dfrac13\left(\dfrac{1\cdot3}{2\cdot4}\right)^2a_4+\dots$.
\item $\displaystyle\dfrac{20}{9\pi}=\dfrac12\cdot1^2+\dfrac13\left(\dfrac12\right)^2+\dfrac14\left(\dfrac{1\cdot3}{2\cdot4}\right)^2+\dfrac15\left(\dfrac{1\cdot3\cdot5}{2\cdot4\cdot6}\right)^2+\dots$.
\item $\displaystyle\dfrac4\pi-1=\dfrac{1}{1^2}\left(\dfrac12\right)^2+\dfrac{1}{3^2}\left(\dfrac{1\cdot3}{2\cdot4}\right)^2+\dfrac{1}{5^2}\left(\dfrac{1\cdot3\cdot5}{2\cdot4\cdot6}\right)^2+\dots$.
\end{enumerate}
\source{Anglin.}

\item Integrate
\begin{enumerate}[label=(\roman*)]
\item $\displaystyle\int \dfrac{\sin2\theta\,d\theta}{\sqrt{\sin^4\theta+4\sin^2\theta\cos^2\theta+2\cos^4\theta}}$.
\item $\displaystyle\int \dfrac{(z^2-1)\,dz}{(z^2+z+1)\sqrt{z^4+4z^3+4z^2+4z+1}}$.
\end{enumerate}

\item If $J$ be the Jacobian of two quadratic functions of $x$, viz.
\[
u_1\equiv a_1x^2+2b_1x+c_1, \qquad u_2\equiv a_2x^2+2b_2x+c_2,
\]
\[
J\equiv\begin{vmatrix} 2(a_1x+b_1), & 2(b_1x+c_1) \\ 2(a_2x+b_2), & 2(b_2x+c_2) \end{vmatrix},
\]
show that if $u_1=0$, $u_2=0$ have no positive roots, then
\[
\int_0^\infty \dfrac{J}{u_1u_2}\,dx=2\log\dfrac{a_1c_2}{a_2c_1}.
\]

\item By means of the identity
\[
\int_0^{\pi/2}(a+\sin^2x)^n\cos x\,dx=\int_0^{\pi/2}(1+a-\sin^2x)^n\sin x\,dx,
\]
prove that
\begin{multline*}
a^n+{}^nC_1\dfrac{a^{n-1}}{3}+{}^nC_2\dfrac{a^{n-2}}{5}+{}^nC_3\dfrac{a^{n-3}}{7}+\dots\\
=(1+a)^n-2\dfrac{n}{3}(1+a)^{n-1}+2^2\dfrac{n(n-1)}{3\cdot5}(1+a)^{n-2}\\
-2^3\dfrac{n(n-1)(n-2)}{3\cdot5\cdot7}(1+a)^{n-3}+\dots.
\end{multline*}
\source{Wolstenholme, \textit{Problems}, No. 1929; Wiggins, E. Times, No. 13323.}

\item Show that
\begin{enumerate}[label=(\roman*)]
\item
\begin{multline*}
a^n+{}^nC_1\dfrac{1}{p+2}a^{n-1}+{}^nC_2\dfrac{1\cdot3}{(p+2)(p+4)}a^{n-2}\\
+{}^nC_3\dfrac{1\cdot3\cdot5}{(p+2)(p+4)(p+6)}a^{n-3}+\dots\\
=(1+a)^n-{}^nC_1\dfrac{p+1}{p+2}(1+a)^{n-1}+{}^nC_2\dfrac{(p+1)(p+3)}{(p+2)(p+4)}(1+a)^{n-2}\\
-{}^nC_3\dfrac{(p+1)(p+3)(p+5)}{(p+2)(p+4)(p+6)}(1+a)^{n-3}+\dots.
\end{multline*}
\item
\begin{multline*}
a^n+{}^nC_1\dfrac{1}{p+2}a^{n-1}b+{}^nC_2\dfrac{1\cdot3}{(p+2)(p+4)}a^{n-2}b^2\\
+{}^nC_3\dfrac{1\cdot3\cdot5}{(p+2)(p+4)(p+6)}a^{n-3}b^3+\dots\\
=(a+b)^n-{}^nC_1\dfrac{p+1}{p+2}(a+b)^{n-1}b+{}^nC_2\dfrac{(p+1)(p+3)}{(p+2)(p+4)}(a+b)^{n-2}b^2\\
-{}^nC_3\dfrac{(p+1)(p+3)(p+5)}{(p+2)(p+4)(p+6)}(a+b)^{n-3}b^3+\dots.
\end{multline*}
\end{enumerate}

\item (i) Integrate
\[
\int \dfrac{2x^3-1}{x^6+2x^3-x^2+1}\,dx. \source{Ox. I. P., 1903.}
\]
(ii) Integrate
\[
\int \dfrac{(3x^4-1)\,dx}{x^8+2x^4-16x^2+1}.
\]
(iii) Prove that
\[
\int_{-\pi/4}^{\pi/4} \dfrac{(1-\sin\theta\cos\theta+\sin^2\theta)\cos^2\theta\,d\theta}{(1+\sin\theta\cos\theta)^3}=\dfrac{10\pi}{9}\sqrt3.
\]

\item (i) Show that
\[
\int_0^a x^{n-1}(1-a+x)^{-n-1}\,dx=\dfrac1n\dfrac{a^n}{1-a}.
\]
(ii) Evaluate
\[
\int_a^{a+h} \dfrac{(a+h-x)^{n-1}}{(n-1)!}\dfrac{d^n}{dx^n}f(x)\,dx.
\]
How could your result be applied to the summation of series? \source{$a$, 1886.}

\item Discuss the integration of
\[
\int f\left[x,\ \left(\dfrac{a+bx}{a'+b'x}\right)^{\frac{m}{n}},\ \left(\dfrac{a+bx}{a'+b'x}\right)^{\frac{p}{q}},\ \left(\dfrac{a+bx}{a'+b'x}\right)^{\frac{r}{s}}\right]dx,
\]
where $f$ denotes a rational integral algebraic function of the quantities indicated. \source{Lacroix, C.I., ii., p. 35.}

\item If $F(x)$ be a rational integral algebraic function of $x$, show that $\displaystyle\int_{-1}^{+1}\dfrac{F(x)}{\sqrt{1-x^2}}=\pi k$, where $k$ is the coefficient of $\dfrac1a$ in the product
\[
F(a)\left[\dfrac1a+\dfrac12\cdot\dfrac{1}{a^3}+\dfrac{1\cdot3}{2\cdot4}\cdot\dfrac{1}{a^5}+\dots\right]; \source{St. John's, 1891.}
\]
or where $k$ is the constant term in the expansion of $\dfrac{xF(x)}{(x^2-1)^{\frac12}}$. \source{Coll., 1892.}

\item If $f(x)$ be an arbitrary algebraic polynomial of degree $n-1$, and
\[
P_n(x)\equiv A\dfrac{d^n}{dx^n}(x-a)^n(x-b)^n,
\]
where $A$ is a constant, then
\[
\int_a^b f(x)P_n(x)\,dx=0. \source{Lond. Univ.}
\]

\item Prove that
\[
\int_0^a \dfrac{x\,dx}{\cos x\cos(a-x)}=\dfrac{a}{\sin a}\log\sec a. \source{Coll., 1896.}
\]

\item Show that if $a$ be less than unity,
\[
\int_0^\pi \dfrac{x\sin x\,dx}{1+a^2\cos^2x}=\pi\dfrac{\tan^{-1}a}{a}. \source{$a$, 1891.}
\]

\item Integrate
\begin{enumerate}[label=(\roman*)]
\item $\displaystyle\int \dfrac{b^4x+x^5}{(b^4-x^4)^2}\sin^{-1}\dfrac{x^2}{b^2}\,dx$. \source{$\delta$, 1881.}
\item $\displaystyle\int \dfrac{d\phi}{\cos\phi}\sqrt{1-\dfrac12\sin^2\phi}$. \source{St. John's, 1885.}
\end{enumerate}

\item From the definition of a Bessel's function, viz.
\[
J_n(x)=\dfrac{x^n}{2^n\Gamma(n+1)}\left[1-\dfrac{x^2}{2(2n+2)}+\dfrac{x^4}{2\cdot4(2n+2)(2n+4)}-\dots\right],
\]
derive the results
\[
\dfrac{\sin x}{x}=\int_0^{\pi/2}J_0(x\sin\theta)\sin\theta\,d\theta,
\]
\[
\dfrac{1-\cos x}{x}=\int_0^{\pi/2}J_1(x\sin\theta)\,d\theta. \source{Coll., 1896.}
\]

\item Integrate
\begin{enumerate}[label=(\roman*)]
\item $\dfrac{x^3}{(x+1)^3(x^3-1)}$.
\item $\dfrac{1}{\sqrt{1+3\sin x\cos x+2\sin^2x\cos^2x}}$.
\item $\dfrac{1}{\sqrt{(1+\sin x)(2+\sin x)}}$. \source{Math. Trip., 1897.}
\end{enumerate}

\item Show that
\[
\int \sin^mx\sin nx\,dx=\sin^2nx\,\dfrac{d}{dx}\left\{\dfrac{\phi(\sin x)}{\sin nx}\right\},
\]
where the form of the function $\phi$ is defined by the relation
\begin{multline*}
\phi(z)=\dfrac{z^m}{n^2-m^2}-\dfrac{m(m-1)}{(n^2-m^2)\{n^2-(m-2)^2\}}z^{m-2}\\
+\dfrac{m(m-1)(m-2)(m-3)}{(n^2-m^2)\{n^2-(m-2)^2\}\{n^2-(m-4)^2\}}z^{m-4}-\dots,
\end{multline*}
$m$ being a positive integer and $n$ not being of the form $\pm(m-2r)$, where $r$ is a positive integer not greater than $\dfrac m2$. \source{Math. Trip., 1897.}

\item Draw graphs of the transformation formula
\[
(a_2x^2+2b_2x+c_2)y^2=a_1x^2+2b_1x+c_1
\]
corresponding to those of Arts. 301 and 309 for
\[
(a_2x^2+2b_2x+c_2)y=a_1x^2+2b_1x+c_1.
\]

\end{enumerate}

\newpage

\begin{center}
    {\LARGE \bfseries \color{accent} Chapter IX: General Theorems --- Exercises}\\
    \vspace{0.2cm}
    {\color{accent}\rule{0.55\linewidth}{0.8pt}}
\end{center}
\vspace{0.35cm}

\section*{Exercises --- Limiting Forms Expressed as Definite Integrals (Arts. 319--320)}

\begin{enumerate}[leftmargin=2em]

\item Determine by integration the limiting values of the sums of the following series when $n$ is infinitely great:
\begin{enumerate}[label=(\roman*)]
\item $\dfrac{1}{n+1}+\dfrac{1}{n+2}+\dfrac{1}{n+3}+\dots+\dfrac{1}{n+n}$. \source{$a$, 1884.}
\item $\dfrac{n}{n^2+1^2}+\dfrac{n}{n^2+2^2}+\dfrac{n}{n^2+3^2}+\dots+\dfrac{n}{n^2+n^2}$. \source{Oxford, 1888.}
\item $\dfrac{1}{\sqrt{2n-1^2}}+\dfrac{1}{\sqrt{4n-2^2}}+\dfrac{1}{\sqrt{6n-3^2}}+\dots+\dfrac{1}{\sqrt{2n^2-n^2}}$. \source{Clare, etc., 1882.}
\item $\dfrac1n\left\{\sin^{2k}\dfrac{\pi}{2n}+\sin^{2k}\dfrac{2\pi}{2n}+\sin^{2k}\dfrac{3\pi}{2n}+\dots+\sin^{2k}\dfrac{\pi}{2}\right\}$, $k$ being a positive integer. \source{St. John's, 1886.}
\end{enumerate}

\item Show that the limit when $n$ is increased indefinitely of
\[
\dfrac{(n-m)^{\frac12}}{n}+\dfrac{(2^2n-m)^{\frac12}}{2n}+\dfrac{(3^2n-m)^{\frac12}}{3n}+\dots+\dfrac{(n^3-m)^{\frac12}}{n^2} \quad\text{is}\quad \dfrac32.
\]
\source{Colleges, 1892.}

\item Find the limit when $n$ is indefinitely great of the series
\[
\dfrac{\sqrt{n-1}}{n}+\dfrac{\sqrt{2n-1}}{2n}+\dfrac{\sqrt{3n-1}}{3n}+\dots+\dfrac{\sqrt{n^2-1}}{n^2}.
\]
\source{Colleges, 1890.}

\item Evaluate
\[
Lt_{n=\infty}\left[\dfrac{1}{\sqrt{2a^2n-1}}+\dfrac{1}{\sqrt{4a^2n-1}}+\dfrac{1}{\sqrt{6a^2n-1}}+\dots+\dfrac{1}{\sqrt{2a^2n^2-1}}\right].
\]

\item Evaluate
\[
Lt_{n=\infty}\left[\dfrac{n^2}{(n^2+1^2)^{\frac32}}+\dfrac{n^2}{(n^2+2^2)^{\frac32}}+\dots+\dfrac{n^2}{\{n^2+(n-1)^2\}^{\frac32}}\right].
\]
\source{C. S., 1901.}

\end{enumerate}

\section*{Miscellaneous Examples (Arts. 321--353)}

\begin{enumerate}[leftmargin=2em]

\item Integrate
\begin{enumerate}[label=(\roman*)]
\item $\displaystyle\int \dfrac{\sin^2x\,dx}{(x\cos x-\sin x)^2}$; \source{L.}
\item $\displaystyle\int \dfrac{\log x\,dx}{x^2(1-\log x)^2}$. \source{L.}
\end{enumerate}

\item Prove that
\[
\int_{-c}^{c} \dfrac{x(c^2-x^2)\,dx}{(b^2+c^2-2bx)^{\frac52}}=\dfrac45\dfrac{c^5}{b^4} \quad (b>c). \source{L.}
\]

\item If $X=a+2bx+cx^2$, show that $\displaystyle\int \dfrac{dx}{X^n}$ can be made to depend upon $\displaystyle\int \dfrac{dx}{X^{n-1}}$.

Find a reduction formula for $\displaystyle\int \cos mx\sin^nx\,dx$, and apply it to the case $n=4$. \source{L.}

\item Evaluate $\displaystyle\int \dfrac{2x-3}{5x^2-16x+14}\dfrac{dx}{\sqrt{3x^2-10x+9}}$. \source{L.}

\item Prove that $\displaystyle\int_0^b u\dfrac{d^nv}{dx^n}\,dx$ can be made to depend upon $\displaystyle\int_0^b v\dfrac{d^nu}{dx^n}\,dx$.

Hence show that if $f(x)$ be an arbitrary polynomial of degree $n-1$, and
\[
P_n(x)=\dfrac{d^n(Ax^2+Bx+C)^n}{dx^n},
\]
then
\[
\int_\alpha^\beta f(x)P_n(x)\,dx=0,
\]
where $\alpha$, $\beta$ are the roots, considered real, of the quadratic $Ax^2+Bx+C=0$.

\item Prove that the effect of the operation $p\dfrac{d}{dt}+q$ on a periodic function $a\cos(nt+\epsilon)$ is to multiply the amplitude $a$ by $\sqrt{p^2n^2+q^2}$, and to increase the angle $nt+\epsilon$ by $\tan^{-1}\dfrac{pn}{q}$.

Write down the effect of the operation
\[
\left(p\dfrac{d}{dt}+q\right)\Big/\left(P\dfrac{d}{dt}+Q\right),
\]
and generally, of the operation
\[
\left(\alpha+\beta\dfrac{d}{dt}+\gamma\dfrac{d^2}{dt^2}+\dots\right)\Big/\left(A+B\dfrac{d}{dt}+C\dfrac{d^2}{dt^2}+\dots\right)
\]
on the same periodic function. \source{Int. Arts, London.}

\item When $y^2\equiv ax^2+2bx+c$, prove that
\[
\int \dfrac{dx}{y}=\dfrac{1}{\sqrt a}\cosh^{-1}\dfrac{y\sqrt a}{\sqrt{ac-b^2}}, \quad \dfrac{1}{\sqrt a}\sinh^{-1}\dfrac{y\sqrt a}{\sqrt{b^2-ac}} \quad\text{or}\quad \dfrac{1}{\sqrt{-a}}\sin^{-1}\dfrac{y\sqrt{-a}}{\sqrt{b^2-ac}},
\]
the real form to be chosen, and deduce the value of the integral in the degenerate case when $a=0$. \source{Int. Arts, London.}

\item Find the limiting value of $(n!)^{\frac1n}/n$, when $n$ is infinite.

\item Find the limiting value when $n$ is infinite of the $n$th part of the sum of the $n$ quantities
\[
\dfrac{n+1}{n}, \dfrac{n+2}{n}, \dfrac{n+3}{n}, \dots \dfrac{n+n}{n},
\]
and show that it bears to the limiting value of the $n$th root of the product of the same quantities the ratio $3e:8$, where $e$ is the base of the Napierian logarithms. \source{Oxford 1886, and I. P., 1911.}

\item If $na$ is always equal to unity, and $n$ is indefinitely great, show that the limiting value of the product
\[
(1+a^4)\{1+(2a)^4\}^{\frac12}\{1+(3a)^4\}^{\frac13}\{1+(4a)^4\}^{\frac14}\dots\{1+(na)^4\}^{\frac1n}
\]
is $e^{\pi^2/48}$. \source{Oxford, 1888.}

\item Show that the limit of the sum of $n$ terms of the series
\[
\dfrac{n^3}{(n^2+1^2)(n^2+2\cdot1^2)}+\dfrac{n^3}{(n^2+2^2)(n^2+2\cdot2^2)}+\dots+\dfrac{n^3}{(n^2+n^2)(n^2+2n^2)},
\]
when $n$ is infinite, is $\sqrt2\tan^{-1}\sqrt2-\dfrac{\pi}{4}$. \source{$\gamma$, 1901.}

\item Find
\[
Lt_{n=\infty}\left[\dfrac{\sqrt{n-a}}{n-c}+\dfrac{\sqrt{2n-a}}{2n-c}+\dfrac{\sqrt{3n-a}}{3n-c}+\dots+\dfrac{\sqrt{n^2-a}}{n^2-c}\right].
\]

\item Find the limiting value, when $n$ is infinite, of
\[
\left\{\tan\dfrac{\pi}{2n}\cdot\tan\dfrac{2\pi}{2n}\cdot\tan\dfrac{3\pi}{2n}\dots\tan\dfrac{n\pi}{2n}\right\}^{\frac1n}. \source{Oxford I. P., 1903.}
\]

\item Show that the limit of the product
\[
\left(1+\dfrac1n\right)\left(1+\dfrac2n\right)^{\frac12}\left(1+\dfrac3n\right)^{\frac13}\dots\left(1+\dfrac nn\right)^{\frac1n},
\]
when $n$ is increased indefinitely, is $e^{\pi^2/12}$. \source{Colleges, 1896.}

\item Find the limit, when $n$ is indefinitely increased, of
\[
\dfrac1n\left\{\sec\dfrac{x}{n}+\sec\dfrac{2x}{n}+\dots+\sec\dfrac{(n-1)x}{n}\right\},
\]
where $x$ is $<\dfrac{\pi}{2}$.

Examine the case when $x>\dfrac{\pi}{2}$.

\item Find the limiting value of
\[
2\log n-\log\left[(1+n^2)^{\frac1n}(2^2+n^2)^{\frac1n}\dots(2n^2)^{\frac1n}\right],
\]
when $n$ is indefinitely increased. \source{Oxford I. P., 1900.}

\item Show from elementary considerations that when $n$ increases indefinitely,
\[
1+\dfrac12+\dfrac13+\dots+\dfrac1n-\log n
\]
approaches a finite limit intermediate between $\dfrac12$ and $1$. \source{St. John's, 1884.}

\item If $f(x)=f(a+x)$, show that
\[
\int_a^{na} f(x)\,dx=(n-1)\int_0^a f(x)\,dx,
\]
and illustrate geometrically. \source{Oxford I. P., 1888.}

\item Prove that $\displaystyle\int_0^a \phi(x)\,dx=\int_0^a \phi(a-x)\,dx$,

and show that
\begin{enumerate}[label=(\arabic*)]
\item $\displaystyle\int_0^\pi \dfrac{x\sin^nx}{1+\cos^2x}\,dx=\dfrac{\pi}{2}\int_0^\pi \dfrac{\sin^nx}{1+\cos^2x}\,dx$,
\end{enumerate}
and evaluate this integral when $n=1$ and when $n=3$.
\begin{enumerate}[label=(\arabic*)]
\setcounter{enumii}{1}
\item $\displaystyle\int_0^{\pi/2} \log\dfrac{(1+\sin^2x)^2}{1+\frac14\sin^22x}\,dx=\dfrac{\pi}{2}\log2$.
\end{enumerate}

\item If $\phi(x)=-\phi(2a-x)$, show that
\[
\int_b^{2a}\phi(x)\,dx=-\int_0^b\phi(x)\,dx. \source{Colleges, 1886.}
\]

\item Prove that
\[
\int_b^c \dfrac{\phi(x-b)}{\phi(c-x)}\,dx=\int_b^c \dfrac{\phi(c-x)}{\phi(x-b)}\,dx,
\]
provided $\phi(x)$ remains finite when $x$ vanishes. \source{St. John's, 1883.}

\item Prove that if $\phi(x)$, $\psi(x)$, $\phi'(x)$, $\psi'(x)$ be continuous and finite from $x=a$ to $x=b$,
\[
\int_a^b \phi'(x)\psi'(x)\,dx=\phi\{a+\theta(b-a)\}[\psi(b)-\psi(a)],
\]
where $\theta$ is a positive proper fraction.

\item Prove that $\displaystyle\int_a^{\pi-a} xf(\sin x)\,dx=\dfrac{\pi}{2}\int_a^{\pi-a} f(\sin x)\,dx$. \source{St. John's, 1883.}

\item Show that
\[
\int_a^b f^n(x)\phi(c-x)\,dx-\int_a^b f(x)\phi^n(c-x)\,dx=\sum_{r=1}^{n}\left[f^{r-1}(x)\phi^{n-r}(c-x)\right]_a^b,
\]
where $f^n(x)$ means the $n$th differential coefficient of $f(x)$. \source{$\gamma$, 1893.}

\item Show that, if $\displaystyle\psi(x)=\int_0^x \phi(x)\phi'(2a-x)\,dx$,

then $\psi(2a)-2\psi(a)=[\phi(a)]^2-\phi(0)\phi(2a)$. \source{Trinity, 1895.}

\item If $f(x,y)$ is symmetrical in $x$ and $y$, prove that
\[
\int_{1-b}^{b} xf(x,1-x)\,dx=\dfrac12\int_{1-b}^{b} f(x,1-x)\,dx. \source{Colleges $a$, 1889.}
\]

\item Examine under what limitations the formula
\[
\int_b^a \phi(x)\,dx=\int_b^c \phi(x)\,dx+\int_c^a \phi(x)\,dx
\]
holds good.

Show that
\[
\int_{-\infty}^{\infty}\left(x+\dfrac1x\right)\phi\left(x-\dfrac1x\right)\dfrac{dx}{x}=2\int_{-\infty}^{\infty}\phi(x)\,dx. \source{Math. Tripos, 1884.}
\]

\item If
\[
A_n=1+\dfrac13+\dfrac15+\dots+\dfrac{1}{2n+1}, \qquad B_m=\dfrac12+\dfrac14+\dots+\dfrac{1}{2m},
\]
show that when $n$ and $m$ are both infinite and the ratio $n:m$ tends to a limit $k^2$,
\[
A_n-B_m=\log2+\log k. \source{Colleges $a$, 1888.}
\]

\item Show that
\[
\dfrac{1}{D^n}e^{ax}\dfrac{\sin}{\cos}(bx+\alpha)=e^{ax}\dfrac{\dfrac{\sin}{\cos}\left(bx+\alpha-n\tan^{-1}\dfrac{b}{a}\right)}{(a^2+b^2)^{\frac{n}{2}}}
+A_0+A_1x+A_2x^2+\dots+A_{n-1}x^{n-1},
\]
$A_0$, $A_1$, etc., being arbitrary constants, and also that it may be written
\[
e^{ax}\dfrac{1}{(D+a)^n}\dfrac{\sin}{\cos}(bx+\alpha)+A_0+A_1x+\dots+A_{n-1}x^{n-1},
\]
and explain how the latter operation is to be conducted.

\item If $\displaystyle I_1=\int_0^{\pi/2}\log(1+a_1\sin^2\theta)\,d\theta$,

show that $I_1=\dfrac{\pi}{4}\log(1+a_1)+\dfrac12I_2$,

where $\displaystyle I_2=\int_0^{\pi/2}\log(1+a_2\sin^2\theta)\,d\theta$

and $4(1+a_2)(1+a_1)=(2+a_1)^2$.

Hence show that
\[
I_1=\dfrac{\pi}{4}\log\left[(1+a_1)(1+a_2)^{\frac12}(1+a_3)^{\frac14}\dots\right],
\]
where $4(1+a_{r+1})(1+a_r)=(2+a_r)^2$.

\item Show that if $n>1$,
\[
\int_0^1 \tanh\dfrac{1}{nx}\,dx<\dfrac1n(1+\log n). \source{Oxford I. P., 1911.}
\]

\item How is the equation
\[
\int_a^b f'(x)\,dx=f(b)-f(a)
\]
to be interpreted when $f(x)$ is not a single-valued function?

Illustrate your answer by evaluating
\[
\int_0^{\frac12n\pi} \dfrac{d\theta}{a^2\cos^2\theta+b^2\sin^2\theta},
\]
where $a$ and $b$ are real and $n$ is a positive integer. \source{Oxford I. P., 1912.}

\item Remembering that $\displaystyle\int_0^\infty$ means the limit tended to by $\displaystyle\int_\epsilon^\eta$ as the first of the two positive quantities $\epsilon$, $\eta$ tends to zero, and the second to infinity, prove that if $a>1$, the value of
\[
\int_0^\infty (a^ne^{-ax}-e^{-x})x^{n-1}\,dx
\]
is zero if $n>0$, but not if $n=0$. \source{Oxford I. P., 1917.}

\item If $f(x)$ be any function of $x$ which can be put into partial fractions of the form $\dfrac{A}{a^2-x^2}$, then will
\[
\int_0^\infty \dfrac{f(x)}{1+x^2}\,dx=\dfrac{\pi}{2}f(\sqrt{-1}). \source{R. P.}
\]

\item If $0<b<a$, $a_1=\dfrac12(a+b)$, $b_1=(ab)^{\frac12}$,

prove that $b<b_1<a_1<a$, $a_1-b_1<\dfrac12(a-b)$.

Show that if $(a+b)\tan\theta\cot\phi=a-b\tan^2\theta$,

then
\[
\int_0^{\pi/2}(a^2\cos^2\theta+b^2\sin^2\theta)^{-\frac12}\,d\theta=\int_0^{\pi/2}(a_1^2\cos^2\phi+b_1^2\sin^2\phi)^{-\frac12}\,d\phi. \source{Math. Trip., Part II., 1915.}
\]

\item Show that
\[
\int_0^{\pi/2} \sin\theta\tan^{-1}(\sin\theta)\,d\theta=\dfrac{\pi}{2}(\sqrt2-1). \source{Math. Trip., 1882.}
\]

\item Show how to evaluate $\displaystyle\int R(x,y)\,dx$, where $R(x,y)$ denotes any rational algebraic function of the coordinates $x$, $y$ of a point on a conic. \source{St. John's, 1891.}

\item Show that if $a$ be greater than unity,
\[
\int_0^\pi \dfrac{x\,dx}{a^2-\cos^2x}=\dfrac{\pi^2}{2a\sqrt{a^2-1}}. \source{Oxf. I. P., 1890.}
\]

\item Prove that
\[
\int_0^\infty \phi(x)\,dx=\int_0^1 \left\{\phi(x)+\dfrac{1}{x^2}\phi\left(\dfrac1x\right)\right\}dx. \source{St. John's Coll., 1882.}
\]

\item Prove that $\displaystyle\int_0^\pi \dfrac{x\sin x\,dx}{2+\cos2x}=\dfrac{\pi}{\sqrt2}\tan^{-1}\sqrt2$. \source{Oxf. I. P., 1889.}

\item Integrate $\displaystyle\int_a^b \dfrac{dx}{c^2-x^2}$, when $c$ lies between $a$ and $b$. \source{R. P.}

\item Prove that
\[
\int_0^1 x^n(2-x)^n\,dx=2^{2n}\int_0^1 x^n(1-x)^n\,dx. \source{Oxf. II. P., 1886.}
\]

\item Prove that
\[
\int_a^{\pi/2-a} 2\cos^2x\,\phi(\sin2x)\,dx=\int_{2a}^{\pi/2}\phi(\sin x)\,dx. \source{St. John's.}
\]

\item Show that
\[
\theta=\int_{a-c}^r \dfrac{(a^2-c^2+r^2)}{r\sqrt{2(a^2+c^2)r^2-(a^2-c^2)^2-r^4}}\,dr
\]
is the equation of a circle. \source{Math. Trip., 1882.}

\item Find the integrals
\begin{enumerate}[label=(\alph*)]
\item $\displaystyle\int \left\{\left(\dfrac{x}{e}\right)^x+\left(\dfrac{e}{x}\right)^x\right\}\log x\,dx$;
\item $\displaystyle\int \dfrac{dx}{x}\sqrt{x^2(x-3)^2-4x}$;
\item $\displaystyle\int \dfrac{dx}{\sin^3x\tan^n\dfrac{x}{2}}$. \source{St. John's, 1887.}
\end{enumerate}

\item Prove that
\[
\int_0^\alpha \log(1+\tan\alpha\tan x)\,dx=\alpha\log\sec\alpha. \source{Colls., 1896.}
\]

\item Evaluate
\begin{enumerate}[label=(\roman*)]
\item $\displaystyle\int \dfrac{xe^x}{(x+1)^2}\,dx$. \source{Trin., 1891.}
\item $\displaystyle\int e^x\dfrac{x^2+1}{(x+1)^2}\,dx$. \source{Hall, I.C.}
\item $\displaystyle\int \dfrac{e^{x\sqrt2}}{1-x\sqrt2}\cdot\dfrac{1-x^2}{\sqrt{1-2x^2}}\,dx$.
\item $\displaystyle\int \dfrac{xe^x\,dx}{(e^x-1)^3}$. \source{Hall, I.C.}
\item $\displaystyle\int \dfrac{dx}{(1+x\tan x)^2}$. \source{Trin., 1891.}
\item $\displaystyle\int_e^{e^2} \dfrac{\sec x\csc x}{\log\tan x}\,dx$. \source{Trin., 1884.}
\item $\displaystyle\int \dfrac{\log(1+x^2)}{\sqrt{1-x}}\,dx$.
\end{enumerate}

\item If $\displaystyle I_n\equiv\int_{-1}^1 (1-x^2)^n\cos ax\,dx$,

show that $a^2I_n=2n(2n-1)I_{n-1}-4n(n-1)I_{n-2}$, provided $n>1$.

Show also that
\[
I_n=\dfrac{n!}{a^{2n+1}}\{f(a)\sin a+g(a)\cos a\},
\]
where $f(a)$ and $g(a)$ are algebraic functions of $a$, of degrees $\leqslant n$, with integral coefficients. \source{Trin., 1892.}

\item Show that
\begin{enumerate}[label=(\roman*)]
\item $\displaystyle\int x^3\sqrt{\dfrac{1+x^2}{1-x^2}}\,dx=\dfrac12\tan^{-1}\sqrt{\dfrac{1+x^2}{1-x^2}}-\dfrac14(2+x^2)\sqrt{1-x^4}$.
\item $\displaystyle\int \dfrac{x^2+1}{x^2-1}\dfrac{dx}{\sqrt{1-ax^2+x^4}}=\dfrac{1}{\sqrt{a-2}}\cos^{-1}\dfrac{x\sqrt{a-2}}{x^2-1}$ \quad $(a>2)$.
\end{enumerate}
\source{Hall, I.C., p. 325 and p. 346.}

\item Show that
\[
\int_0^1 x^{cx^a}\,dx=1-\dfrac{c}{(a+1)^2}+\dfrac{c^2}{(2a+1)^3}-\dfrac{c^3}{(3a+1)^4}+\dots. \source{Anglin.}
\]

\item Prove that
\begin{enumerate}[label=(\roman*)]
\item $\displaystyle\dfrac{2}{\pi}\int_0^{\pi/2} \dfrac{d\theta}{\sqrt{1+\sin^2\theta}}=1^2-\left(\dfrac12\right)^2+\left(\dfrac{1\cdot3}{2\cdot4}\right)^2-\dots$.
\item $\displaystyle1+\dfrac{1}{2^2}\left(\dfrac12\right)^2+\dfrac{1}{3^2}\left(\dfrac{1\cdot3}{2\cdot4}\right)^2+\dots=\dfrac{11}{\pi}-4$.
\end{enumerate}
\source{Anglin.}

\item If $\phi(x)=a_1x+\dfrac23a_3x^3+\dfrac{2\cdot4}{3\cdot5}a_5x^5+\dots$,

prove that
\begin{enumerate}[label=(\roman*)]
\item $\displaystyle\int_0^{\pi/2} \cos^2\theta\,\phi(\sin\theta)\,d\theta=\dfrac13a_1+\dfrac15\left(\dfrac23\right)^2a_3+\dfrac17\left(\dfrac{2\cdot4}{3\cdot5}\right)^2a_5+\dots$.
\item $\displaystyle\dfrac{\pi}{2}=1+\dfrac13\cdot1^2+\dfrac15\left(\dfrac23\right)^2+\dfrac17\left(\dfrac{2\cdot4}{3\cdot5}\right)^2+\dots$.
\item $\displaystyle\pi-3=\dfrac{1}{3^2}\cdot1^2+\dfrac{1}{5^2}\left(\dfrac23\right)^2+\dfrac{1}{7^2}\left(\dfrac{2\cdot4}{3\cdot5}\right)^2+\dots$.
\end{enumerate}
\source{Anglin.}

\end{enumerate}

\newpage

\begin{center}
    {\LARGE \bfseries \color{accent} Chapter X: Differentiation, etc., under an Integration Sign --- Exercises}\\
    \vspace{0.2cm}
    {\color{accent}\rule{0.55\linewidth}{0.8pt}}
\end{center}
\vspace{0.35cm}

\section*{General Examples (Arts. 354--366)}

\begin{enumerate}[leftmargin=2em]

\item Prove that
\[
\frac{d}{da}\int_{a^q}^{a^p} a^m x^n\,dx=a^{m-1}\left[\left(\frac{m}{n+1}+p\right)a^{(n+1)p}-\left(\frac{m}{n+1}+q\right)a^{(n+1)q}\right],
\]
and verify the result by performing the integration first.

\item If $A$ be the area bounded by a parabola and its latus rectum ($4a$), prove

(1) by differentiating the integral $4\displaystyle\int_0^a \sqrt{ax}\,dx$ with regard to $a$,

(2) by first integrating and then differentiating with regard to $a$,

that
\[
\frac{dA}{da}=\frac{16a}{3}.
\]

\item Apply the method of Art. 355 to prove that
\[
\frac{d}{dc}\int_{\frac{c}{2}}^{\frac{c\sqrt3}{2}}\sqrt{c^2-x^2}\,dx=\frac16\pi c,
\]
and explain geometrically each step of the process.

Obtain the same result by first integrating and then differentiating the result with regard to $c$; and also geometrically.

\item Show that if $u=\displaystyle\int_0^\infty e^{-ax^3-bx^2}\,dx$,

then
\[
3ab\frac{\partial^2u}{\partial b^2}-3a\frac{\partial u}{\partial b}-2b^2\frac{\partial u}{\partial a}=1,
\]
provided $a$ be positive. \source{Trinity, 1888.}

\item Show that $\dfrac{d^n}{dc^n}\displaystyle\int_{-c}^c f(x+c)\,dx=2^nf^{(n-1)}(2c)$. \source{$a$, 1883.}

\item If $f(x+c)=f(x)$ for all values of $x$, show that
\[
\int_0^c f(y+az)\,dy
\]
is independent of $z$. \source{$a$, 1887.}

\item Prove that
\[
\int_0^{\pi/2}\frac{dx}{(a^2\cos^2x+\beta^2\sin^2x)^{n+1}}=\frac{\pi}{2^{n+1}}\frac{(-1)^n}{n!}\left(\frac1a\frac{\partial}{\partial a}+\frac1\beta\frac{\partial}{\partial\beta}\right)^n(a\beta)^{-1}.
\]

\item Prove that if $u=(a^n+b^n)\displaystyle\int_\beta^\alpha F\{(a-b)x\}\,dx$,

where $F$ denotes any function, $\beta$ and $\alpha$ being independent of $a$ and $b$, and $n$ being a positive integer, then
\[
\left\{(a^n+b^n)\left(\frac{\partial}{\partial a}+\frac{\partial}{\partial b}\right)^n-2\cdot n!\right\}u=0.
\]
\source{Oxford, 1886.}

\item If $u=\displaystyle\int_0^c\phi(x,t)\,dt$,

where $c$ is a function of $u$ and $x$, prove that
\[
\frac{du}{dx}=\frac{\displaystyle\int_0^c\dfrac{\partial\phi}{\partial x}\,dt+\dfrac{\partial c}{\partial x}\phi(x,c)}{1-\dfrac{\partial c}{\partial u}\phi(x,c)}.
\]
\source{$\delta$, 1885.}

\item If $u=\displaystyle\int_\alpha^\beta\phi(x,y)\,dy$,

where $\alpha$ and $\beta$ are functions of $x$ and $u$, prove that
\[
\frac{du}{dx}=\frac{\displaystyle\int_\alpha^\beta\dfrac{\partial\phi}{\partial x}\,dy+\phi(x,\beta)\dfrac{\partial\beta}{\partial x}-\phi(x,\alpha)\dfrac{\partial\alpha}{\partial x}}{1-\phi(x,\beta)\dfrac{\partial\beta}{\partial u}+\phi(x,\alpha)\dfrac{\partial\alpha}{\partial u}}.
\]

\item Comment upon the application of the rule of Art. 355 to the case
\[
\frac{d}{da}\int_{-a}^a\frac{\phi(x)\,dx}{\sqrt{a^2-x^2}}.
\]
Prove that in this case the true result is
\[
\frac1a\int_{-a}^a\frac{x\phi'(x)\,dx}{\sqrt{a^2-x^2}}.
\]

\item If
\[
u=\int_{f(a)}^{F(a)}\phi(\theta,a)\,d\theta,
\]
we have
\[
\frac{du}{da}=\int_{f(a)}^{F(a)}\frac{\partial\phi(\theta,a)}{\partial a}\,d\theta+\phi\{F(a),a\}\frac{dF}{da}-\phi\{f(a),a\}\frac{df}{da}.
\]
Do you consider that this formula fails in the case in which
\[
F(a)=a,\quad f(a)=0\quad\text{and}\quad\phi(\theta,a)=\frac{1}{\sqrt{\cos\theta-\cos a}}\ ?
\]
If so, to what extent and in what respect?

Prove that in this case
\[
\frac{du}{da}=\frac{\sin a}{2\sqrt2}\int_0^{\pi/2}\frac{\sin^2\theta\,d\theta}{\left(1-\sin^2\dfrac{a}{2}\sin^2\theta\right)^{3/2}}.
\]
Make any remarks that occur to you as to the reasons for the peculiar form which the general formula assumes in this case. \source{$e$, 1884.}

\item Show that the equation
\[
\frac{d}{dx}\int_0^1 x^3e^{x^2\left(1-\frac1t\right)}\frac{dt}{t^2}=\int_0^1\frac{\partial}{\partial x}\left[x^3e^{x^2\left(1-\frac1t\right)}\frac{1}{t^2}\right]dt
\]
ceases to hold for $x=0$. \source{Math. Tripos, 1897.}

\item Find a curve in which the abscissa of the centroid of the area of that portion bounded by the curve, the coordinate axes and an ordinate is proportional to the abscissa of the bounding ordinate. \source{Colleges, 1878.}

\item A vessel in the form of a right circular cylinder with vertical axis and a flat horizontal base is filled to varying depths with liquid of varying density. If the depth of the centre of gravity of the liquid be always $\dfrac1n$ of the immersed portion of the axis, show that the density varies as $(\text{depth})^{\frac{2-n}{n-1}}$.

\item Find the general equation of all solids of revolution for which the distance from the vertex of the centroid of a segment made by a plane perpendicular to the axis, is proportional to the height of the segment. \source{Todhunter, \textit{Integral Calculus}, p. 198.}

\item Find the form of the curve for which the area bounded by the curve, the coordinate axes and an ordinate is such that the moments of inertia of this area about the coordinate axes are in a constant ratio.

\item A body moves from rest at a distance $a$ towards a centre of attraction varying inversely as the distance. Show that the time of describing the space between $\beta a$ and $\beta^na$ will be a maximum when
\[
\beta n^{\frac{1}{2(n-1)}}=1
\]
\[
\left[\text{It may be assumed that }\left(\frac{dx}{dt}\right)^2\propto\log\frac{a}{x}.\right]
\]
\source{Tait and Steele, \textit{Dynamics of a Particle}.}

\item Find the density of a parabolic plate as a function of the abscissa in order that the distance of the centroid from the vertex may vary as the square root of the length of the plate. \source{$a$, 1881.}

\item Find the equation of the curve such that the area included by the ordinate at any point, the axis of $x$ and the curve is in a constant ratio to the area included by the ordinate, the axis of $x$ and the tangent. \source{Math. Tripos, 1882.}

\item Prove that
\[
\frac{d}{da}\int_0^a\frac{F(x)\,dx}{\sqrt{a-x}}=\int_0^a\frac{F(x)+2x\dfrac{dF(x)}{dx}}{2a\sqrt{a-x}}\,dx.
\]
Under what circumstances will $\displaystyle\int_0^a\frac{F(x)\,dx}{\sqrt{a-x}}$ be independent of $a$? \source{Todhunter, \textit{Int. Calc.}}

\item If $\tan\dfrac{\pi}{8}\sin\phi=x\sqrt{\dfrac{1-x^2}{1+x^2}}=\sin\psi$,

verify that
\[
\int_0^x\frac{dx}{\sqrt{1-x^8}}=\frac{1}{2\sqrt2}\int_0^\phi\frac{d\phi}{\sqrt{1-\tan^2\dfrac{\pi}{8}\sin^2\phi}}+\sin^2\frac{\pi}{8}\int_0^\psi\frac{d\psi}{\sqrt{1-\tan^2\dfrac{\pi}{8}\sin^2\psi}}.
\]
\source{Math. Tripos, 1896.}

\item Prove that
\[
b^2\int_0^\infty\frac{\cos bx}{a^2+x^2}\,dx+6\int_0^\infty\frac{\cos bx}{(a^2+x^2)^2}\,dx-8a^2\int_0^\infty\frac{\cos bx}{(a^2+x^2)^3}\,dx=0.
\]
\source{$\epsilon$, 1884.}

\item Verify that
\[
y=A\int_0^1v^{\beta-1}(1-v)^{\gamma-\beta-1}(1-xv)^{-\alpha}\,dv
\]
satisfies the differential equation of the hypergeometric series, viz.
\[
x(1-x)\frac{d^2y}{dx^2}+\{\gamma-(\alpha+\beta+1)x\}\frac{dy}{dx}-\alpha\beta y=0,
\]
when $\beta>0$ and $\gamma>\beta$.

\item If $u=\displaystyle\int_0^\pi e^{qx\cos\theta}\{A+B\log(x\sin^2\theta)\}\,d\theta$,

verify that
\[
x\frac{d^2u}{dx^2}+\frac{du}{dx}-q^2xu=0.
\]

\item Prove that $y=\dfrac1\pi\displaystyle\int_0^\pi(x+\cos\phi\sqrt{x^2-1})^n\,d\phi$

satisfies the equation
\[
\frac{d}{dx}\left[(x^2-1)\frac{dy}{dx}\right]=n(n+1)y.
\]
\source{St. John's, 1892.}

\item Show that the differential equation
\[
\frac{d^2u}{dx^2}+u=\frac{a}{x}+\frac{b}{x^2}
\]
is satisfied by
\[
u=\int_0^\infty\frac{a\sin z+b\cos z}{x+z}\,dz.
\]
Write down the complete solution. \source{St. John's, 1883.}

\item If $y=\displaystyle\int_0^\pi\sqrt{x}e^{nx\cos\theta}\cos\{\sqrt7\log(\sqrt x\sin\theta)+\alpha\}\,d\theta$,

prove that
\[
\frac{d^2y}{dx^2}=\left(n^2-\frac{2}{x^2}\right)y.
\]
\source{St. John's, 1889.}

\item Prove that
\[
\int_0^\pi(\cosh x-\sinh x\cos\phi)^{\frac{2n-1}{2}}\,d\phi=\int_0^\pi\frac{d\phi}{(\cosh x-\sinh x\cos\phi)^{\frac{2n+1}{2}}}.
\]
Prove also that if
\[
P=\int_0^\pi(\cosh x-\sinh x\cos\phi)^{\frac12}\,d\phi,
\]
\[
\frac{dP}{dx}=\frac12\int_0^\pi\frac{\cos\phi\,d\phi}{(\cosh x-\sinh x\cos\phi)^{\frac32}}.
\]
\source{$a$, 1886.}

\item If $\dfrac{y}{x}=\displaystyle\int_0^{\pi/2}\cos(mx^n\sin\phi)\cos^{\frac1n}\phi\,d\phi$, prove that $y$ satisfies the equation
\[
\frac{d^2y}{dx^2}+m^2n^2x^{2n-2}y=0.
\]
\source{$a$, 1886.}

\item Verify that
\[
y=\int e^{ux}V[A+B\log\{U_1(a_2+b_2x)\}]\,du,
\]
where $U_1=b_2u^2+b_1u+b_0$, $\log VU_1=\displaystyle\int\frac{a_2u^2+a_1u+a_0}{U_1}\,du$,

and the limits are given by $e^{ux}VU_1=0$, satisfies the equation
\[
(a_2+b_2x)\frac{d^2y}{dx^2}+(a_1+b_1x)\frac{dy}{dx}+(a_0+b_0x)y=0,
\]
provided $a_1b_2-a_2b_1=b_2^2$. \source{Spitzer, \textit{Crelle}, vol.\ liv.}

\item Verify that if $x$ be positive
\[
u=C_1\int_{-q}^q e^{xt}(t^2-q^2)^{\frac{a}{2}-1}\,dt+C_2\int_{-\infty}^{-q}e^{xt}(t^2-q^2)^{\frac{a}{2}-1}\,dt,
\]
and if $x$ be negative
\[
u=C_1\int_{-q}^q e^{xt}(t^2-q^2)^{\frac{a}{2}-1}\,dt+C_2\int_q^\infty e^{xt}(t^2-q^2)^{\frac{a}{2}-1}\,dt
\]
solves the differential equation
\[
x\frac{d^2u}{dx^2}+a\frac{du}{dx}-q^2xu=0.
\]
\source{Petzval.}

\item Prove that
\[
\int_{\frac{\pi}{2}-a}^{\frac{\pi}{2}}\sin\theta\arccos\frac{\cos a}{\sin\theta}\,d\theta=\frac{\pi}{2}(1-\cos a).
\]
\source{Trinity, 1886.}

\item Prove that
\[
\int_0^1\log\frac{1+ax}{1-ax}\frac{dx}{x\sqrt{1-x^2}}=\pi\sin^{-1}a.
\]
\source{Oxford, 1888.}

\item Establish the known result
\[
\int_0^\infty\frac{\log x}{(a+x)^2}\,dx=\frac{\log a}{a},
\]
and hence prove that when $n$ is a positive integer

(a) \[
\int_0^\infty\frac{\log x}{(a+x)^{n+2}}\,dx=\frac{1}{(n+1)a^{n+1}}\left\{\log a-\frac11-\frac12-\frac13-\dots-\frac1n\right\},
\]

($\beta$) \[
\int_0^\infty\frac{(\log x)^2}{(a+x)^{n+2}}\,dx=\frac{1}{(n+1)a^{n+1}}\left\{\left(\frac{\pi^2}{3}-\frac{1}{1^2}-\frac{1}{2^2}-\frac{1}{3^2}-\dots-\frac{1}{n^2}\right)+\left(\log a-\frac11-\frac12-\frac13-\dots-\frac1n\right)^2\right\}.
\]
\source{Math. Tripos, 1883.}

\item If the operator $\Delta$, applied to a function of $a$, has the effect of changing $a$ to $a+1$, and subtracting the original function, show that
\[
\Delta\int_a^b\phi(x,a)\,dx=\int_a^b\Delta\phi(x,a)\,dx,
\]
where $a$ and $b$ are independent of $a$.

Prove that
\[
\int_0^\infty e^{-ax}(e^{-x}-1)^n\,dx=\frac{(-1)^n}{a(a+1)\dots(a+n)}.
\]
\source{Bertrand, \textit{C.I.}, p. 183.}

\item Given $u=\displaystyle\int_0^\infty\frac{\cos ax}{1+x^2}\,dx$, differentiating twice we have
\[
\frac{d^2u}{da^2}=-\int_0^\infty\frac{x^2\cos ax}{1+x^2}\,dx.
\]
But this is indeterminate when $x$ is infinite. Discuss the validity of the differentiation. \source{Bertrand, \textit{Cal. Int.}, p. 181.}

\item Is it true that
\[
\int_0^1\left[\int_0^1\frac{a^2-x^2}{(a^2+x^2)^2}\,dx\right]da=\int_0^1\left[\int_0^1\frac{a^2-x^2}{(a^2+x^2)^2}\,da\right]dx\ ?
\]
If not, why not? [See Art. 1899, Vol. II.]

Evaluate each side separately and compare the results. \source{Bertrand, \textit{Cal. Int.}, p. 187.}

\item If $P+\iota Q=\phi(x+\iota y)$, show that in general
\[
\int_a^b\int_\alpha^\beta\frac{\partial Q}{\partial y}\,dx\,dy=\int_\alpha^\beta\int_a^b\frac{\partial P}{\partial x}\,dy\,dx
\]
and
\[
\int_a^b\int_\alpha^\beta\frac{\partial P}{\partial x}\,dx\,dy=-\int_\alpha^\beta\int_a^b\frac{\partial Q}{\partial y}\,dy\,dx.
\]
Examine the case $\phi(x+\iota y)=e^{-(x+\iota y)^2}$, taking $a=0$, $\alpha=0$ and $b=\infty$. \source{Bertrand.}

\item If $f(x)=(2-x)^{\frac12}\displaystyle\int_0^{\pi/2}(1-x\sin^2\theta)^{-\frac12}\,d\theta$, show that
\[
\frac{df(x)}{dx}=\frac12(2-x)^{-\frac12}\int_0^{\pi/2}\cos2\theta\left[(1-x\cos^2\theta)^{-\frac32}-(1-x\sin^2\theta)^{-\frac32}\right]d\theta.
\]
Hence show that as $x$ increases from $0$ to $1$, $f(x)$ increases from $\dfrac{\pi}{\sqrt2}$ to $\infty$. \source{C. S., 1898.}

\item Prove that
\[
\int_0^u du\int_0^u du\int_0^u du\dots\int_0^u du\,f(u)=\frac{1}{(n-1)!}\int_0^u(u-z)^{n-1}f(z)\,dz,
\]
there being $n$ integration signs in the left member of the equality. \source{R. P.}

\item Show that
\[
\frac{d^2}{dc^2}\left\{\int_0^c\int_0^c\phi(x+y+c)\,dx\,dy\right\}=9\phi(3c)-8\phi(2c)+\phi(c).
\]
\source{Oxf. II. P., 1890.}

\item Show that the quartic function
\[
Q\equiv ax^4+4bx^3+6cx^2+4dx+e
\]
can, in general, be expressed in three different ways as the sum of two squares $P^2+R^2$, where
\[
P\equiv a^{-\frac32}\left[(ax+b)^2+3(ac-b^2)-2\lambda\right]
\]
and
\[
R\equiv a^{-\frac52}\lambda^{-\frac12}\left[2(ax+b)\lambda+a^2d-3abc+2b^3\right],
\]
$\lambda$ having any one of three determinate values $\lambda_1$, $\lambda_2$, $\lambda_3$.

Verify the evaluation of the integral $\displaystyle\int\frac{dx}{Q}$ in the form
\[
\frac{a^2}{4}\left[(\lambda_2-\lambda_3)\lambda_1^{\frac12}\tan^{-1}\frac{R_1}{P_1}+(\lambda_3-\lambda_1)\lambda_2^{\frac12}\tan^{-1}\frac{R_2}{P_2}+(\lambda_1-\lambda_2)\lambda_3^{\frac12}\tan^{-1}\frac{R_3}{P_3}\right]\Big/\Lambda,
\]
where $\Lambda=(\lambda_2-\lambda_3)(\lambda_3-\lambda_1)(\lambda_1-\lambda_2)$. \source{Math. Trip., 1897.}

\item Show that
\[
\int_0^ax^{n-1}(1-a+x)^{-n-r}\,dx=\frac{(n-1)!}{(n+r-1)!}\left(\frac{d}{da}\right)^{r-1}\frac{a^{n+r-1}}{1-a}
\]
\source{I. C. S., 1892.}

\end{enumerate}

\newpage

\begin{center}
    {\LARGE \bfseries \color{accent} Chapter XI: Preliminary to Integration of $\displaystyle\int\frac{M}{N}\frac{dx}{\sqrt{Q}}$. Definitions of Elliptic Functions --- Exercises}\\
    \vspace{0.2cm}
    {\color{accent}\rule{0.55\linewidth}{0.8pt}}
\end{center}
\vspace{0.35cm}

\providecommand{\sn}{\operatorname{sn}}
\providecommand{\cn}{\operatorname{cn}}
\providecommand{\dn}{\operatorname{dn}}
\providecommand{\tnn}{\operatorname{tn}}
\providecommand{\amp}{\operatorname{am}}

\section*{Exercises --- Legendre's Formulae (Art. 391)}

Prove the following twelve results:

\begin{enumerate}[leftmargin=2em]

\item $\displaystyle\int_0^\theta\frac{d\theta}{\Delta^3}=\frac{1}{k'^2}E(\theta,k)-\frac{k^2}{k'^2}\frac{\sin\theta\cos\theta}{\Delta}$.

[Putting $P=\dfrac{\sin\theta\cos\theta}{\Delta}$ and differentiating, we obtain, after a little reduction, $k^2\dfrac{dP}{d\theta}=\Delta-\dfrac{k'^2}{\Delta^3}$, then integrating we obtain the result stated.]

\item $\displaystyle\int_0^\theta\frac{\sin^2\theta\,d\theta}{\Delta}=\frac{1}{k^2}(F-E)$.

\item $\displaystyle\int_0^\theta\frac{\cos^2\theta\,d\theta}{\Delta}=\frac{1}{k^2}(E-k'^2F)$.

\item $\displaystyle\int_0^\theta\frac{\sec^2\theta\,d\theta}{\Delta}=\frac{1}{k'^2}(\Delta\tan\theta+k'^2F-E)$.

\item $\displaystyle\int_0^\theta\frac{\tan^2\theta\,d\theta}{\Delta}=\frac{1}{k'^2}(\Delta\tan\theta-E)$.

\item $\displaystyle\int_0^\theta\frac{\tan^2\dfrac{\theta}{2}\,d\theta}{\Delta}=2\Delta\tan\frac{\theta}{2}+F-2E$.

\item $\displaystyle\int_0^\theta\frac{\sec^2\dfrac{\theta}{2}\,d\theta}{\Delta}=2\Delta\tan\frac{\theta}{2}+2F-2E$.

\item $\displaystyle\int_0^\theta\Delta\sec^2\theta\,d\theta=\Delta\tan\theta+F-E$.

\item $\displaystyle\int_0^\theta\Delta\tan^2\theta\,d\theta=\Delta\tan\theta+F-2E$.

\item $\displaystyle\int_0^\theta\Delta^3\,d\theta=\frac{k^2}{3}\Delta\sin\theta\cos\theta+\frac{4-2k^2}{3}E-\frac{k'^2}{3}F$.

\item $\displaystyle\int_0^\theta\Delta\sin^2\theta\,d\theta=-\frac13\Delta\sin\theta\cos\theta+\frac{2k^2-1}{3k^2}E+\frac{k'^2}{3k^2}F$.

\item $\displaystyle\int_0^\theta\Delta\cos^2\theta\,d\theta=\frac13\Delta\sin\theta\cos\theta+\frac{1+k^2}{3k^2}E-\frac{k'^2}{3k^2}F$.

\end{enumerate}

\section*{General Examples (Arts. 367--392)}

\begin{enumerate}[leftmargin=2em]

\item By putting $x=\dfrac{1-\sin\theta}{1+\sin\theta}$, shew that
\[
u=\int_x^1\frac{dx}{\sqrt{x(1+6x+x^2)}}=\frac{1}{\sqrt2}\int_0^\theta\frac{d\theta}{\sqrt{1-\frac12\sin^2\theta}}=\frac{1}{\sqrt2}F\left(\theta,\frac{1}{\sqrt2}\right);
\]
and that
\[
x=\frac{1-\sn(u\sqrt2)}{1+\sn(u\sqrt2)},\ \left(\text{mod. }\frac{1}{\sqrt2}\right);\quad\text{i.e. } u=\frac{1}{\sqrt2}\sn^{-1}\left(\frac{1-x}{1+x},\frac{1}{\sqrt2}\right).
\]

\item Prove that
\[
F_1(\theta,k)=\frac{\pi}{2}\left(1+\frac{1^2}{2^2}k^2+\frac{1^2\cdot3^2}{2^2\cdot4^2}k^4+\frac{1^2\cdot3^2\cdot5^2}{2^2\cdot4^2\cdot6^2}k^6+\dots\right);
\]
and that $F_1\left(\theta,\dfrac1{10}\right)=1.574745$ very nearly.

\item Prove that
\[
E_1(\theta,k)=\frac{\pi}{2}\left(1-\frac{1}{2^2}k^2-\frac{1^2\cdot3}{2^2\cdot4^2}k^4-\frac{1^2\cdot3^2\cdot5}{2^2\cdot4^2\cdot6^2}k^6-\dots\right).
\]

\item Prove that
\[
\Pi_1(\theta,k,n)=\frac{\pi}{2}\left[1+\left(\frac12k^2-n\right)\frac12+\left(\frac{1\cdot3}{2\cdot4}k^4-\frac12k^2n+n^2\right)\frac{1\cdot3}{2\cdot4}+\left(\frac{1\cdot3\cdot5}{2\cdot4\cdot6}k^6-\frac{1\cdot3}{2\cdot4}k^4n+\frac12k^2n^2-n^3\right)\frac{1\cdot3\cdot5}{2\cdot4\cdot6}+\dots\right]\quad\text{if } n \text{ be}<1.
\]

\item Establish the truth of
\begin{enumerate}[label=(\alph*)]
\item $\left(\sn u+\dfrac1{\cn u}\right)^2+\left(\cn u+\dfrac1{\sn u}\right)^2=\left(1+\dfrac1{\sn u\,\cn u}\right)^2$.

\item $\dfrac{\dfrac{\cn u}{\sn u}-\dfrac{\sn u}{\cn u}}{\cn u+\sn u}=\dfrac1{\sn u}-\dfrac1{\cn u}$.

\item $\left(\dfrac1{\sn u}-\sn u\right)\left(\dfrac1{\cn u}-\cn u\right)\left(\dfrac{\sn u}{\cn u}+\dfrac{\cn u}{\sn u}\right)=1$.

\item $\dfrac{1-\cn u}{1+\cn u}=\left(\dfrac1{\sn u}-\dfrac{\cn u}{\sn u}\right)^2$.
\end{enumerate}

\item Prove that
\begin{enumerate}[label=(\arabic*)]
\item $\dn^2u-k^2\cn^2u=k'^2$,
\item $\dfrac1{\cn^2u}=1+\tnn^2u$,
\item $\dfrac1{\sn^2u}=1+\dfrac1{\tnn^2u}$.
\end{enumerate}

\item Prove that
\begin{enumerate}[label=(\arabic*)]
\item $\dfrac{d}{du}\sn^2u=2\sn u\,\cn u\,\dn u$,
\item $\displaystyle\int \sn^pu\,\cn u\,\dn u\,du=\dfrac{\sn^{p+1}u}{p+1}$,
\item $\displaystyle\int \dfrac{\dn u}{a+b\cn u}\,du=\dfrac1{\sqrt{a^2-b^2}}\cos^{-1}\left(\dfrac{a\cn u+b}{a+b\cn u}\right)$, $\quad a>b$.
\end{enumerate}

\item Prove
\[
2\sn u\,\cn v=\sin(\amp u+\amp v)+\sin(\amp u-\amp v),
\]
\[
2\cn u\,\cn v=\cos(\amp u+\amp v)+\cos(\amp u-\amp v).
\]

\item By putting $x=a\cos\theta$, show that
\[
\int_x^a\frac{dx}{\sqrt{a^4-x^4}}=\frac{1}{a\sqrt2}\cn^{-1}\left(\frac{x}{a},\frac{1}{\sqrt2}\right).
\]

\item Prove $\displaystyle\int_a^x\frac{dx}{\sqrt{x^4-a^4}}=\frac{1}{a\sqrt2}\cn^{-1}\left(\frac{a}{x},\frac{1}{\sqrt2}\right)$.

\item By putting $x=a\sqrt{\dfrac{1+z}{1-z}}$, show that
\[
\int_x^\infty\frac{dx}{\sqrt{a^4+x^4}}=\frac{1}{2a}\cn^{-1}\left(\frac{x^2-a^2}{x^2+a^2},\frac{1}{\sqrt2}\right).
\]

\item Prove that
\[
\int_x^\infty\frac{dx}{\sqrt{a^4+2a^2x^2\cos2\alpha+x^4}}=\frac{1}{2a}\cn^{-1}\left(\frac{x^2-a^2}{x^2+a^2},\sin\alpha\right).
\]

\item Prove that
\[
\sn K=1,\quad \cn K=0,\quad \dn K=k',\quad \tnn K=\infty.
\]

\item Prove that
\begin{enumerate}[label=(\arabic*)]
\item $\dfrac{d}{du}(\sn u+\cn u)^n=n(\sn u+\cn u)^{n-1}(\cn u-\sn u)\,\dn u$,
\item $\displaystyle\int (k^2\sn u+\cn u)^n(k^2\cn u-\sn u)\,\dn u\,du=\dfrac{(k^2\sn u+\cn u)^{n+1}}{n+1}$.
\end{enumerate}

\item Draw graphs of $y=\Delta\theta$ and $y=\dfrac1{\Delta\theta}$, showing that the former consists of an undulating curve lying entirely below the line $y=1$ and the other of an undulating line lying entirely above the line $y=1$. Take the cases $k^2=\dfrac12$ and $k^2=\dfrac14$.

Show that the areas bounded by these curves, the $x$-axis, the $y$-axis and any ordinate at a point whose abscissa is $\theta$ represent $E(\theta)$ and $F(\theta)$ completely. Examine what happens in the limiting cases $k=0$ and $k=1$.

\item Show that the complete elliptic integrals of the First and Second Species may be expressed as
\[
F_1=\frac{\pi}{2}f\left(\frac12,\frac12,1,k^2\right),
\]
\[
E_1=\frac{\pi}{2}f\left(-\frac12,\frac12,1,k^2\right),
\]
where $f(a,b,c,x)$ is the hypergeometric series
\[
1+\frac{a\cdot b}{1\cdot c}x+\frac{a\,\overline{a+1}\ b\,\overline{b+1}}{1\cdot2\ c\,\overline{c+1}}x^2+\dots.
\]

\item Show by differentiating $F(\theta,k)$ and $E(\theta,k)$ with regard to $k$
\begin{enumerate}[label=(\arabic*)]
\item $\dfrac{dE}{dk}=\dfrac1k(E-F)$,
\item $\dfrac{dF}{dk}=\dfrac1{kk'^2}(E-k'^2F)-\dfrac{k}{k'^2}\dfrac{\sin\theta\cos\theta}{\Delta}$.
\end{enumerate}

Hence, eliminating $E$ and $F$ alternately, show that
\[
(1-k^2)\frac{d^2F}{dk^2}+\frac{1-3k^2}{k}\frac{dF}{dk}-F+\frac{\sin\theta\cos\theta}{\Delta^3}=0,
\]
\[
(1-k^2)\frac{d^2E}{dk^2}+\frac{1-k^2}{k}\frac{dE}{dk}+E-\frac{\sin\theta\cos\theta}{\Delta}=0,
\]
and for the complete functions $F_1$, $E_1$
\[
(1-k^2)\frac{d^2F_1}{dk^2}+\frac{1-3k^2}{k}\frac{dF_1}{dk}-F_1=0,
\]
\[
(1-k^2)\frac{d^2E_1}{dk^2}+\frac{1-k^2}{k}\frac{dE_1}{dk}+E_1=0.
\]

\end{enumerate}

\newpage

\begin{center}
    {\LARGE \bfseries \color{accent} Chapter XII: Quadrature (I). Plane Surfaces, Cartesian and Polar Equations --- Exercises}\\
    \vspace{0.2cm}
    {\color{accent}\rule{0.55\linewidth}{0.8pt}}
\end{center}
\vspace{0.35cm}

\section*{Exercises --- Cartesian Equations (Arts. 393--406)}

\begin{enumerate}[leftmargin=2em]

\item Obtain the area bounded by a parabola and its latus rectum. A series of ordinates are drawn between the vertex and the latus rectum, parallel to the latter, viz. $x=\left(\dfrac{r}{n}\right)^{\frac23}a$, where $r=1,2,3,\dots n-1$. Show that they divide the aforementioned area into $n$ equal parts.

\item Obtain the areas bounded by the curve, the $x$-axis, and the specified ordinates in the following cases:
\begin{enumerate}[label=(\alph*)]
\item The catenary $y=c\cosh\dfrac{x}{c}$, from $x=0$ to $x=h$.
\item The logarithmic curve $y=e^x$, from $x=0$ to $x=h$.
\item The logarithmic curve $y=\log_ex$, from $x=1$ to $x=h$ $(h>1)$.
\item The ellipse $y=\dfrac{b}{a}\sqrt{a^2-x^2}$, from $x=\sqrt{a^2-b^2}$ to $x=a$.
\item The hyperbola $xy=k^2$, from $x=a$ to $x=b$, $a$ and $b$ both $>0$; first, if the hyperbola be rectangular, second, if the angle between the asymptotes be $\omega$.
\item The curve $y=xe^{x^2}$, from $x=0$ to $x=h$.
\end{enumerate}

\item Obtain the area (1) bounded by the parabolas $y^2=4ax$, $x^2=4ay$;

(2) bounded by the parabolas $y^2=4ax$, $x^2=4by$.

In what ratio is this area divided by the common chord in each case?

\item Find the areas of the portions into which the ellipse $x^2/a^2+y^2/b^2=1$ is divided (1) by the straight line $y=c$;

(2) by the two straight lines $y=c$, $x=d$, supposed to cut within the ellipse.

\item Trace the curve $x^2y^2=a^2(y^2-x^2)$, and find the whole area included between the curve and its asymptotes.

\item Find the area between the curve $y^2(a+x)=(a-x)^3$ and its asymptote.

\item Find the area of the loop of the curve
\[
y^2x+(x+a)^2(x+2a)=0.
\]

\item Two curves in which $y\propto x^m$ and two in which $y\propto x^n$ form a quadrilateral; show that its area is
\[
\frac{m-n}{(m+1)(n+1)}(x_1y_1-x_2y_2+x_3y_3-x_4y_4),
\]
where $(x_1,y_1)$, $(x_2,y_2)$, $(x_3,y_3)$, $(x_4,y_4)$ are the coordinates of the corners taken in order. \source{Trinity, 1891.}

\item By means of the integral $\displaystyle\int y\,dx$ taken round the contour of the triangle formed by the intersecting lines,
\[
y=a_1x+b_1,
\]
\[
y=a_2x+b_2,
\]
\[
y=a_3x+b_3,
\]
show that they enclose the area
\[
\frac{(b_1-b_3)^2}{2(a_1-a_3)}+\frac{(b_2-b_1)^2}{2(a_2-a_1)}+\frac{(b_3-b_2)^2}{2(a_3-a_2)}.
\]
\source{Smith's Prize, 1876.}

\item A four-sided figure is formed by the three parabolas,
\[
y^2-9ax+81a^2=0,
\]
\[
y^2-4ax+16a^2=0,
\]
\[
y^2-ax+a^2=0,
\]
and the axis of $x$. Prove that its area is $12a^2$, and is equal to the area enclosed by the chords of the area. \source{Colleges $\alpha$, 1886.}

\item Find the curvilinear area enclosed between the parabola $y^2=4ax$ and its evolute. \source{Oxf. I. P., 1889.}

\item Show that the area cut off from a semi-cubical parabola by a tangent is divided by the tangent at the cusp in the ratio 64 : 17. \source{Oxford II. P., 1889.}

\item (i) Find the area of a loop of the curve
\[
ay^2=x^2(a-x).
\]
\source{I. C. S., 1882.}

(ii) Find the whole area of the curve
\[
a^2y^2=a^2x^2-x^4.
\]
\source{I. C. S., 1881.}

\item Trace the curve $a^2x^2=y^3(2a-y)$, and prove that its area is equal to that of the circle whose radius is $a$. \source{I. C. S., 1887 and 1890.}

\item Trace the curve $a^4y^2=x^5(2a-x)$, and prove that its area is to that of the circle of radius $a$ as 5 to 4.

\item Find the area of the curve
\[
y(x^2+1)=x^3-1
\]
from $x=0$ to $x=1$. \source{St. John's, 1881.}

\item (i) Find the area between $y^2=\dfrac{x^3}{a-x}$ and its asymptote.

(ii) Show that the whole area between
\[
y^2(x-a)(b-x)=c^2x^2
\]
and its asymptote is $\pi c(a+b)$. \source{Ox. II. P., 1903.}

(iii) Show that the area between the curve
\[
y^2x=a^3-a^2x
\]
and its asymptote is that of a circle of radius $a$. \source{St. John's, 1889.}

\item Find the area between the axis of $x$, the hyperbola $x^2/a^2-y^2/b^2=1$, and the line $y=x\tan\alpha$, where
\[
\frac{b}{a}>\tan\alpha>0.
\]
\source{Ox. I. P., 1901.}

If $A$ be the vertex, $O$ the centre, and $P$ any point on the hyperbola $x^2/a^2-y^2/b^2=1$, prove that
\[
x=a\cosh\frac{2S}{ab},\quad y=b\sinh\frac{2S}{ab},
\]
where $S$ is the sectorial area $AOP$. \source{Math. Tripos, 1885.}

\item Find by integration the area lying on the same side of the axis of $x$ as the positive part of the axis of $y$, and which is contained by the lines
\[
y^2=4ax,
\]
\[
x^2+y^2=2ax,
\]
\[
x=y+2a.
\]
Express the area both when $x$ is the independent variable and when $y$ is the independent variable. \source{Colleges, 1882.}

\item Prove that the area of the loop of
\[
a(x-y)(x-2y)=y^3\quad\text{is}\quad\frac{a^2}{60}.
\]
\source{Coll. $\beta$, 1891.}

\item Find the areas of the two regions of space bounded by the straight line $y=c$, and the curves whose equations are
\[
x^2+y^2=c^2,
\]
\[
y^3+4x^2=4c^3.
\]
\source{I. C. S., 1891.}

\item Prove that the area contained between the curve
\[
(x+3a)(x^2+y^2)=4a^3
\]
and its asymptote is $3a^2\sqrt3$. \source{Oxf. I. P., 1901.}

\item Prove that the area of the curve
\[
x^4-3ax^3+a^2(2x^2+y^2)=0
\]
is $\dfrac38\pi a^2$. \source{Math. Trip., 1893.}

\item Find the area of one loop of the curve
\[
y^4-y^2+x^2=0.
\]
\source{Colleges $\alpha$, 1885.}

\item Through the cusp of the evolute of a parabola, a line is drawn perpendicular to the axis. Show that it divides the area between the parabola and the evolute in the ratio 17 : 5. \source{C. S., 1896.}

\item Show that the ordinate $x=a$ divides the area between $y^2(2a-x)=x^3$ and its asymptote into two parts in the ratio
\[
3\pi-8:3\pi+8.
\]
\source{Math. Trip. I., 1912.}

\end{enumerate}

\section*{Exercises --- Polar Coordinates (Arts. 407--414)}

Find the areas bounded by

\begin{enumerate}[leftmargin=2em]

\item $r^2=a^2\cos^2\theta-b^2\sin^2\theta$, the central pedal of a hyperbola.

\item One loop of $r=a\sin4\theta$. Also state the total area.

\item One loop of $r=a\sin5\theta$. Also state the total area.

\item One loop of $r=a\sin n\theta$.

Give the total area in the cases, (i) $n$ even; (ii) $n$ odd.

\item The portion of $r=ae^{\theta\cot\alpha}$ bounded by the radii vectores
\[
\theta=\beta,\quad \theta=\beta+\gamma\quad(\gamma<2\pi).
\]

\item Any sector of $r^{\frac12}\theta=a^{\frac12}$ $(\theta=\alpha \text{ to } \theta=\beta)$.

\item Any sector of the reciprocal spiral $r\theta=a$ $(\theta=\alpha \text{ to } \theta=\beta)$.

\item The cardioide $r=a(1-\cos\theta)$.

\item The Lima\c{c}on $r=a+b\cos\theta$, (i) if $a>b$; (ii) if $a<b$ obtain the two areas of outer and inner portions.

\item Find the area included between the two loops of the curve
\[
r=a(2\cos\theta+\sqrt3).
\]
\source{Oxf. I. P., 1889.}

\item Prove that the area in the positive quadrant of the curve
\[
(x^2+y^2)^5=(a^2x^3+b^2y^3)^2\quad\text{is}\quad\frac13(a^2+b^2).
\]
\source{$\gamma$, 1899.}

\item Find the area of the closed part of the Folium
\[
r=\frac{3a\sin\theta\cos\theta}{\sin^3\theta+\cos^3\theta}.
\]
\source{I. C. S., 1884.}

\item Show that the area of a loop of the curve
\[
ax^{2n+1}-b^2x^ny^n+cy^{2n+1}=0
\]
is $\dfrac{1}{2(2n+1)}\dfrac{b^4}{ac}$, $a$ and $c$ being positive. \source{Colleges, 1881.}

\item Trace the curve whose equation is
\[
r^4=a^4\sec\theta\tan\theta,
\]
and find the area between the curve and any pair of radii vectores drawn from the pole. \source{Trinity, 1882.}

\item Trace the lemniscate $r^2=a^2\cos2\theta$ and its first positive pedal, and show that the area of a loop of the latter is double the area of a loop of the former.

Find the areas of each of the two small lozenge-shaped portions common to the two loops of the pedal.

\item Show that the area contained between the curve
\[
r=a\cos5\theta
\]
and the circle $r=a$ is three-fourths of the area of the circle. \source{Oxf. I. P., 1888.}

\item Find the area between the curve $r=a(\sec\theta+\cos\theta)$ and its asymptote. \source{St. John's, 1881.}

\item Prove that the area of the curve
\[
r^2(2c^2\cos^2\theta-2ac\sin\theta\cos\theta+a^2\sin^2\theta)=a^2c^2
\]
is equal to $\pi ac$. \source{I. C. S., 1879.}

\item Find the area of the curve
\[
r=3a\cos\theta+a\cos3\theta.
\]
\source{Math. Trip., 1882.}

\item Find the area of the loop of the curve
\[
r^2=a^2\theta\cos\theta
\]
between $\theta=0$ and $\theta=\dfrac{\pi}{2}$.

\end{enumerate}

\section*{General Problems on Quadrature (Cartesians and Polars)}

\begin{enumerate}[leftmargin=2em]

\item Find the area bounded by
\[
x^2+y^2=4a^2,\quad x^2+y^2=2ay\quad\text{and}\quad x=a.
\]
\source{H. C. S.}

Also the area of the loop of the curve
\[
by^2=x^2(a-x)
\]
($a$ and $b$ both positive). \source{I. C. S., 1882.}

\item Find the whole area of the curve
\[
y^2=x^2\frac{a^2-x^2}{a^2+x^2}.
\]
\source{I. C. S., 1885; Colleges, 1892.}

\item A parabola $y^2=ax$ cuts the hyperbola $x^2-y^2=2a^2$ at the points $P$, $Q$; and the tangent at $P$ to the hyperbola cuts the parabola again at $R$. Find the area of the curvilinear triangle $PQR$.

\item Find the area included between one of the branches of the curve $x^2y^2=a^2(x^2+y^2)$ and its asymptotes.

Find the whole area of the curve
\[
x^4+y^4=a^2(x^2+y^2).
\]
\source{Colleges $\alpha$, 1887.}

\item Trace the curve $a^2y^2=x^3(2a-x)$, and prove that its area is equal to that of the circle whose radius is $a$. \source{I. C. S., 1887.}

\item Prove that the whole area of
\[
(x^2+a^2)y^2+3a^3y+2a^4=0
\]
is
\[
(3-2\sqrt2)\pi a^2.
\]
\source{Colleges $\beta$, 1891.}

\item Find the area of the loops of the curve
\[
y^4-x^4-a^2y^2+b^2x^2=0\quad\text{when } b^2>a^2.
\]
\source{Oxford I. P., 1902.}

\item Find the area bounded by the cycloid
\[
x=a(\theta+\sin\theta),
\]
\[
y=a(1-\cos\theta),
\]
and the straight line joining two consecutive cusps.

\item Show that the coordinates of a point $P$ on the Folium of Descartes $x^3+y^3=axy$ can be expressed as
\[
x=\frac{at}{1+t^3},\quad y=\frac{at^2}{1+t^3}.
\]
Show that as $t$ varies from $0$ to $\infty$ $P$ traces out a closed loop, and that its area is $\dfrac{a^2}{6}$. \source{Colleges, 1896.}

\item Prove that the area of either loop of the curve
\[
x^5+y^5-5a^2x^2y=0
\]
is
\[
\frac{2\pi a^2}{\sqrt{10+2\sqrt5}}.
\]
\source{$\gamma$, 1893.}

\item Show that in that part of the curve $(x+y-3c)xy+c^3=0$ for which $x$ is positive, the area between the curve, the axis of $x$, and the ordinate which touches the curve is $\dfrac12c^2$. \source{St. John's, 1886.}

\item Trace the curve $y^4+x^3y=a^2x^2$,

and show that the area of the segment which lies between the axis of $y$ and the straight line whose equation is $y=x$ is $\dfrac16a^2\log2$. \source{Colleges $\epsilon$, 1883.}

\item Pairs of ordinates of the hyperbola $xy=a^2$ are determined by the condition that the area included by any pair, the curve, and the $x$-axis is constant; show that the lengths of any such pair are in a constant ratio. \source{Oxford I. P., 1888.}

\item Show that the area between the curve
\[
x(x^2+y^2-a^2)+\frac29a^3\sqrt3=0
\]
and its asymptote is $\pi a^2$. \source{St. John's, 1892.}

\item Show that the area between the inner branch of the curve
\[
(x^2+y^2-a^2)^2=\frac14x^2(x^2+y^2)
\]
and the positive parts of the two axes is $\pi a^2/3\sqrt3$. \source{St. John's, 1888.}

\item Prove that the whole area of the epicycloid generated by a point on a circle of radius $\dfrac{a}{4}$ rolling on a fixed circle of radius $a$ is to the area of the fixed circle in the ratio of 15 to 8.

\item Find the whole area of the curve whose equation is
\[
(x^2+y^2)(x+y+a)(x+y-a)+x^2y^2=0.
\]
\source{Colleges, 1886.}

\item Find the area of a loop of the curve
\[
x^4+y^4=2a^2xy.
\]
\source{Oxford I. P., 1888.}

\item Find the area cut off from an ellipse by a focal chord. \source{Colleges $\alpha$, 1883.}

\item Prove that the areas cut off by the equiangular spiral $r=ae^{\theta\cot\alpha}$ from the space bounded by any two fixed lines through the pole are in geometrical progression. \source{Oxford I. P., 1900.}

\item Find the area of the curve $r=a\theta e^{b\theta}$ enclosed between two given radii vectores and two successive branches of the curve. \source{Trinity, 1881.}

\item Find the area of the loop of the curve $r=a\theta\cos\theta$ between $\theta=0$ and $\theta=\dfrac{\pi}{2}$. \source{Oxford II. P., 1890.}

\item Find the area of the curve
\[
(r-a\cos\theta)^2=a^2\cos2\theta.
\]
\source{Colleges $\alpha$, 1887.}

\item Show that the area of the loop of the folium $x^3+y^3=3axy$ is divided by the parabola $y^2=ax$ in the ratio 5 : 4.

In what ratio does the line $x+y=2a$ cut the loop in the above folium. \source{Oxford I. P., 1889.}

\item Find the area included between the axis of $y$ and the curve
\[
y^2+2y-2x(y+1)=x^4-3x^3+3,
\]
the curve being supposed to stop at the node. \source{St. John's, 1884.}

\item Determine by integration the area of the ellipse
\[
x^2+xy+y^2=1.
\]

\item (i) Find the whole area enclosed by the hypocycloid
\[
x^{\frac23}+y^{\frac23}=a^{\frac23}.
\]
\source{Oxford I. P., 1888.}

(ii) Prove that the area of the locus of intersection of pairs of tangents at right angles for this curve is $\dfrac14\pi a^2$. \source{Math. Tripos, 1888.}

\item Prove that the locus of the points of bisection of the intercepts on the normals of a cycloid between the cycloid and its base divides the area between the cycloid and its base into two parts in the ratio 7 : 5. \source{Oxford II. P., 1886.}

\item Trace the curve $x^{2n+1}+y^{2n+1}=(2n+1)ax^ny^n$, when $n$ is even, and when $n$ is odd, $n$ being a positive integer; and prove that the area of the loop is $(2n+1)\dfrac{a^2}{2}$. Prove that this is also the area between the infinite branches of the curve and the asymptote. \source{St. John's, 1882.}

\item Find the whole area contained between the curve
\[
x^2(x^2+y^2)=a^2(y^2-x^2)
\]
and its asymptotes. \source{Oxford I. P., 1887.}

\item Find the area bounded by the circle $x=a\cos\theta$, $y=a\sin\theta$ and the hyperbola $x=b\cosh u$, $y=b\sinh u$; that area being taken which lies within the circle and on the convex side of the hyperbola, and $b$ being less than $a$. \source{Trinity, 1888.}

\item (a) Show that in the Archimedean Spiral $r=a\theta$, if $A_1,A_2,A_3,A_4,\dots$ be the areas of the inner loop and the successive heart-shaped figures formed by the convolutions of the curve
\[
A_1=\frac{\pi^2a^2}{4},\quad A_{n+1}=2n\pi^2a^2.
\]

(b) In the Reciprocal Spiral $r\theta=a$, if $A_1,A_2,A_3\dots$ be the areas of the successive closed loops,
\[
A_n=\frac{4a^2}{\pi}\frac{1}{4n^2-1}.
\]

\item Find the area of the loop of the curve
\[
(x+y)(x^2+y^2)=2axy.
\]
\source{Oxford I. P., 1890.}

\item At all points of the first negative pedal of the curve $r=\cosh(m\theta\cot\alpha)$ lines are drawn making a constant angle $\alpha$ with the tangent. Show that the area bounded by any pair of such lines, the curve enveloped and the first negative pedal is
\[
A\{1+(m^2-1)\cos^2\alpha\},
\]
where $A$ is the area of the corresponding portion of the first negative pedal bounded by radii vectores from the pole. \source{Colleges $\alpha$, 1891.}

\item Find the area of that portion of the loop of the curve
\[
r^2=p\cos\theta+q\sin\theta,
\]
which is not enclosed by the curve
\[
r^2=b+a\cos\theta.
\]
If a family of such curves be taken (by varying $p$ and $q$), such that this area is constant, show that the envelope of the system is a curve whose equation is
\[
r^2=c+a\cos\theta.
\]
\source{Colleges $\beta$, 1889.}

\item Show that the whole area enclosed by the outer line of the curve $r^{\frac23}=a^{\frac23}\cos\dfrac23\theta$ is $\dfrac98a^2\sqrt3$. \source{Colleges, 1876.}

\item In a hyperbola, $C$ is the centre, $A$ the end of the transverse axis and $P$ any point $(x,y)$ on the same branch of the curve as $A$; prove that twice the area of the sector $CAP$ is
\[
ab\log\left(\frac{x}{a}+\frac{y}{b}\right).
\]

\item Show that the area contained between a hyperbola, any tangent and a line parallel to the asymptote which bisects the part of the tangent intercepted between the curve and the asymptote
\[
=\frac{ab}{2}\left(\log2-\frac58\right),
\]
and is constant. \source{Trinity, 1886.}

\item Prove that the area of the curve
\[
x=\frac{ap}{(1+p^2)^2},\quad y=\frac12\frac{ap^2(1-p^2)}{(1+p^2)^2}
\]
is $\dfrac{1}{32}\pi a^2$. \source{Math. Tripos, 1882.}

\item Show that the area cut off from the ellipse
\[
ax^2+2hxy+by^2=1
\]
by the line $lx+my=1$ is
\[
\alpha\beta(\theta-\sin\theta\cos\theta),
\]
where $\alpha$, $\beta$ are the semiaxes of the ellipse and
\[
\cos\theta=\frac{\sqrt{ab-h^2}}{\sqrt{am^2+bl^2-2hlm}}.
\]
\source{Colleges, 1892.}

\item Trace the curve whose equation is
\[
x(x^{2n}+y^{2n})=a^2y^{2n-1},
\]
and prove that the area between the curve, the axis of $x$ and a tangent parallel to the axis of $y$ is
\[
\frac{a^2}{4n}(2n-1-\log2n).
\]
\source{St. John's, 1885.}

\item Show that in the curve
\[
r^2=\sec2\theta\log(2\cos^2\theta)
\]
the area between the curve and the lines $\theta=\pm\dfrac14\pi$ is $\left(\dfrac14\pi\right)^2$. \source{St. John's, 1886.}

\item Find in integral form, and completely, the area enclosed between two confocal conics and two given radii from the centre. \source{Trinity, 1881.}

\item Prove that the area of each of the two equal and similar pieces of the ellipse $x^2/a^2+y^2/b^2=1$ which are cut off by the hyperbola $x^2/\alpha^2-y^2/\beta^2=1$ $(\alpha<a)$ is
\[
ab\sin^{-1}\frac{\beta(a^2-\alpha^2)^{\frac12}}{(a^2\beta^2+\alpha^2b^2)^{\frac12}}-\alpha\beta\sinh^{-1}\frac{b(a^2-\alpha^2)^{\frac12}}{(\alpha^2\beta^2+a^2b^2)^{\frac12}}.
\]
\source{St. John's, 1887.}

\item Prove that the areas of the two loops of the curve
\[
r^2-2ar\cos\theta-8ar+9a^2=0
\]
are
\[
(32\pi+24\sqrt3)a^2\quad\text{and}\quad(16\pi-24\sqrt3)a^2.
\]
\source{Math. Tripos, 1875.}

\item The area between two tangents to the same convolution of an equiangular spiral at right angles to one another, and the curve, is
\[
pp'+\frac12(p^2-p'^2)\cot2\gamma,
\]
where $p$, $p'$ are the perpendiculars from the pole on the tangents and $\gamma$ is the angle of the spiral. \source{Colleges, 1882.}

\item A circle with centre at the origin cuts the loop of the Folium $x^3+y^3-3axy=0$. If the angle subtended at the origin by the common chord equals
\[
2\tan^{-1}\frac{2^{\frac13}-1}{2^{\frac13}+1},
\]
prove that the area between the loop and the circle is
\[
\frac{a^2}{2}\left[1-(2^{\frac23}+2^{\frac43})\tan^{-1}\frac{2^{\frac13}-1}{2^{\frac23}}\right].
\]
\source{Colleges, 1885.}

\item The centre of a circle of constant radius $a$ moves along a fixed straight line $AB$ in its plane, and from $A$ a fixed point in the line a tangent $AP$ is drawn to the circle. Show that the area included between the locus of $P$ and its asymptotes is $\pi a^2$. \source{Math. Tripos, 1882.}

\item Show that the curve
\[
r=a\left(\frac12\sqrt3+\cos\frac{\theta}{2}\right)
\]
has three loops, whose areas are
\[
a^2\left(\frac54\pi+2\sqrt3\right),\quad a^2\left(\frac56\pi-\frac54\sqrt3\right),\quad a^2\left(\frac{5}{12}\pi-\frac34\sqrt3\right)
\]
respectively. \source{Colleges, 1892.}

\item Show that the area of the Cassinian
\[
r^4-2a^2r^2\cos2\theta+a^4=b^4
\]
is
\[
2\int_0^{\pi/2}\sqrt{b^4-a^4\sin^2\phi}\,d\phi,\quad\text{provided } b>a,
\]
but is
\[
2\int_0^{\pi/2}\frac{b^4\cos^2\phi\,d\phi}{\sqrt{a^4-b^4\sin^2\phi}},\quad\text{when } a>b.
\]
\source{Math. Tripos, 1883.}

\item Prove that the area of the first negative pedal of an ellipse with respect to the focus is
\[
\frac{\pi a^2(2-3e^2)}{2\sqrt{1-e^2}}\quad(e<\tfrac12),
\]
where $a$ and $e$ are the semi major axis and the eccentricity of the ellipse. \source{Colleges, 1892.}

How do you interpret this result if $e<\tfrac12$?

\item Find the area of the curve whose Cartesian equation is
\[
a^2(y-x)^2=(a+x)^3(a-x).
\]
\source{Math. Tripos, 1896.}

\item Find the value of $\displaystyle\int_0^1 v_x\,dx$, $v_x$ being the real root of the cubic
\[
v_x^3+v_x^2v_1+v_xv_1^2-\frac{c}{x}=0.
\]
\source{Colleges, 1872; R. P.}

\item Find the area in the first quadrant bounded by the axes of coordinates and the curve
\[
\sinh^{-1}\frac{x}{a}+\sinh^{-1}\frac{y}{b}=c,
\]
taking $a$, $b$, $c$ all positive. \source{I. C. S., 1897.}

\item Trace the whole curve
\[
x^2y^2=c^2(a-x)(x-b),
\]
where $0<b<a$, and find its whole area. \source{I. C. S., 1898.}

\item It is given that the abscissa $ON$ and ordinate $NP$ of a point on any arch of a cycloidal arc are $a(\theta-\sin\theta)$ and $a(1-\cos\theta)$. $NP$ is produced to $K$ so that $NK=2a$, and the rectangle $ONKA$ is completed. Prove that the area included by $ON$, $NP$ and the arc $OP$ never differs from three-fourths of $ONKA$ by more than
\[
\frac{3a^2}{8}\sqrt3;
\]
and find for what positions of $P$ the difference vanishes. \source{I. C. S., 1912.}

\item Trace on squared centimetre paper the curves
\[
x^4+y^4=4a^2xy,
\]
\[
x^4+y^4=4ax^2y,
\]
taking $a=10$ cm., and estimate the area of a loop of each curve.

Prove that
\[
\int_0^\infty\frac{t^2}{(1+t^4)^2}\,dt=\frac14\int_0^\infty\frac{t^2}{1+t^4}\,dt=\frac{\pi}{8\sqrt2},
\]
and hence calculate the area of a loop of the second curve. Find also the area of a loop of the first curve. Give each area to the nearest square centimetre when $a$ is 10 centimetres. \source{C. S., 1913.}

\item Obtain the area contained between the two curves
\[
r^2\cos2\theta=4a^2\cos^4\theta\quad\text{and}\quad r^2\cos2\theta=a^2.
\]
\source{Oxf. I. P., 1912.}

\item Show that the area of the loop of the curve
\[
x^7+y^7=ax^3y^3
\]
is equal to $a^2/14$. \source{Oxf. I. P., 1914.}

\item Prove by any method that the area of the ellipse
\[
\{a(x-2)+3y\}^2+4(x+1)(x-2)=0
\]
is independent of $a$, and find the area.

Prove also that the straight line $y=x$ divides the ellipse $x^2+3y^2=6y$ into two areas which are in the ratio
\[
4\pi-3\sqrt3:8\pi+3\sqrt3.
\]
\source{Oxf. I. P., 1916.}

\item Trace the curve
\[
r\cos\theta=a\sin3\theta,
\]
and show that the area of a loop is
\[
\frac18a^2(9\sqrt3-4\pi).
\]
\source{Math. Trip. I., 1919.}

\item Show that the curve $r=a(2\cos\theta+\cos3\theta)$ has three loops, the area of the larger loop being $\dfrac{10\pi+9\sqrt3}{12}a^2$, and the areas of the two smaller loops being $\dfrac{5\pi-9\sqrt3}{24}a^2$. \source{Math. Trip. I., 1916.}

\item Show that the coordinates of any point on the curve
\[
y^2(a+x)=x^2(3a-x)
\]
may be taken as
\[
x=a\sin3\theta/\sin\theta,\quad y=a\sin3\theta/\cos\theta,
\]
and prove that the area of the loop and the area between the curve and its asymptote are both equal to $3\sqrt3a^2$. \source{Math. Trip. I., 1915.}

\item Show that the area of the loop of the curve
\[
(x^2+y^2)^2-4axy^2=0
\]
in the positive quadrant is $\dfrac14\pi a^2$. \source{Math. Trip. I., 1920.}

\item Having established Simpson's Rule, that if
\[
y=y(x)\equiv a_0+a_1x+a_2x^2+a_3x^3,
\]
then
\[
\int_0^1 y\,dx=\frac16\{y(0)+y(1)+4y(\tfrac12)\},
\]
prove that if $y(x)$ also contains a term $a_4x^4$ the error in still using Simpson's Rule is
\[
\frac{1}{120}a_4.
\]
\source{Math. Trip. I., 1920.}

\end{enumerate}

\newpage

\begin{center}
    {\LARGE \bfseries \color{accent} Chapter XIII: Quadrature (II). Tangential Polars, Pedal Equations and Pedal Curves, Intrinsic Equations, etc. --- Exercises}\\
    \vspace{0.2cm}
    {\color{accent}\rule{0.55\linewidth}{0.8pt}}
\end{center}
\vspace{0.35cm}

\section*{Problems on Quadrature (Arts. 416--450)}

\begin{enumerate}[leftmargin=2em]

\item Interpret geometrically $\displaystyle\int_{p_0}^{p_1}\sqrt{r^2-p^2}\,dp$ in the case of the curve $r=f(p)$.

Prove that the value of $\displaystyle\int\sqrt{r^2-p^2}\,dp$, taken all round an ellipse whose semiaxes are $a$, $b$, and whose centre is the pole, is $\pi(a-b)^2$. \source{Oxford I. P., 1903.}

\item Use the pedal equation of an ellipse, viz. $\dfrac{a^2b^2}{p^2}=a^2+b^2-r^2$, to show that the area of the portion of an ellipse included between the curve, the semi-major axis and a central radius vector $r$, is
\[
\frac{ab}{2}\tan^{-1}\sqrt{\frac{a^2-r^2}{r^2-b^2}},
\]
$a$, $b$ being the semiaxes of the ellipse. \source{Colleges, 1882.}

\item Find the area of the part of the ellipse $p^2(2a-r)=b^2r$ included between two focal radii vectores drawn, one to an extremity of the minor axis and the other to the nearer extremity of the major axis. \source{Oxford I. P., 1889.}

\item Find the area included between an ellipse and its evolute and bounding radii of curvature, the one coinciding with the major axis and the other inclined at an angle of $\dfrac{\pi}{4}$ to it. \source{Colleges, 1884, and $\beta$, 1888.}

\item Through every point of an ellipse a line is drawn outwards normal to the ellipse and equal to the radius of curvature at the point. Show that the area of the curve thus obtained is
\[
\pi\,\frac{9a^4+14a^2b^2+9b^4}{2ab}.
\]
\source{Colleges $\alpha$, 1891.}

\item Show that the area of that part of the evolute of an ellipse $\left(\text{eccentricity}>\dfrac{1}{\sqrt2}\right)$ which lies outside the ellipse is
\[
a^4b^4\int_{b^2}^{\frac{a^2+b^2}{3}}\frac{(a^2+b^2-3\rho)^2}{\rho^3(a^2+b^2-\rho)^2}\frac{d\rho}{\sqrt{(\rho-b^2)(a^2-\rho)}}.
\]
\source{Colleges, 1882.}

\item Find the area of the pedal of the curve
\[
(ax)^{\frac23}+(by)^{\frac23}=(a^2-b^2)^{\frac23},
\]
the origin being taken at $x=\sqrt{a^2-b^2}$, $y=0$. \source{Oxford I. P., 1888.}

\item Show that the area of the space between the epicycloid $p=A\sin B\psi$ and its pedal curve taken from cusp to cusp is $\dfrac14\pi A^2B$. \source{Colleges, 1878.}

\item Show that the area between an epicycloid and the arc of the fixed circle included between two consecutive cusps is
\[
\frac{\pi b^2}{a}(3a+2b),
\]
where $a$ and $b$ are the radii of the fixed and rolling circles respectively. \source{Colleges $\alpha$, 1884.}

Show also that the area of the corresponding sector of the fixed circle is that of an ellipse with semiaxes the radii of the two circles. \source{Oxford I. P., 1913.}

\item Show that the $p$-$\psi$ equation to a cycloid when one of the cusps is taken as origin is
\[
p=2a(\sin\psi-\psi\cos\psi),
\]
where $a$ is the radius of the generating circle; and find the area between the curve from cusp to cusp and the corresponding arc of the pedal with regard to a cusp. \source{Oxford II. P., 1903.}

\item Show that the area bounded by that portion of the cardioide $r^{\frac12}=a^{\frac12}\sin\dfrac12\theta$, which lies in the first quadrant, the terminal tangents, and the corresponding portion of the locus of the extremity of the polar subtangent, is
\[
3a^2(10-3\pi)/16.
\]
\source{Math. Tripos, 1896.}

\item Show that in the curve in which the area bounded by the curve and the radii vectores from a certain fixed point varies as the square of the length of the bounding arc, the radius of curvature varies as the projection of the radius vector on the tangent. \source{Colleges $\alpha$, 1891.}

\item The pedal of a cycloid with regard to any point on its axis meets the cycloid at the vertex $A$ and cuts the tangent at the cusp in $Q$; find the area between it and the chord $AQ$; and prove that this area is least when the origin is the middle point of the axis. \source{St. John's, 1883.}

\item An elliptic wire is pushed in one plane through a very short straight tube; find the equation to the locus of the centre, and prove that the area of each loop is $\dfrac{\pi}{2}(a-b)^2$, where $a$ and $b$ are the semiaxes. \source{Colleges, 1886.}

\item A point $Q$ is taken on the normal drawn outward at a point $P$ of a catenary, the parameter of which is $c$. Prove that if $PQ$ is equal to the length of the arc of the catenary measured from the vertex to $P$, the area between the locus of $Q$ and the catenary, and bounded by the normal at the vertex and by another normal inclined at an angle $\psi$ to this, is
\[
\frac{c^2}{2}(\tan^2\psi+\tan\psi-\psi).
\]
\source{Colleges $\gamma$, 1882.}

\item Prove that the pedal of the cardioide $r=a\cos^2\dfrac{\theta}{2}$ with respect to the cusp consists of two closed regions of areas $A$ and $B$, $A$ consisting of the inner loop and $B$ being external to $A$ and bounded by the outer line of the curve and such that
\[
2A+B=\frac{15\pi a^2}{32}.
\]
\source{Colleges $\gamma$, 1899.}

\item Prove that the area of the pedal of the curve $x^{\frac23}+y^{\frac23}=a^{\frac23}$ with respect to the point $(a,0)$ is five times as great as the area of its pedal with respect to the origin. \source{Oxford II. P., 1899.}

\item The tangent at a point $P$ of a lemniscate cuts the curve again at $Q$, $R$. Prove that the middle point of $QR$ is at the same distance from the nodal point as $P$; and that the equation to its locus is
\[
a^{10}(x^2-y^2)=r^4\{a^8+4(a^4-r^4)(a^4-4r^4)\},
\]
where $r^2\equiv x^2+y^2$.

Show that it can be written $r^2=a^2\cos\dfrac25\theta$.

Trace the curve completely, and prove that the portion corresponding to the upper half of one branch of the lemniscate divides the other branch into two parts whose areas are in the ratio of
\[
6-3\sqrt3:3\sqrt3-4.
\]
\source{St. John's, 1884.}

\item Show that the area of a loop of the curve
\[
(x^2-a^2)^2+(y^2-3a^2)^2=a^4
\]
is
\[
a^2\sqrt2\left(\frac{\pi}{3}-\log_e\frac{\sqrt3+1}{2}\right).
\]
\source{Math. Tripos, 1882.}

\item The tangent at every point $P$ of a closed finite curve is produced to $Q$ so that $PQ$ is constant. Find the area between the locus of $Q$ and the original curve. How is the result to be explained, (i) if the curvature of the first curve is sometimes in one direction, sometimes in the opposite direction; (ii) if the curve cuts itself a given number of times. \source{St. John's Coll., 1881.}

\item A straight line of constant length $c$ is drawn from each point of a closed oval curve making a given angle $\alpha$ with the normal at that point. Prove that the area of the curve traced out by the end of the line is
\[
S+\pi c^2\pm lc\cos\alpha,
\]
where $S$ is the area of the given oval curve and $l$ is its length. \source{Coll. $\gamma$, 1893.}

\item Show that the area of the polar reciprocal of a curve whose equation is given in rectangular coordinates is
\[
\frac12k^4\int\frac{\dfrac{d^2y}{dx^2}}{\left(y-x\dfrac{dy}{dx}\right)^2}dx,
\]
$x$, $y$ being the coordinates of a point on the original curve.

Apply this to find the area of the ellipse $\dfrac{x^2}{a^2}+\dfrac{y^2}{b^2}=1$. \source{Colleges, 1886.}

\item The area of a given closed oval curve is $A$; the bisectors of the internal and external angles between tangents to it which meet at a given constant angle $2\alpha$ envelop curves whose areas are $A_1$ and $A_2$; show that
\[
A_1\cos^2\alpha+A_2\sin^2\alpha=A.
\]
\source{Colleges $\gamma$, 1888.}

\item Prove that for any closed curve which has a centre, the area of the locus of intersection of tangents at right angles, and the area of the locus of intersection of normals at right angles differ by twice the area of the curve. \source{Math. Tripos, 1888.}

\item $O$ being a fixed point, $OP$ a radius vector of any curve, $OP$ is produced to $Q$ so that $OP.PQ=a^2$, and $A$ is the area between the locus of $Q$ and the given curve. If $A'$ be the area of the inverse of the curve with respect to $O$, the constant of inversion being $a$, show that $A'-A$ is independent of the form of the curve.

If the given curve be a circle, and $O$ a point on its circumference, find the area of any part bounded by the locus of $Q$, the circle and two radii vectores from $O$. \source{St. John's, 1891.}

\item A circle rolls on the outside of an oval curve, the pedals of the curve, of the locus of the centre of the circle and of the envelope of the circle are of areas $A_0$, $A_1$, $A_2$, respectively; prove that $A_2-2A_1+A_0$ depends only on the rolling circle.

Show that if the area of the oval curve, of the locus of the centre of the circle and of the envelope of the circle be $S_0$, $S_1$, $S_2$ respectively,
\[
A_2-2A_1+A_0=S_2-2S_1+S_0.
\]
\source{Trinity, 1878.}

\item One of the curves given by the equation
\[
y=a^2\frac{d}{ds}\left\{\frac{dy}{dx}+\frac13\left(\frac{dy}{dx}\right)^3\right\}
\]
cuts the axis of $x$ twice at the angle $\alpha$. Prove that the area between the curve and the axis is
\[
a^2\{\tan\alpha\sec\alpha+\log(\sec\alpha+\tan\alpha)\}.
\]
\source{Oxf. I. P., 1912.}

\item A curve concave to the axis of $x$ is such that the product of the ordinate and radius of curvature at any point is constant and equal to $c^2$ (The Elastica, or Bent Bow). Prove that the maximum value of the ordinate is $2c\sin\dfrac{\alpha}{2}$, where $\alpha$ is the angle at which the curve crosses the axis of $x$. \source{Ox. I. P., 1903.}

Show that the area which lies between the bow and the bow-string is $2c^2\sin\alpha$.

\item Show that the area of a closed curve, which is the envelope of the line $x\cos\psi+y\sin\psi=p$, is the value of the integral
\[
-\frac12\int\left(\frac{dp}{d\psi}+p\right)\left(\frac{dp}{d\psi}-p\right)d\psi
\]
taken completely round the curve. \source{Math. Trip., 1898.}

\item The integral $-\dfrac12\displaystyle\int\left(\dfrac{dp}{d\psi}+np\right)^2d\psi$ is taken round a closed curve, $n$ being taken equal to $\tan\psi$ or to $-\cot\psi$, according as the one or the other is numerically less than unity. Show that the value of the integral differs from the area of the curve by the sum of the squares of the perpendiculars from the origin upon the tangents at the points where the integral changes form. \source{Math. Trip., 1898.}

\item In the cycloid prove that the conic locus of points with regard to which the area of the pedal is constant, is in general a circle, and find the point for which the area of the pedal is a minimum. \source{Ox. I. P., 1900.}

\item In a catenary, $A$ is the vertex, $P$ any point on the curve, $AO$, $PN$ perpendiculars upon the directrix, $PY$ a tangent and $NY$ perpendicular to it. Show that the area of the figure $ONPA$ is double that of the triangle $YNP$.

\item Show that the area of the first positive pedal of the curve $p=f(r)$ may be obtained by the formula
\[
\frac12\int\frac{p^2}{\sqrt{r^2-p^2}}\frac{dp}{dr}\,dr,
\]
where the letters $p$ and $r$ are the pedal coordinates of a point on the original curve.

Apply this method to find the area of the cardioide, which is the first positive pedal of the circle $r^2=ap$.

\item Employ the formula
\[
\frac12\int\frac{pr}{\sqrt{r^2-p^2}}\,dr
\]
to find the area of
\[
r^2+a^2=b^2+2ap\quad(a>b).
\]
To what curve does this pedal equation belong?

\item In the epicycloid
\[
p^2=a^2\frac{r^2-a^2}{c^2-a^2},
\]
where $a$ and $\dfrac{c-a}{2}$ are the radii of the fixed and rolling circles respectively, obtain a formula for the area of any sectorial portion with centre of the sector at the origin. Hence show that the area between one foil of the curve and the fixed circle is
\[
\pi(c-a)^2(c+2a)/4a.
\]

\item When $a<b$ the conchoid of Nicomedes, viz.
\[
x^2y^2=(a+y)^2(b^2-y^2)\quad\text{or}\quad r=a\csc\theta\pm b
\]
has a loop. Find its area.

\item Let $S$ be the focus of a parabola, $SP_1$, $SP_2$ two focal radii vectores of lengths $r_1$, $r_2$. The latus rectum is $4a$ and $P_1P_2=c$. Prove Lambert's expression for the sectorial area $SP_1P_2$, viz.
\[
\frac{\sqrt a}{3}\left[s^{\frac32}-(s-c)^{\frac32}\right],
\]
where $2s=r_1+r_2+c$.

Show that the segment cut off by a focal chord of length $c$ is
\[
\frac13a^{\frac12}c^{\frac32}.
\]

\item In the case of the Cotes's spirals, whose equations are of the form
\[
\frac{1}{p^2}=\frac{A}{r^2}+B,
\]
show that the area of the sectorial portion bounded by the curve and the radii vectores $r_1$ and $r_2$ is
\[
\frac{1}{2B}\left\{(Br_1^2+A-1)^{\frac12}\sim(Br_2^2+A-1)^{\frac12}\right\},\quad B\neq0.
\]

Examine in detail the particular cases of

(i) the equiangular spiral;

(ii) the reciprocal spiral;

(iii), (iv) and (v) the cases which reduce to the polar forms,
\[
u=a\sinh n\theta,\quad u=a\cosh n\theta,\quad u=a\sin n\theta,\quad\text{respectively.}
\]

\item Riccati's \textit{Syntractory} is generated as follows. The tractory is an involute of a common catenary of parameter $c$, starting from the vertex. $PT$ is a tangent at any point $P$ of the tractory, cutting the directrix of the catenary at $T$. $Q$ is a point on $PT$ or $PT$ produced such that $QT=c'$. The locus of $Q$ is the syntractory.

Show that the areas between the two branches and the directrix are
\[
\frac{\pi}{2}c'(2c\pm c').
\]

\item If $A$ be the area of the `Helmet,'
\[
(k+1)\{(x^2+ka^2)y^2-2ay(a^2-x^2)\}+(a^2-x^2)^2=0,\quad(k\neq-1),
\]
and $V$ the volume formed by its revolution about the $y$-axis, prove that
\[
A=\frac{\pi a^2}{\sqrt{k(k+1)}}\left[2(k+1)^{\frac32}-(2k+3)k^{\frac12}\right],
\]
\[
V=\frac{2\pi a^3}{3\sqrt{k+1}}\left[3(k+1)^{\frac32}\log\frac{\sqrt{k+1}+1}{\sqrt{k+1}-1}-2(3k+4)\right].
\]

[For the first part of the example, and for several others of similar character, see Wolstenholme's \textit{Problems}, Nos. 1866 to 1870.]

\end{enumerate}

\newpage

\begin{center}
    {\LARGE \bfseries \color{accent} Chapter XIV: Quadrature, etc. (III). Surface Integrals, Areals, Corresponding Curves --- Exercises}\\
    \vspace{0.2cm}
    {\color{accent}\rule{0.55\linewidth}{0.8pt}}
\end{center}
\vspace{0.35cm}

\section*{Exercises --- Surface Integrals, Centroids and Moments of Inertia (Arts. 451--459)}

\begin{enumerate}[leftmargin=2em]

\item In the first quadrant of the circle $x^2+y^2=a^2$ the surface density varies at each point as $xy$.

Find (i) the mass of the quadrant,

(ii) its centroid,

(iii) its moment of inertia about the $y$-axis.

\item Work out the corresponding results for the portion of the parabola $y^2=4ax$ bounded by the axis and the latus rectum, the surface density varying as $x^py^q$.

\item Find the centroid of a fine rod of uniform sectional area and of which the line-density varies as the $n$th power of the distance from one end. Also its moment of inertia about that end, about the other end, and about the middle point.

\item Find the centroid of the triangle bounded by the lines $y=mx$, $x=a$ and the $x$ axis when the surface density at each point varies as the square of the distance from the origin. Also find the moment of inertia about the $y$-axis.

\item Find the centroid of

(i) either of the areas bounded by the circle $(x-a)^2+y^2=a^2$ and the parabola $y^2=ax$;

(ii) the centroid of the area bounded by the parabolas
\[
y^2=4ax,\quad x^2=4by;
\]

(iii) the centroid of the area bounded by
\[
y^2=4ax,\quad y=2x,
\]
the surface density being uniform in each case.

\item Find the moment of inertia of a triangle of uniform surface density

(i) about one of its sides;

(ii) about an axis perpendicular to its plane through an angular point.

\item Find the area remote from the pole between the circles
\[
r=a,\quad r=2a\cos\theta;
\]
and assuming a surface density varying inversely as the distance from the pole, find

(1) the centroid;

(2) the moment of inertia about a line through the pole perpendicular to the plane.

\item Find for the area included between the curves
\[
y^2=4ax,\quad x^2=4ay,
\]

(i) the moment of inertia about the $x$-axis;

(ii) the moment of inertia about an axis through the origin and at right angles to the plane of the area.

\item Find the coordinates of the centroid of the area bounded by the catenary $y=c\cosh\dfrac{x}{c}$, an ordinate, and the coordinate axes.

\item If the density at any point of a circular disc whose radius is $a$ vary directly as the distance from the centre and a circle described on a radius as diameter be cut out, prove that the centroid of the remainder will be at a distance $\dfrac{6a}{5(3\pi-2)}$ from the centre. \source{Math. Trip., 1875.}

\end{enumerate}

\section*{Exercises --- Corresponding Points and Areas (Arts. 460--465)}

\begin{enumerate}[leftmargin=2em]

\item Find the whole area of a loop of each of the curves

(i) $x(x^2+y^2)=a(x^2-y^2)$.

(ii) $(m^2x^2+n^2y^2)^2=a^2x^2-b^2y^2$.

\source{St. John's, 1887.}

\item Trace the shape of the following curves, and find their areas:

(i) $(x^2+y^2)^3=axy^4$.

(ii) $(x^2+2y^2)^3=axy^4$.

\source{Barnes Scholarships, 1887.}

\item Prove that the area of
\[
\frac{x^2}{a^4}+\frac{y^2}{b^4}=\frac{1}{c^2}\left(\frac{x^2}{a^2}+\frac{y^2}{b^2}\right)^2\quad\text{is}\quad\frac{\pi c^2}{2ab}(a^2+b^2).
\]

\item Prove that the area in the positive quadrant of the curve
\[
(a^2x^2+b^2y^2)^{\frac32}=mx^3+ny^3\quad\text{is}\quad\frac{1}{3ab}\left(\frac{m}{a^3}+\frac{n}{b^3}\right).
\]
\source{$a$, 1890.}

\item Prove that the area of the curve
\[
(a^2x^2+b^2y^2)^2=c^6(x^2-y^2)\quad\text{is}\quad\frac{c^6}{a^3b^3}\left\{ab+(b^2-a^2)\tan^{-1}\frac{b}{a}\right\}.
\]
\source{St. John's, 1883.}

\item Show that the area of the loop of the curve
\[
\frac{x^5}{a^5}+\frac{y^5}{b^5}=5\frac{x^2y^2}{a^2b^2}\quad\text{is}\quad\frac52ab.
\]

\item Find the area of the curve
\[
\frac{l}{\sqrt{r^2-c^2}}=1+e\cos\theta\quad(e<1).
\]

\item Show that the area bounded by
\[
(x^2+y^2-c^2)(x^2+y^2)=4a^2x^2\quad\text{is}\quad(2a^2+c^2)\pi.
\]

\item Find the area included within the curve whose equation is
\[
\left(\frac{x}{a}\right)^{\frac23}+\left(\frac{y}{b}\right)^{\frac23}=1.
\]
\source{Colleges, 1885.}

\item Trace the curve
\[
\left(\frac{x}{a}+\frac{y}{b}\right)^{\frac23}+\left(\frac{x}{a}-\frac{y}{b}\right)^{\frac23}=2,
\]
and show that its area is half as great again as that of the ellipse
\[
\frac{x^2}{a^2}+\frac{y^2}{b^2}=1.
\]
\source{Math. Tripos, 1884.}

\item Prove that the area of the curve
\[
(x^2+y^2)^5=a^2y(x^7+y^7)\quad\text{is}\quad\frac{105}{512}\pi a^2+\frac{125}{384}a^2.
\]
\source{St. John's, 1889.}

\item Prove that the area of the curve
\[
(a^2x^2+b^2y^2)^5=8a^4b^4xy(a^4x^6+b^4y^6)\quad\text{is}\quad a^2+b^2.
\]
\source{St. John's, 1889.}

\item Show that the area in the first quadrant of the curve
\[
c^4\left(\frac{x^2}{a^2}+\frac{y^2}{b^2}\right)^5=\left(\frac{x^3}{a}+\frac{y^3}{b}\right)^2\quad\text{is}\quad\frac{ab(a^2+b^2)}{3c^2}.
\]

\item Trace the curve $4(x^2+2y^2-2ay)^2=x^2(x^2+2y^2)$, proving that the area of a loop is $4\pi(2-\sqrt3)a^2/\sqrt3$, and that the area included between the loops is
\[
8a^2(2\pi-3\sqrt3)/3\sqrt3.
\]
\source{Trinity, 1896.}

\item Find the whole area of the curve
\[
\left(\frac{x^2}{a^2}+\frac{y^2}{b^2}\right)^2=\frac{2xy}{ab},
\]
\source{Oxford I. P., 1890.}

and of a loop of the curve
\[
\frac{x^4}{a^4}+\frac{y^4}{b^4}=\frac{2xy}{ab}.
\]
\source{Oxford II. P., 1900.}

\item Show that the area of either oval of
\[
x^2\{x^2/a^2+y^2/b^2-1\}+c^2=0\quad\text{is}\quad\frac12\pi b(a-2c).
\]
\source{St. John's, 1890.}

\item If $f(x,y)=0$ be a closed curve, show that its area is $mn$ times the area of the closed curve $f(mx,ny)=0$. Trace the curve $(4x^2+9y^2)^4=axy^6$, and find its area. \source{Oxford II. P., 1890.}

\item Trace the curve $\dfrac{x^3}{a^3}+\dfrac{y^3}{b^3}=\dfrac{3xy}{ab}$, and show that the area of its loop is $\dfrac32ab$.

\item A curve is defined by the equations
\[
x=6a\sin^2\phi,\quad y=6a\sin^2\phi\tan\phi,
\]
where $\phi$ is a variable parameter. Show that the centroid of the portion enclosed between the infinite branches and the asymptote is situated on the $x$-axis at a distance $5a$ from the origin. \source{Oxford II. P., 1889.}

\item (i) In an involute of a circle, show that the area swept out by the radius vector drawn from the centre of the circle to a point on the curve varies as the cube of the central perpendicular upon the tangent, the initial line being the radius to the point where the involute meets the circle.

(ii) In the Conchoid of Nicomedes $r=a\sec\theta-b$ in the case when $a<b$, show that the area of the loop is
\[
a^2(a\sec^2\alpha-2\sec\alpha\cosh^{-1}\sec\alpha+\tan\alpha),
\]
and that the distance of the centroid of the loop from the node is
\[
\frac23a\,\frac{3a\sec\alpha-3\cosh^{-1}\sec\alpha-\sin\alpha\tan^2\alpha}{a\sec\alpha-2\cosh^{-1}\sec\alpha+\sin\alpha},
\]
where $\alpha=\cos^{-1}a/b$.

\item Prove that the area contained by the curve
\[
x^4+2x^2y^2+4ax^2y+2a^2(y^2-x^2-2ay)+a^4=0\quad\text{is}\quad\pi a^2(4-5/\sqrt2).
\]
Find also the distance from the axis of $y$ of the centre of gravity of that portion of the area which lies in the first quadrant. \source{Colleges $\beta$, 1890.}

\item Show that the area included between the curve
\[
s=a\tan\psi,
\]
its tangent at $\psi=0$ and its tangent at $\psi=\phi$
is
\[
\frac12a^2\tan\phi+a^2\tan\frac12\phi-a^2\log(\sec\phi+\tan\phi).
\]
\source{Trinity, 1892.}

\item Show that an expression for the element of area in trilinear coordinates is
\[
\csc C\,d\alpha\,d\beta.
\]
Show that the area of the conic whose trilinear equation is
\[
a^{-1}\beta\gamma+b^{-1}\gamma\alpha+c^{-1}\alpha\beta=0
\]
is to that of the triangle of reference as
\[
4\pi:3\sqrt3.
\]
\source{Oxford II. P., 1890.}

\item Show that the coordinates of the centroid of the area bounded by half the cycloid $x=a(\theta+\sin\theta)$, $y=a(1-\cos\theta)$, the line of cusps and the $y$-axis are given by
\[
\frac{9\pi\bar x}{(3\pi-4)(3\pi+4)}=\frac{3\bar y}{7}=\frac{a}{2}.
\]
\source{Wallis.}

\item $OB$ and $OC$ are any two semi-diameters of an ellipse conjugate to each other; find the locus of the intersection of the normals at $B$ and $C$, and show that the area of the curve is
\[
\frac{\pi(a^2-b^2)^2}{4ab}
\]
\source{R. P.}

\item Tangents to a system of similar and similarly situated concentric ellipses are drawn such that the distance of each from the centre is the same. Find the area of the curve formed by the points of contact. \source{Trinity, 1885.}

\item Show that the moment of inertia of the portion of a uniform parabolic lamina cut off by the latus rectum about the tangent at an extremity of the latus rectum, is equal to $\dfrac{12Ma^2}{7}$, $4a$ being the latus rectum and $M$ the mass of the lamina. \source{Oxf. I. P., 1914.}

\item Prove by integration that the moment of inertia of a uniform triangular lamina $ABC$ of mass $M$ about a perpendicular axis at $A$ is
\[
\frac{1}{12}M(3b^2+3c^2-a^2).
\]
\source{Ox. I. P., 1915.}

\end{enumerate}

\newpage

\begin{center}
    {\LARGE \bfseries \color{accent} Chapter XV: Quadrature, etc. (IV). Miscellaneous Theorems, Connexion of a Line-Integral and a Surface-Integral, Mechanical Integration, etc. --- Exercises}\\
    \vspace{0.2cm}
    {\color{accent}\rule{0.55\linewidth}{0.8pt}}
\end{center}
\vspace{0.35cm}

\section*{Examples (Arts. 466--509)}

\begin{enumerate}[leftmargin=2em]

\item $Ox$, $Oy$ being perpendicular axes, $A$, $B$ fixed points on $Oy$ and $AMBA$ any closed region of area $S$ lying in the positive quadrant, show that the integral
\[
\int\big[\{\phi(y)e^x-my\}\,dx+\{\phi'(y)e^x-m\}\,dy\big],
\]
taken round the curve from $A$ to $B$, is equal to
\[
m(S+a-b)+\phi(b)-\phi(a),
\]
$\phi(y)$, $\phi'(y)$ being finite and continuous, $m$ a constant and $OA=a$, $OB=b$. \source{J. Math. Schol. Oxford, 1904.}

\item $P_1$, $P_2$ are points on a closed oval of area $A$, such that $P_1$, $P_2$ subtends a right angle at a fixed point $O$. Show that the area of the curve traced out by the middle point of $P_1P_2$ is equal to
\[
\frac12A+\frac18\int_0^{2\pi}\left(r_1\frac{dr_2}{d\theta_1}-r_2\frac{dr_1}{d\theta_2}\right)d\theta_1,
\]
where $\theta_2=\theta_1+\dfrac{\pi}{2}$ and $OP_1=r_1$, $OP_2=r_2$. \source{Colleges $\beta$, 1889.}

\item A fixed point $O$ is taken on a central oval which is such that through any point inside it other than the centre one and only one chord can be drawn which is bisected at that point; prove that the locus of the middle point of the chord $PQ$ for a constant sum $2\sigma$ of the arcs $OP$, $OQ$ cuts at right angles the same locus for a constant difference $2\sigma'$ of these arcs; and deduce that the area of the oval is
\[
\frac12\int_0^l d\sigma\int_0^{\frac12l}d\sigma'\sin\theta,
\]
where $l$ is the length of the oval, and $\theta$ is the angle between the tangents at $P$ and $Q$. \source{Math. Tripos, 1889.}

\item A bar $AB$ carries at a point of its length a small wheel having $AB$ for axis and which turns about $AB$: the end $A$ is constrained to move in a given straight line; show that if the end $B$ is carried round any closed curve without singular points and which does not cut the straight line on which $A$ moves, the area of the curve is measured by the product of $AB$ into the whole length registered by the revolving wheel. \source{Colleges, 1892.}

[This is the principle of construction of Coffin's Planimeter. A full description will be found on p. 159, \textit{Practical Electrical Engineering}, by Briggs and others. It is the case when the rod $OP$ of Fig. 102 is of infinite length, so that $P$ describes a straight line instead of a circle.]

\item A straight line of given length moves with its extremities on the arcs of two closed curves of given areas, and a point is attached to the moving line.

Prove that when the area traced by this attached point has a minimum value for different positions of the point on the line, the difference of the areas of the circles whose radii are the segments into which the point divides the line is equal to the difference of the areas of the given curves. \source{St. John's, 1882.}

\item Show that the path of the mid-point of a rod of constant length $2c$, whose ends lie upon an ellipse, is an oval of area $\pi(ab-c^2)$.

If, instead of both ends being on the ellipse, one end lies on the ellipse and the other on the major axis, or if one end lies on the ellipse and the other on the auxiliary circle, find the areas of the paths described by the centre of the rod in both cases.

\item A rigid cyclic quadrilateral $ABCD$ moves in its plane so as to return to its original position after turning through four right angles. Show that if $(A)$, etc., denote the areas of the curves described by $A$, etc., and if $S_1$, $S_2$, etc., denote the areas of the triangles $BCD$, $CDA$, etc., then
\[
S_1(A)+S_3(C)=S_2(B)+S_4(D).
\]
Find also the equation connecting the areas described by any three vertices with that described by the centre of the circumcircle of the triangle. \source{I. C. S., 1909.}

\item Two bars $OP$, $RPQ$, of lengths $OP=c$, $RPQ=b+a$, respectively turn round a fixed pin at $O$ and a joint at $P$. $dS_1$, $dS_2$ denote the polar elements of area about $O$ of the curves traced by $P$ and $Q$ respectively; prove that
\[
dS_2-dS_1=a\,d\zeta+a\left(\frac12a+b\right)d\theta-\frac12a\,dp,
\]
where $PQ=a$, $RP=b$, $p$ is the perpendicular from $O$ on $RPQ$, $d\zeta$ is the displacement of $R$ perpendicular to $RPQ$ and $\theta$ is the inclination of $RPQ$ to a fixed line $OA$. \source{Math. Trip., Pt. I., 1914.}

\end{enumerate}

\newpage

\begin{center}
    {\LARGE \bfseries \color{accent} Chapter XVI: Rectification (I). Elementary --- Exercises}\\
    \vspace{0.2cm}
    {\color{accent}\rule{0.55\linewidth}{0.8pt}}
\end{center}
\vspace{0.35cm}

\section*{Exercises --- Cartesian Rectification (Arts. 511--519)}

\begin{enumerate}[leftmargin=2em]

\item Find the length of the arc of the curve $y^2(2a-x)=x^3$, the cissoid of Diocles. \source{Huygens, 1625--1695.}

\item Find the curve for which the length of the arc measured from the origin varies as the square root of the ordinate.

\item The major axis of an ellipse is 1 foot in length, and its eccentricity is $\dfrac{1}{10}$. Prove that its circumference is 3.1337 feet nearly. \source{Trinity, 1883.}

\item Find the length of any arc of the curve
\[
x^{\frac23}-y^{\frac23}=a^{\frac23}.
\]

\item Show that in the ``catenary of equal strength,'' $y=a\log\sec\dfrac{x}{a}$,
\[
s=a\log\tan\left(\frac{x}{2a}+\frac{\pi}{4}\right),
\]
and that the intrinsic equation of the curve is $s=a\,\mathrm{gd}^{-1}\psi$.

\item Show that in the common catenary, or chainette, $y=c\cosh\dfrac{x}{c}$,
\[
s=\sqrt{y^2-c^2},\quad s=c\tan\psi,\quad s^2=c(\rho-c),\quad s=c\sinh\frac{x}{c}.
\]
The area bounded by the curve, the directrix, the $y$-axis and an ordinate is $A=cs$.

The centroid of the arc has coordinates
\[
\bar x=x-c\tan\frac{\psi}{2}\qquad \bar y=\frac12(y+x\cot\psi)
\]
\[
=x-(y-c)\frac{c}{s},\qquad=\frac12\left(y+\frac{cx}{s}\right).
\]

The centroid of the area bounded by the curve, the directrix, the $y$-axis and an ordinate is given by
\[
\bar x=x-(y-c)\frac{c}{s},\quad \bar y=\frac14\left(y+\frac{cx}{s}\right),
\]
and that both centroids lie on the ordinate through the intersection of the terminal tangents.

\item Show that the length of the curve $y=\log\coth\dfrac{x}{2}$ from the point $(x_1,y_1)$ to the point $(x_2,y_2)$ is $\log\dfrac{\sinh x_2}{\sinh x_1}$.

\item Show that in the epi- or hypo-cycloid
\[
x=(a+b)\cos\theta-b\cos\frac{a+b}{b}\theta,
\]
\[
y=(a+b)\sin\theta-b\sin\frac{a+b}{b}\theta.
\]
(i) $s=\dfrac{4b}{a}(a+b)\cos\dfrac{a\theta}{2b}$,   (ii) $s=\dfrac{4b}{a}(a+b)\cos\dfrac{a}{a+2b}\psi$,

$s$ being measured from the point where $\theta=\dfrac{\pi b}{a}$, i.e. a vertex.

\item For the four-cusped hypocycloid
\[
x^{\frac23}+y^{\frac23}=a^{\frac23},
\]
show (i) that $s=\dfrac{3a}{4}\cos2\psi$, $s$ being measured from a vertex;

(ii) the whole length of the curve is $6a$;

(iii) $s^3\propto x^2$, $s$ being measured from the cusp which lies on the $y$-axis.

\item In the tractrix
\[
x=\sqrt{c^2-y^2}+\frac{c}{2}\log\frac{c-\sqrt{c^2-y^2}}{c+\sqrt{c^2-y^2}},
\]
show that $s=c\log\dfrac{c}{y}$.

\item Show that the distance from the vertex of the centroid of a wire in the form of portion of a cycloid, of which the vertex is the middle point, is $\dfrac13$ of the greatest ordinate of the arc.

\item Show that the arc of a parabola of latus rectum $4a$ measured from the vertex, and the radius vector from the focus, are expressible in terms of a parameter $t$ in the respective forms
\[
\frac{s}{a}=\frac{t}{1-t^2}+\frac12\log\frac{1+t}{1-t},\quad \frac{r}{a}=\frac{1}{1-t^2}.
\]
\source{Math. Trip. Pt. II., 1915.}

Prove also that $s=\sqrt{r(r-a)}+a\tanh^{-1}\sqrt{\dfrac{r-a}{r}}$.

\end{enumerate}

\section*{Exercises --- Polar Rectification (Arts. 520--526)}

\begin{enumerate}[leftmargin=2em]

\item Find the length of any arc of the curve from the formula
\[
s=\int\sqrt{r^2+r'^2}\,d\theta
\]
for the following cases:

(i) $r=a\cos\theta$ (circle).   (ii) $r=ae^{m\theta}$ (equiang. spiral).

(iii) $r=a\sin^2\dfrac{\theta}{2}$ (cardioide).   (iv) $\dfrac{2a}{r}=1+\cos\theta$ (parabola).

(v) $r=a\dfrac{\sin^2\theta}{\cos\theta}$ (cissoid).   (vi) $r=\dfrac{3a}{2}\dfrac{\sin^2\theta}{\cos^3\theta}$ (semicub. parab.).

\item Show that the length of the arc of that part of the cardioide
\[
r=a(1+\cos\theta),
\]
which lies on the side of the line $4r=3a\sec\theta$ remote from the pole, is equal to $4a$. \source{Oxford.}

\item Show that the whole length of the limaçon $r=a\cos\theta+b$ is equal to that of an ellipse whose semiaxes are equal in length to the maximum and minimum radii vectores of the limaçon. Hence show how to divide the arc of the limaçon into four equal parts. \source{Colleges $a$, 1888.}

\item Prove that the length of the $n$th pedal of a loop of the curve
\[
r^m=a^m\sin m\theta
\]
is
\[
a(mn+1)\int_0^{\pi/m}\sin m\theta\,d\theta,\quad\text{where }m(k-n+1)=1.
\]

\item Show that the length of a loop of the curve
\[
3x^2y-y^3=(x^2+y^2)^3
\]
\[
=2\int_0^1\frac{d\xi}{\sqrt{1-\xi^6}}.
\]
\source{St. John's, 1881.}

\item Show that the rectification of the curve $r^n=a^n\sin n\theta$ is given by the integral
\[
s=a\int_0^\xi\frac{d\xi}{\sqrt{1-\xi^{2n}}}.
\]
\source{Math. Trip., 1896.}

\item Two radii vectores $OP$, $OQ$ of the curve
\[
r=2a\cos^3\left(\frac{\pi}{4}+\frac{\theta}{3}\right)
\]
are drawn equally inclined to the initial line; prove that the length of the intercepted arc is $a\alpha$, where $\alpha$ is the circular measure of the angle $POQ$. \source{Asparagus, Educ. Times.}

\item Show that the centroid of a wire bent into the form of a cardioide $r=a(1+\cos\theta)$, and with a line density $k\sec\dfrac{\theta}{2}$, $k$ being a constant, is on the axis of the cardioide at distance $\dfrac{a}{2}$ from the cusp.

\end{enumerate}

\section*{Exercises --- The Converse Problem (Arts. 527--529)}

\begin{enumerate}[leftmargin=2em]

\item Find the curves in which

(i) $s=a\sin^{-1}\dfrac{y}{a}$.   (ii) $s=\sqrt{y^2-c^2}$.

(iii) $s=\dfrac{r^2}{2a}$.   (iv) $y=ce^{-\frac{s}{c}}$.

(v) $s\propto r$.   (vi) $s\propto\sqrt{x}$.

(vii) $s=2\sqrt{2ar}$.

\item Show that the equation
\[
nx\frac{d^2s}{dx^2}+\frac{ds}{dx}=0
\]
leads to a cycloid or a four-cusped hypocycloid according as $n=2$ or $n=3$.

\end{enumerate}

\section*{Exercises --- Length of the Arc of an Evolute (Arts. 530--535)}

\begin{enumerate}[leftmargin=2em]

\item Show that in the parabola $y^2=4ax$, the length of the arc of the evolute intercepted within the parabola is
\[
4a(3\sqrt3-1).
\]

\item Find the whole length of the evolute of the cardioide
\[
r=a(1+\cos\theta).
\]

\item Show that the length of the evolute of the portion of the Folium of Descartes $x^3+y^3=3axy$, which corresponds to the loop, is $\dfrac{3a}{4}(4-\sqrt2)$.

\end{enumerate}

\section*{Problems (Arts. 536--565)}

\begin{enumerate}[leftmargin=2em]

\item Show that the whole length of the curve
\[
y^2(a^2-y^2)=8a^2x^2\quad\text{is}\quad\pi a\sqrt2.
\]
\source{Oxford I. P., 1890.}

\item Find the whole length of the loop of the curve
\[
3ay^2=x(x-a)^2.
\]
\source{Oxford I. P., 1889.}

\item Show that the arcs of an equiangular spiral, measured from the pole to the different points of its intersection with another equiangular spiral having the same pole but a different angle, will form a series in geometrical progression. \source{Trinity, 1884.}

\item Show that the length of an arc of the curve $y^n=x^{m+n}$ can be found in finite terms in the cases when $\dfrac{n}{2m}$ or $\dfrac{n}{2m}+\dfrac12$ is an integer.

\item Evaluate the expressions,
\[
\text{(i) }\int y\frac{dx}{ds}\,ds,\quad\text{(ii) }\int x\frac{dy}{ds}\,ds,\quad\text{(iii) }\int\left(\frac{x}{r^2}\frac{dy}{ds}-\frac{y}{r^2}\frac{dx}{ds}\right)ds,
\]
wherein the line-integrals are taken round the perimeter of a closed curve. \source{St. John's, 1890.}

\item If $s$ be the length of the curve
\[
r=a\tanh\frac{\theta}{2}
\]
between the origin and $\theta=2\pi$, and $A$ be the area between the same points, show that $A=a(s-a\pi)$. \source{Oxford I., 1888.}

\item Show that if the arc of the curve
\[
r=a\tanh^n\frac{\theta}{2n}\quad(n\text{ being integral}),
\]
measured from the origin, be called $s$, and if $A$ be the corresponding area swept out by the radius vector from the origin,
\[
A=\frac{a^2\theta}{2}-na^2\left[\left(\frac{r}{a}\right)^{\frac1n}+\frac13\left(\frac{r}{a}\right)^{\frac3n}+\dots+\frac{1}{2n-1}\left(\frac{r}{a}\right)^{\frac{2n-1}{n}}\right],
\]
\[
s+r=a\theta-2na\left[\left(\frac{r}{a}\right)^{\frac1n}+\frac13\left(\frac{r}{a}\right)^{\frac3n}+\dots+\frac{1}{n-2}\left(\frac{r}{a}\right)^{\frac{n-2}{n}}\right],\quad\text{if }n\text{ be odd and }>2,
\]
\[
s+r=2na\left[\log\cosh\frac{\theta}{2n}-\left\{\frac12\left(\frac{r}{a}\right)^{\frac2n}+\frac14\left(\frac{r}{a}\right)^{\frac4n}+\dots+\frac{1}{n-2}\left(\frac{r}{a}\right)^{\frac{n-2}{n}}\right\}\right],\quad\text{if }n\text{ be even,}
\]
the results for $r=a\tanh\dfrac{\theta}{2}$ giving $2A=a(s-r)=a(2s-a\theta)$.

\item Show that the length of an arc of the Cissoid of Diocles
\[
r=a\frac{\sin^2\theta}{\cos\theta}
\]
is $a\sqrt3(z-\tanh^{-1}z)$ taken between limits $\theta_1$ and $\theta_2$ where
\[
3z^2\cos\theta=1+3\cos^2\theta.
\]

\item Show that the intrinsic equation of the semicubical parabola
\[
3ay^2=2x^3\quad\text{is}\quad9s=4a(\sec^3\psi-1).
\]

\item In a certain curve
\[
x=e^\theta\sin\theta,\quad y=e^\theta\cos\theta.
\]
Show that $s=e^\theta\sqrt2+C$.

Also that for the curve
\[
x=e^{a\theta}\sin b\theta,\quad y=e^{a\theta}\cos b\theta,\quad s=e^{a\theta}\sqrt{a^2+b^2}+C.
\]
Name these curves.

\item Show that the length of an arc of the curve
\[
x\sin\theta+y\cos\theta=f'(\theta),
\]
\[
x\cos\theta-y\sin\theta=f''(\theta),
\]
is given by $s=f(\theta)+f''(\theta)+C$.

\item Trace the curve $y^2=\dfrac{x}{3a}(a-x)^2$, and find the length of that part of the evolute which corresponds to the loop. \source{St. John's, 1881 and 1891.}

\item Show that the curve whose pedal equation is $p^2=r^2-a^2$ has for its intrinsic equation $s=a\dfrac{\psi^2}{2}$.

What curve is this?

\item The coordinates of a point on a plane curve are given by the relations
\[
x=a[(1-\theta^2)\cos\theta+2\theta\sin\theta-1],
\]
\[
y=a[(1-\theta^2)\sin\theta-2\theta\cos\theta];
\]
prove that $3sa^{\frac13}=(2a+\rho)(\rho-a)^{\frac13}$,

$\rho$ being the radius of curvature at the point and $s$ the arcual distance from the origin. \source{Oxford II. P., 1888.}

\item The evolute of a parabola whose vertex is $A$ meets the axis in $C$, and the parabola in $Q$. Find the perimeter of the figure bounded by $AC$, the parabolic arc $AQ$, and the arc of the evolute $CQ$. \source{Oxford I. P., 1889.}

\item Prove that the length of the first negative pedal, taken with respect to the origin, of the loop of the folium of Descartes
\[
x^3+y^3-3axy=0
\]
is equal to $6a-a\{\pi-\sqrt2\log(\sqrt2+1)\}$.

\item Find the length of the arc between two consecutive cusps of the curve
\[
(c^2-a^2)p^2=c^2(r^2-a^2).
\]
\source{Oxford I. P., 1889.}

\item Show that the length of the arc of the hyperbola $xy=a^2$ between the limits $x=b$ and $x=c$ is equal to the arc of the curve $p^2(a^4+r^4)=a^4r^2$, between the limits $r=b$, $r=c$. \source{Oxford I. P., 1888.}

\item By means of the formula $s=\displaystyle\int\frac{r\,dr}{\sqrt{r^2-p^2}}$, find the length of the curve $r=a\sin^2\dfrac{\theta}{2}$. \source{Colleges $a$, 1887.}

\item If $s$ be the arc of an ellipse $\dfrac{x^2}{a^2}+\dfrac{y^2}{b^2}=1$ measured from the end of the major axis to a point whose eccentric angle is $\phi$, prove that
\[
s+ae^2\cos\phi\sin\theta=\int_0^\theta\sqrt{a^2\cos^2\theta+b^2\sin^2\theta}\,d\theta,
\]
where $\theta=\tan^{-1}\left(\dfrac{a}{b}\tan\phi\right)$. \source{Colleges $a$, 1883.}

\item Show that the circumference of an ellipse can be expressed either as
\[
4a\int_0^{\pi/2}(1-e^2\sin^2\theta)^{\frac12}\,d\theta
\]
or as
\[
4a(1-e^2)\int_0^{\pi/2}(1-e^2\sin^2\theta)^{-\frac32}\,d\theta,
\]
where $a$ is the semi-major axis, $e$ the eccentricity. \source{Trinity, 1887.}

\item Show that the three-cusped hypocycloid has equations of the forms

(i) $p=b\cos3\psi$,

(ii) $r^4+8br^3\cos3\theta+18b^2r^2=27b^4$.

Show that the length of an arc of the inverse of this curve with respect to the centre is proportional to
\[
\tan^{-1}(2\sqrt2\sin3\psi).
\]
\source{St. John's, 1887.}

\item Prove that the intrinsic equation of the curve
\[
r^{\frac15}=a^{\frac15}\cos\frac15\theta\quad\text{is}\quad\frac{s}{a}=\frac{5}{16}\psi+\frac54\sin\frac{\psi}{3}+\frac{5}{32}\sin\frac{2\psi}{3},
\]
where $s$ and $\psi$ are measured from the point $a,0$, and the tangent at that point. \source{St. John's, 1889.}

\item A circle of perimeter $S$ and area $A$ rolls externally in its plane entirely round an oval curve of perimeter $S$ and area $B$. Prove that its centre describes an oval of perimeter $2S$ and area $3A+B$. \source{Oxford I. P., 1918.}

\item Find the centroid of a sector bounded by two radii vectores, and an arc of the curve whose polar equation is
\[
r^2=a^2(1-\sin2\theta)(1+\sin2\theta)^{-1},
\]
and show that an arc of this curve is expressible as
\[
\frac{5a}{2}\int\frac{\cos^2\chi\,d\chi}{(1+\sqrt5\sin\chi)\sqrt{1-5\sin^2\chi}}.
\]
\source{Matz, Educ. Times.}

\item A rod moves always to pass through a fixed point and have one extremity on a straight line distant $h$ from the point. Show that the arc of the curve traced out by its centre of instantaneous rotation, as the rod moves from the perpendicular position to one inclined at $45^\circ$ to the line, is
\[
\frac14\{\log(\sqrt5+2)+2\sqrt5\}h.
\]
\source{Math. Tripos, 1883.}

\item On the tangent at any point $P$ of a curve, $PT$ is taken equal to the radius vector of $P$; show how to find the length of any arc of the locus of $T$. For example, take the equiangular spiral and verify the result geometrically. \source{St. John's, 1884.}

\item Find the arc of the curve enveloped by the line
\[
x\cos\phi+y\sin\phi=(a\cos^2\phi+b\sin^2\phi)^{-3}
\]
between the points corresponding to $\phi=0$, $\phi=\dfrac{\pi}{2}$. \source{St. John's, 1891.}

\item Find the whole area of the curve
\[
x=a\sin\theta-b\sin2\theta,
\]
\[
y=a\cos\theta-b\cos2\theta,
\]
and show that the whole length of its perimeter is equal to that of an ellipse whose semiaxes are $a+2b$, $a-2b$. \source{Colleges $a$, 1885.}

\item Prove that if $s$ be the arc of the curve
\[
r=a\sec\alpha,
\]
\[
\theta=\tan\alpha-\alpha,
\]
where $\alpha$ is a variable parameter, measured from the initial line to a point $P$ on the curve, and if $A$ be the area bounded by the curve, the initial line, and the radius vector to $P$, then
\[
9A^2=2as^3.
\]
Find the area swept out in any portion of its progress by the intercept of the tangent to the curve between the curve and the first positive pedal with regard to the origin. \source{Trinity, 1890.}

\item If a curve be given by $r^{2n}=\sin^2\phi+m^2\cos^2\phi$, where $r$ is the radius vector and $\phi$ the angle it makes with the tangent, show that
\[
m\tan\frac1m(\phi-n\theta)=\tan\phi,
\]
$\theta$ being the angle the radius vector makes with the initial line (which is to be appropriately chosen).

Obtain also a formula for the rectification of the curve. (The result is not obtainable in finite terms.) \source{I. C. S., 1898.}

\item Consider the nature of the curves

(i) $s\psi^2=a$,  (ii) $\dfrac{\psi}{2\pi}=\sin\dfrac{s}{a}$,  (iii) $s=l\sin m\psi$,

when $m<1$ and when $m>1$.

\item Given a closed oval of continuous curvature without any singularities: a series of parallel curves is drawn. Prove that if $A$ denote the area of any one of them and $l$ its perimeter, then
\[
4\pi A-l^2
\]
is the same for all. \source{I. C. S., 1895.}

\item In the equation of the curve $r=a+\epsilon u$, $a$ and $\epsilon$ are constants, the latter being small; and $u$ is a function of $\theta$ finite for all values of $\theta$ and periodic, with a period $2\pi$. Show that if $A$ denote the area of the curve, then its length is $2\sqrt{\pi A}$ accurately as far as small quantities of the first order inclusive. \source{I. C. S., 1896.}

\item The area of an ellipse differs from that of its auxiliary circle by 10 per cent. of the area of the latter. Show that the perimeter of the ellipse differs from that of the auxiliary circle by 4.93 per cent. approximately of the perimeter of the latter. \source{I. C. S., 1910.}

\item Assuming that for the catenary formed by a hanging elastic wire
\[
\frac{x}{c}=u+k\sinh u,\quad \frac{y}{c}=\cosh u+\frac12k\cosh^2u,
\]
prove that
\[
\frac{s}{c}=\frac12ku+\sinh u+\frac14k\sinh2u,
\]
reducing to the common catenary when $k=0$ and approximating to a parabola when $k$ is large. \source{B.A. Hon. Lond., 1899.}

\item In the cycloid $y=\dfrac{s^2}{8a}$. Show that the only curve for which both $x$ and $y$ are finite integral functions of $s$ is a straight line. \source{Oxf. I. P., 1913.}

\item Find the Cartesian equation (choosing convenient axes of coordinates) of the curve in which
\[
\rho^2/a^2=(d\rho/ds)^2+1.
\]
\source{Oxf. I. P., 1917.}

\item Find the intrinsic equation of the curve $27ay^2=4x^3$. Prove that the involutes of the curve $27ay^2=4x^3$ are given by the equations
\[
x=a\tan^2\psi+c\cos\psi-2a,
\]
\[
y=-2a\tan\psi+c\sin\psi,
\]
$c$ being an arbitrary constant.

What happens when $c=0$? \source{Oxf. I. P., 1915.}

\item Show that the length of a quadrant of the curve $x^{\frac23}+y^{\frac23}=a^{\frac23}$ is equal to $\dfrac{3a}{2}$, and find the length of one quadrant of the curve
\[
\left(\frac{x}{a}\right)^{\frac23}+\left(\frac{y}{b}\right)^{\frac23}=1.
\]
\source{Math. Tripos, Part I., 1910.}

\item Trace the form of the curve
\[
x(1+t^2)=1-t^2,\quad y(1+t^2)=2t,
\]
as $t$ increases from $-\infty$ to $\infty$, and show that its area is $\pi$.

Find also the length of any arc of the curve in terms of $t$. \source{Math. Trip., Part I., 1910.}

\item Show that the length of the arc of the parabola $y^2=4ax$ which is intercepted between the point of intersection of the parabola and $3y=8x$ is
\[
a\left(\log2+1\tfrac{5}{16}\right).
\]
\source{Math. Trip. I., 1908.}

\item Prove that the perimeter of an ellipse of small eccentricity $e$ and semiaxes $a$, $b$ is equal to
\[
\frac12\pi\{3(a+b)-2\sqrt{ab}\},
\]
neglecting $e^7$ and higher powers. \source{Math. Trip. I., 1917.}

\item Prove that the length of an ellipse may be expressed by
\[
\iint\frac{dS}{\rho}
\]
taken over the area, where $dS$ is an element of the area of the ellipse and $\rho$ the radius of curvature of the similar, similarly situated and concentric ellipse passing through the element $dS$. \source{Colleges, 1892.}

\item Find the intrinsic equations of a circle, a catenary, and a cycloid, and trace the curves $s=a\phi^3$, $s\phi=a$ and $s^2=a^2\phi$.

At any point $P$ of a cycloid the tangent is produced to a length $PT$ equal to the arc measured from the vertex, and at $T$ a perpendicular is drawn equal to the radius of curvature at $P$. Prove that the locus of the extremity of this perpendicular is the same cycloid moved parallel to its axis through a distance equal to twice the diameter of the generating circle. \source{St. John's College, 1882.}

\end{enumerate}

\newpage

\begin{center}
    {\LARGE \bfseries \color{accent} Chapter XVII: Rectification (II). Central Conic, Limaçon, Lemniscate, Trochoids, etc. Application of Elliptic Functions --- Exercises}\\
    \vspace{0.2cm}
    {\color{accent}\rule{0.55\linewidth}{0.8pt}}
\end{center}
\vspace{0.35cm}

\section*{Exercises --- Fagnano's Theorem (Arts. 580--581)}

\begin{enumerate}[leftmargin=2em]

\item Show that if coaxial ellipses be drawn with a given centre such that the areas enclosed between them and their respective director circles is constant, the locus of the Fagnano points is a circle of the same area.

\item Show that the locus of the Fagnano points for similar and similarly situated concentric ellipses is a pair of straight lines.

\item Show that the locus of the Fagnano points which lie on confocal ellipses is
\[
(x^{\frac23}+y^{\frac23})^2(x^{\frac23}-y^{\frac23})=c^2,
\]
$2c$ being the distance between the foci.

\item Show that if $F$ be the Fagnano point on an ellipse of semiaxes $OA=a$, $OB=b$,
\[
2\,\text{arc }BF=aE_1+a-b,
\]
\[
2\,\text{arc }AF=aE_1-a+b,
\]
where $E_1$ is the complete elliptic integral of the second kind
\[
\int_0^{\pi/2}\sqrt{1-e^2\sin^2\phi}\,d\phi.
\]

\item Show that the central perpendicular upon the tangent at a Fagnano point is a geometric mean between the semiaxes, and equal to the semi-diameter conjugate to the radius to the Fagnano point. Further, that the radius of curvature at this point is also equal to the perpendicular, and that the normals at the corresponding point on the evolute pass through the centre. Finally, that the arc of the evolute is at such a point divided in the ratio
\[
b^{\frac32}:a^{\frac32}.
\]

\item Show that if a straight rod $LM$ of length $a+b$ slides with its ends on two axes $Ox$, $Oy$ at right angles and carries a point $F$ whose distances from $L$ and $M$ are respectively $a$ and $b$, which thus describes an ellipse, then at the instant when $LM$ is tangential to the path of $F$, $F$ is a Fagnano point on the described ellipse, and the circle on $LM$ for diameter passes through the point on the normal at $F$ where that normal touches the evolute.

\item Show that the tangents at the points $P(x_1,y_1)$, $Q(x_2,y_2)$ on an ellipse $\dfrac{x^2}{a^2}+\dfrac{y^2}{b^2}=1$, which are related to each other so that $\dfrac{x_1x_2}{a^2}=\dfrac{y_1y_2}{b^2}$, intersect on a confocal hyperbola which passes through the Fagnano points.

[Many properties of these points will be found in Greenhill's \textit{Elliptic Functions}, pages 182, 183.]

\end{enumerate}

\section*{Exercises --- Central Conics (Arts. 583--589)}

\begin{enumerate}[leftmargin=2em]

\item In the hyperbola $\dfrac{x^2}{a^2}-\dfrac{y^2}{b^2}=1$, put $a=b\tan\alpha$, $\Delta=\sqrt{1-\sin^2\alpha\sin^2\phi}$, and show that we may take $x=b\tan\alpha\sec\phi\,\Delta$, $y=b\cos\alpha\tan\phi$, and that
\[
\frac{ds}{d\phi}=\frac{b\cos\alpha}{\Delta\cos^2\phi},\quad t=\frac{b}{\cos\alpha}\Delta\tan\phi,
\]
and
\[
s=b\sec\alpha\tan\phi\,\Delta+b\cos\alpha\,F(\phi,\sin\alpha)-b\sec\alpha\,E(\phi,\sin\alpha).
\]

\item From the polar equation $r^2=a^2\sec2\theta$ deduce the rectification of the rectangular hyperbola, viz.
\[
s=a\sqrt2[\Delta\tan\omega+\tfrac12F-E].
\]

\item If $PQ$ be a chord of one branch of a hyperbola, touching a confocal ellipse at $F$, and the confocal cutting that branch of the hyperbola at $A$ and $B$, and if $PR$, $QS$ be the other tangents from $P$ and $Q$ to the ellipse, show that the elliptic arcs $AR$, $BS$ exceed the elliptic arc $AFB$ by the excess of the tangents $PR$, $QS$ over the chord $PQ$, i.e. that
\[
\text{arc }AR+\text{arc }BS-\text{arc }AFB
\]
is rectifiable in terms of known lines.

In particular, examine what happens:

(1) When $F$ is the vertex of the confocal ellipse.

(2) When $F$ is at $B$.

(3) When $PR$ and $QS$ are at right angles to $PQ$ and $F$ the vertex of the ellipse.

\end{enumerate}

\section*{Exercises --- The Lemniscate (Art. 592)}

\begin{enumerate}[leftmargin=2em]

\item Find the length of the arc of a lemniscate $r^2=a^2\cos2\theta$ from $\theta=0$ to $\theta=\dfrac{\pi}{6}$.

\item Find the area of the curve $y^2=\dfrac{1}{1-x^4}$ for the portion in the first quadrant. What connection is there between this problem and the evaluation of the perimeter of the lemniscate?

\item Draw a careful polar graph of the lemniscate $r^2=25\cos2\theta$, taking one inch as unit of length, and deduce a Cartesian graph of
\[
y=5\,\mathrm{cn}\,\frac{x\sqrt2}{5}\quad\left(\text{mod. }\frac{1}{\sqrt2}\right).
\]

\item Show that the difference between the lengths of the asymptote and the infinite arc of the hyperbola $x^2/a^2-y^2/b^2=1$ in the first quadrant is
\[
t-s=\frac{\pi a}{2}\left[\frac12\cdot\frac{1}{e}+\frac{1\cdot1^2}{2^2\cdot4}\cdot\frac{1}{e^3}+\frac{1\cdot1^2\cdot3^2}{2^2\cdot4^2\cdot6}\cdot\frac{1}{e^5}+\frac{1\cdot1^2\cdot3^2\cdot5^2}{2^2\cdot4^2\cdot6^2\cdot8}\cdot\frac{1}{e^7}+\dots\right].
\]

\end{enumerate}

\section*{Miscellaneous Problems}

\begin{enumerate}[leftmargin=2em]

\item Prove that the three equations
\[
x=c\log\sec\psi,\quad y=c(\tan\psi-\psi),\quad s=c(\sec\psi-1),
\]
represent one and the same curve. \source{I. C. S., 1893.}

\item Find the area of the curve
\[
r_1r_2=b^2,
\]
considering all cases which may arise.

\item Prove that the value of the integral
\[
\int pxy\left(\frac{x}{a^3}+\frac{y}{b^3}\right)^2ds,
\]
taken round the ellipse $x^2/a^2+y^2/b^2=1$, is $\dfrac{\pi}{2}$, $p$ denoting the central perpendicular on the tangent at $(x,y)$ and $ds$ an element of arc. \source{I. C. S., 1912.}

\item If the point $x, y$ lies on the curve
\[
y^2=x^2+2px+q,
\]
prove that
\[
\frac{dx}{y}=\frac{dy}{x+p}=\frac{dx+dy}{x+y+p},
\]
and hence obtain the integral of $\dfrac{dx}{\sqrt{x^2+2px+q}}$.

If, however, the point $(x,y)$ lie on the circle $x^2+y^2=a^2$, show that the corresponding relation is
\[
\frac{dx}{y}=-\frac{dy}{x}=\pm\frac{ds}{a},
\]
where $s$ is the length of the arc measured to the point $(x,y)$.

Deduce the known formula for the integral of $\dfrac{dx}{\sqrt{a^2-x^2}}$. \source{I. C. S., 1908.}

\item Show that if $\delta$ stands for $\dfrac{d}{dx}$,
\[
\delta^{-n}(uv)=v\delta^{-n}u-n\frac{dv}{dx}\delta^{-n-1}u+\frac{n(n+1)}{1.2}\frac{d^2v}{dx^2}\delta^{-n-2}u-\dots.
\]

\item If $\dfrac1R\dfrac{dR}{dx}=\dfrac{fx+g}{ax^2+bx+c}$,

and if $b^2-4ac$ be positive and the roots of $ax^2+bx+c=0$ be $\lambda$ and $\mu$, prove that $R=(x-\lambda)^p(x-\mu)^q$, where
\[
\frac{p}{q}=\frac{f}{2a}\pm\frac{2ag-bf}{2a^2(\lambda-\mu)}.
\]
And if $a=-1$, $b=0$, $c=1$,
\[
R=(1-x^2)^{-\frac{f}{2}}\left(\frac{1+x}{1-x}\right)^{\frac{g}{2}}.
\]
If $b^2-4ac$ be negative,
\[
R=(ax^2+bx+c)^{\frac{f}{2a}}e^{\frac{2ag-bf}{2a\sqrt{4ac-b^2}}\tan^{-1}\frac{2ax+b}{\sqrt{4ac-b^2}}}.
\]
If $b^2-4ac=0$,
\[
R=\left(\frac{2ax+b}{2a}\right)^{\frac{f}{a}}e^{-\frac{2ag-bf}{a(2ax+b)}}.
\]
\source{E. J. Routh, Proc. L.M.S., vol. xvi., p. 250.}

\item Show that
\[
\iiint\dots\int(x^2-1)^{-k-1}\left(\frac{x-1}{x+1}\right)^ldx\,dx\dots dx,
\]
there being $2k+1$ integrations, $2k+1$ being an integer, though $k$ may be a fraction, is equal to
\[
\frac{1}{2^{2k+1}M}(x^2-1)^k\left(\frac{x-1}{x+1}\right)^l,
\]
where $M=(k+l)(k+l-1)\dots$ to $2k+1$ factors. \source{Cf. Routh, Proc. L.M.S., vol. xvi., p. 249.}

\item $ABC$ is a triangle with the corner $A$ fixed and with sides $AC$, $CB$ respectively $\sqrt n$ and $\sqrt{n+1}$, given lengths.

The side $AB$ ($=r$) makes an angle $\theta=nA-(n+1)B$ with a fixed straight line $AX$.

Show (1) that the path of $B$ is rectifiable by the formula
\[
s=\sqrt n\int_0^A\frac{d\theta}{\sqrt{1-\dfrac{n}{n+1}\sin^2\theta}}=\sqrt n\,\mathrm{am}^{-1}A,\quad\text{mod. }\sqrt{\frac{n}{n+1}}.
\]

(2) When $n=1$ the rectification is the same as that of a Bernoulli's Lemniscate.

(3) The inclination of the normal to the radius vector is $A+B$.

(4) The area of the triangle is equal to the area of a sector of the curve starting from the axis $AX$. \source{M. Serret's Problem, Calc. Int., p. 269.}

\item $C$ is a point of maximum curvature on the Limaçon
\[
r=a\cos\theta+b,\quad b>a;
\]
$A$ and $A'$ are the two vertices. Prove that the difference between the arcs $AC$, $A'C$ is $4a$. \source{St. John's, 1891.}

\item If $y=x^3-3a^2x$, prove that
\[
\frac{dx}{\sqrt{x^2-4a^2}}=\frac{dy}{3\sqrt{y^2-4a^6}},
\]
and by integration express $x$ explicitly in terms of $y$. \source{Oxford I. P., 1916.}

Apply this method to solve the cubic
\[
x^3-3x^2-45x-473=0.
\]

\item Prove that
\[
\int_0^{\pi/2}e^{-\sin2x}\cos x\,dx=e^{-1}\left\{1+\frac{1}{3.1!}+\frac{1}{5.2!}+\frac{1}{7.3!}+\dots\right\}.
\]
\source{Oxford I. P., 1916.}

\item Prove that if $n$ be an odd positive integer greater than 3,
\[
n\int_0^{\pi/2}\sin^nx\,dx=(2-\sqrt3)\frac{n-1}{n-2}\cdot\frac{n-3}{n-4}\dots\frac43
\]
\[
-\frac{\sqrt3}{8}\left[\frac{1}{2^{n-3}}+\frac{n-1}{n-2}\frac{1}{2^{n-5}}+\dots+\frac{n-1}{n-2}\cdot\frac{n-3}{n-4}\cdots\frac43\right].
\]
\source{Oxford I. P., 1916.}

\item The parameters $t_1$, $t_2$ of two points $A$, $B$ of the unicursal curve
\[
x/(1-t^2)=y/(t-t^3)=a/(1+t^2)
\]
are equal to $\tan\alpha$, $\tan\beta$, where
\[
-\tfrac12\pi<\alpha<-\tfrac14\pi,\quad \tfrac14\pi<\beta<\tfrac12\pi.
\]
Prove that the area of the curvilinear triangle $AOB$, where $O$ is the double point, is
\[
a^2\left[2-\frac{\pi}{2}+\beta-\alpha-\sec\beta\sec\alpha\sin(\beta-\alpha)+\frac12\tan\beta\tan\alpha\sin2(\beta-\alpha)\right].
\]
\source{Oxford I. P., 1916.}

\item If $n$ be a positive integer, show that
\[
\int_0^\pi x\sin^2nx\,\csc^2x\,dx=\tfrac12n\pi^2.
\]
\source{Oxford I. P., 1912.}

\item By assuming
\[
\int\frac{x^3}{\sqrt{1+x^4}}\,dx\quad\text{to be of the form}\quad\frac{a+bx+cx^2+dx^3+ex^4}{\sqrt{1+x^4}},
\]
obtain the integration by differentiation and equating coefficients, also obtain the result directly by putting $x^4=z$.

\item Given a rational integral relation between $x$ and $y$ of the form
\[
y^n+A_1y^{n-1}+A_2y^{n-2}+\dots+A_n=0=F(x,y),\text{ say,}
\]
where $A_1, A_2,\dots A_n$ are rational functions of $x$, prove that when $\displaystyle\int y\,dx$ can be expressed algebraically in terms of $x$, then
\[
\int y\,dx=B_0+B_1y+B_2y^2+B_3y^3+\dots+B_{n-1}y^{n-1},
\]
where $B_0, B_1, B_2\dots$ are rational functions of $x$. \source{Abel.}

\item Assuming $X$ to be a rational function of $x$, and $y^m=X$, and that $\displaystyle\int y\,dx$ is integrable in algebraic form and expressible as
\[
\int y\,dx=P_0+P_1y+P_2y^2+\dots+P_{m-1}y^{m-1},
\]
where $P_0, P_1,\dots P_{m-1}$ are rational functions of $x$, show that
\[
P_0=P_2=P_3=\dots=P_{m-1}=0,
\]
that is the integration must contain one term only, and that $\dfrac1y\displaystyle\int y\,dx$ is a rational algebraic function of $x$. \source{Liouville.}

\item If $M$ and $T$ be two rational polynomials, then, provided $\displaystyle\int\dfrac{M}{\sqrt[m]{T}}dx$ can be integrated in algebraic form at all, the form of the integral is $\dfrac{\theta}{\sqrt[m]{T}}$, where $\theta$ is a function of $x$.

Show also that

(1) $\theta$ is a rational function of $x$.

(2) That $MT=T\dfrac{d\theta}{dx}-\dfrac1m\theta\dfrac{dT}{dx}$.

(3) That $\theta$ is an integral polynomial expression and not of such form as $\dfrac{U}{V}$, where $U$ and $V$ are complete polynomials, i.e. not such that $V$ contains $x$.

(4) That the degree of the polynomial $\theta$ is greater by unity than the degree of $M$.

Use these facts to show that $\displaystyle\int\dfrac{x^2\,dx}{\sqrt{1+x^4}}$ is not expressible algebraically. \source{Bertrand, C. I., p. 94.}

\item $P, Q, R$ being any rational algebraic polynomials, and assuming that when $\displaystyle\int\dfrac{P}{Q}\dfrac{dx}{\sqrt R}$ is integrable by means of the ordinary elementary functions, the integral must be of the form
\[
\int\frac{P}{Q}\frac{dx}{\sqrt R}=\eta+\frac{\theta}{\sqrt R}+A\log(\alpha_1+\beta_1\sqrt R)+B\log(\alpha_2+\beta_2\sqrt R)+\dots,
\]
where $\eta,\theta,\alpha_1,\beta_1,\alpha_2,\dots$ etc., are rational functions of $x$, a result established by Abel, show that when the integration can be reduced to one term, the general type of the result is either algebraic of form $\theta/\sqrt R$ or may be written as
\[
\int\frac{P}{Q}\frac{dx}{\sqrt R}=A\tanh^{-1}\frac{\beta\sqrt R}{\alpha}.
\]
In the latter case show that

(1) $\alpha^2-\beta^2R=Q$.

(2) $\alpha\beta R\left(\dfrac{\beta'}{\beta}-\dfrac{\alpha'}{\alpha}+\dfrac12\dfrac{R'}{R}\right)=\dfrac{P}{A}$.

(3) $2\alpha'Q-Q'\alpha=\dfrac{2P}{A}\beta$,

where accents denote differentiation with regard to $x$.

\item Show that
\[
\int\frac{x^6(11x^2+37x+28)}{x^{14}-x^3-4x^2-5x-2}\frac{dx}{\sqrt{x+2}}=-2\tanh^{-1}\frac{(x+1)\sqrt{x+2}}{x^7}.
\]

\item Prove that
\[
\int\frac{x\,dx}{\sqrt{(x^2+2x-5)(x^2+4x-8)}}=\frac12\tanh^{-1}\frac{x+4}{x+5}\sqrt{\frac{x^2+4x-8}{x^2+2x-5}}.
\]

\item Prove that

(1) $\displaystyle\int\frac{5x^2+3x+1}{2x+1}\frac{dx}{\sqrt{x^4+2x^2+2x+1}}=\tanh^{-1}\frac{x\sqrt{x^4+2x^2+2x+1}}{x^3+x+1}$.

(2) $\displaystyle\int\frac{3\tan^2\theta+2}{2\tan^4\theta+5\tan^2\theta+4}\sec\theta\,d\theta=\tan^{-1}\frac{\sin\theta}{1+\cos^2\theta}$.

(3) $\displaystyle\int\frac{\sinh2x\,dx}{\sqrt{\cosh^2x+1}\sqrt{\cosh^2x+2}}=2\tanh^{-1}\sqrt{\frac{\cosh^2x+1}{\cosh^2x+2}}$.

\item Show that

(i) $\dfrac{2n+1}{2}a^2\displaystyle\int\sqrt{\dfrac{a^{2n-1}+x^{2n-1}}{a^{2n+1}+x^{2n+1}}}\dfrac{dx}{a^2-x^2}-\dfrac{2n-1}{2}\displaystyle\int\sqrt{\dfrac{a^{2n+1}+x^{2n+1}}{a^{2n-1}+x^{2n-1}}}\dfrac{dx}{a^2-x^2}$
\[
=\tanh^{-1}x\sqrt{\frac{x^{2n-1}+a^{2n-1}}{x^{2n+1}+a^{2n+1}}}.
\]

(ii) $\dfrac{2n+1}{2}\displaystyle\int\sqrt{\dfrac{1+\sin^{2n-1}\theta}{1+\sin^{2n+1}\theta}}\sec\theta\,d\theta-\dfrac{2n-1}{2}\displaystyle\int\sqrt{\dfrac{1+\sin^{2n+1}\theta}{1+\sin^{2n-1}\theta}}\sec\theta\,d\theta=\tanh^{-1}\left(\sin\theta\sqrt{\dfrac{1+\sin^{2n-1}\theta}{1+\sin^{2n+1}\theta}}\right)$.

\item Integrate the following:

(i) $\displaystyle\int\frac{(2x+1)dx}{\sqrt{x^4+2x^3-3x^2-4x+3}}$.   (ii) $\displaystyle\int\frac{x\,dx}{\sqrt{x^4+2x^3-3x^2-ax+a}}$. \source{Abel.}

(iii) $\displaystyle\int\frac{5x^2+15x+12}{5x^2+15x+9}\frac{dx}{\sqrt{(x+1)(x+2)}}$.

(iv) $\displaystyle\int\frac{(1+8x)\,dx}{\sqrt{1+6x+4x^2}\sqrt{1-2x+4x^2}}$.

(v) $\displaystyle\int\frac{(2x+a)\,dx}{\sqrt{x^4+2ax^3+3a^2x^2+2a^3x-a^4}}$.

(vi) $\displaystyle\int\frac{3x^4-2x^3+1}{2x^3-x^2+1}\frac{dx}{\sqrt{x^4+1}}$.   (vii) $\displaystyle\int\frac{3x+1}{1-x-2x^2-x^3}\frac{dx}{\sqrt x}$.

(viii) $\displaystyle\int\frac{3x^4+1}{1-x^2-x^6}\frac{dx}{\sqrt{x^4+1}}$.   (ix) $\displaystyle\int\frac{x^2-(a+b)x-ab}{\sqrt{x^2+a^2}\sqrt{x^2+b^2}}dx$.

(x) $\displaystyle\int\frac{1+x-x^2}{1-x+x^2+x^3}\frac{dx}{\sqrt{x^2-1}}$.   (xi) $\displaystyle\int\frac{1+x^4}{1+x^2-x^4}\frac{dx}{\sqrt{x^4-1}}$.

(xii) $\displaystyle\int\frac{1+3x^4+2x^5}{1+2x-x^6}\frac{dx}{\sqrt{1+x^4}}$.

\item Show that the whole perimeter and area of a single loop of the curve $r=2a\cos n\theta$ ($n>1$) are respectively equal to the whole perimeter and area of the ellipse $x^2+n^2y^2=a^2$. \source{Oxf. I. P., 1911.}

\item If an element $ds$ of a curve lie at distance $r$ from the origin, and subtends an angle $d\theta$ there, it is known that unit electric current flowing along $ds$ produces a magnetic force at the origin at right angles to the plane of the curve proportional to $\dfrac{d\theta}{r}$.

Show that if unit current flows through a thin endless wire of given length in the form of an ellipse, the magnetic force due to the current at the centre of the ellipse is inversely proportional to the area of the ellipse. \source{Oxford II. P., 1913.}

\item A current of electricity is flowing round a fine wire $ABCD\dots KA$ bent into a plane polygon. $O$ is any point within the polygon, and perpendiculars $OP, OQ, OR,\dots$ are drawn to the sides $KA, AB, BC$, etc., respectively, and again perpendiculars whose lengths are $\alpha,\beta,\gamma,\dots$ from $O$ upon the sides $PQ, QR, RS,\dots$ of the inscribed polygon $PQRS\dots$. Show that the magnetic force on unit particle situated at $O$ is
\[
\iota\sum\frac{\sin A}{\alpha},
\]
where $\iota$ is the current strength.

\item Show that the perimeter of an ellipse of axes $2a$, $2b$ and small eccentricity $e$ is approximately equal to the perimeter of a circle of diameter $a+b$, with an error which is only about 0.0025 per cent. when $e$ is as great as 0.2. \source{Math. Trip. Part II., 1913.}

\end{enumerate}

\newpage

\begin{center}
    {\LARGE \bfseries \color{accent} Chapter XVIII: Rectification (III). Miscellaneous Theorems --- Exercises}\\
    \vspace{0.2cm}
    {\color{accent}\rule{0.55\linewidth}{0.8pt}}
\end{center}
\vspace{0.35cm}

\section*{Miscellaneous Problems}

\begin{enumerate}[leftmargin=2em]

\item Show that any point on the Lemniscate $r^2=a^2\cos2\theta$ may be represented by
\[
x=a\frac{z+z^3}{1+z^4},\quad y=a\frac{z-z^3}{1+z^4},
\]
and hence obtain the rectification of the curve. \source{Serret.}

Show that the integral obtained for $s$ reduces to the standard Legendrian form by the further substitution
\[
\cos\phi=\frac{z\sqrt2}{\sqrt{1+z^4}}.
\]
\source{Cayley, Ell. Functions, Art. 63.}

\item By the transformation $\dfrac{z-\iota}{z+\iota}=\dfrac{a-\iota}{a+\iota}u$,

show that the equation
\[
x+\iota y=Ce^{\iota\omega}\int\frac{(z-a)^{n+1}(z-\iota)^m}{(z-\alpha)^{n+1}(z+\iota)^{m+2}}\,dz
\]
takes the form
\[
x+\iota y=A\int\frac{u^m(u-1)^{n+1}}{(u-\xi)^{n+1}}\,du,
\]
where
\[
\xi=\frac{(a+\iota)(\alpha-\iota)}{(\alpha-\iota)(a+\iota)},\quad A=\frac{C}{2\iota}e^{\iota\omega}\left(\frac{a+\iota}{\alpha+\iota}\right)^{n+1}\left(\frac{a-\iota}{a+\iota}\right)^{m+1}
\]

Hence show that the condition that $x+\iota y$ should be purely algebraic is
\[
\frac{d^n}{d\xi^n}\xi^m(\xi-1)^{n+1}=0,
\]
$a$ and $\alpha$ being supposed conjugate, and $m$, $n$ positive integers.

Discuss the roots of this equation. \source{Serret, Calc. Intég., p. 254.}

\item In Bernoulli's Lemniscate $r^2=2a^2\cos2\theta$,

show that if $x+\iota y=u$ and $x-\iota y=v$,

the equation of the curve may be written
\[
(u^2-a^2)(v^2-a^2)=a^4.
\]

Further, expressing $u^2$ and $v^2$ as $a^2(1+t^2)$ and $a^2\left(1+\dfrac{1}{t^2}\right)$ respectively, show that the tangent of the angle which the tangent at any point makes with the $x$-axis is
\[
\iota\,\frac{1+t^3}{1-t^3}.
\]

Hence, putting the coordinates of two points at which the tangents are parallel, as $\omega\mu$, $\omega^2\mu$ where $\omega^3=1$, show that the locus of the mid-points of chords joining such points is
\[
[16u^2v^2-8a^2(u^2+v^2)+3a^4]^2
\]
\[
=4a^4[16\{u^4+v^4-u^2v^2\}-12a^2(u^2+v^2)+9a^4],
\]
i.e. a curve of the eighth degree. \source{R. A. Roberts, Proc. L.M. Soc., vol. xviii.}

\item Obtain an integral for the rectification of the inverse of the parabola $y^2=4ax$, with regard to a point on the axis whose coordinates are $(h,0)$.

If $h=-3a$, show that
\[
s=\frac{1}{6a\sqrt2}\log\frac{3+2\sqrt2\sin\omega}{3-2\sqrt2\sin\omega},
\]
where $a\tan^2\omega$, $2a\tan\omega$ are taken as the current coordinates of a point on the parabola, and the arc of the inverse is measured from the point corresponding to the vertex of the parabola. \source{Mr. Roberts, loc. cit.}

Show that the semiperimeter is bisected at the point $\omega=\sin^{-1}\tfrac34$.

\item Show that the tangents to the parabola $y^2=4a(x+a)$ at the points
\[
\{a\sinh^2(u\pm v)-a,\ 2a\sinh(u\pm v)\},
\]
where $u$ is variable but $v$ is a constant, intersect on a confocal parabola; and that if $T$ be a point on this second parabola, and $TP_1$, $TP_2$ the tangents to the first, then
\[
TP_1+TP_2-\text{arc }P_1P_2=a(\sinh2v-2v),
\]
and is constant. \source{Oxford I. P., 1911.}

\item Show that $\displaystyle\iint\frac{dA}{r}$, taken over the area cut from a parabola of latus rectum $4a$ by an ordinate distant $c$ from the vertex ($c<a$), where $r$ denotes the distance from the focus, is equal to
\[
4\sqrt{ac}-2(a-c)\log\frac{\sqrt a+\sqrt c}{\sqrt a-\sqrt c}.
\]
\source{Oxford I. P., 1911.}

\item Show that $\displaystyle\int_0^{\pi/4}\frac{\sin15\theta}{\sin\theta}\,d\theta=\frac{\pi}{4}+\frac11-\frac13+\frac15-\frac17$.

\item If
\[
u=e^{-\int\phi\,dx}\int e^{\int\phi\,dx}(c_0+c_1x+c_2x^2+\dots+c_nx^n)f\,dx+Ce^{-\int\phi\,dx},
\]
where $c_0, c_1, c_2,\dots, c_n, C$ are $(n+2)$ arbitrary constants, and
\[
\frac{\phi}{f}=a_0+a_1x+a_2x^2+\dots+a_mx^m,
\]
where $a_0, a_1,\dots, a_m$ are $(m+1)$ given constants, show that if $m$ be not greater than $n$, $\dfrac{du}{dx}$, obtained by the direct differentiation of $u$ with regard to $x$, contains only $(n+1)$ arbitrary constants. \source{Math. Tripos, 1878.}

\item If $f(m,n)=\displaystyle\int_0^\infty x^m(\cosh x)^{-n}dx$, where $m$ and $n$ are positive integers, each greater than 2, prove that
\[
(n-1)(n-2)f(m,n)=(n-2)^2f(m,n-2)-m(m-1)f(m-2,n-2).
\]
\source{Oxford I. P., 1914.}

\item Given that $a$ and $c$ are positive, show that the limit when $m\to\infty$ and $n\to\infty$ of
\[
\frac1n\left[\frac{1}{a^r}+\frac{1}{\left(a+\dfrac{c}{n}\right)^r}+\frac{1}{\left(a+\dfrac{2c}{n}\right)^r}+\frac{1}{\left(a+\dfrac{3c}{n}\right)^r}+\dots+\frac{1}{\left(a+mn\dfrac{c}{n}\right)^r}\right]
\]
is finite when $r>1$; and find this limit. \source{Oxf. I. P., 1914.}

\item The increase $dS$ in a man's satisfaction $S$ by an increased expenditure $dx$ on a certain commodity, is expressed by the law $dS=\dfrac{\lambda}{x-a}dx$. Similar laws, viz.
\[
dS=\frac{\mu}{y-b}dy,\quad dS=\frac{\nu}{z-c}dz,
\]
hold for two other commodities, where $\lambda,\mu,\nu,a,b,c$ are all positive. Find how the man should expend a given sum $E$ ($>a+b+c$) so that his total satisfaction is greatest. \source{Oxford I. P., 1914.}

Show that the maximum satisfaction is measured by
\[
S=\log\frac{\lambda^\lambda\mu^\mu\nu^\nu(E-a-b-c)^{\lambda+\mu+\nu}}{(\lambda+\mu+\nu)^{\lambda+\mu+\nu}}.
\]

\item Evaluate
\[
\int_{e^u}^\infty\frac{x-1}{x^2-2x\cos v+1}\frac{dx}{\sqrt{x^2-2x\cosh u+1}}.
\]
\source{Oxford II. P., 1914.}

\item Show that the tangent to the curve
\[
3a^2(y-x)+x^3=0,
\]
at the point whose abscissa is $h$, cuts the curve again at the point whose abscissa is $-2h$, and that the area included between the curve and the tangent is $9h^4/4a^2$. \source{Oxf. I. P., 1918.}

\item If $f_1(x)$ and $f_2(x)$ are both polynomials in $x$, show that the integral of $f_1(x)/f_2(x)$ with respect to $x$ can always be written in the form
\[
\phi_1(x)/\phi_2(x)+\log\phi_3(x)/\phi_4(x),
\]
$\phi_1,\phi_2,\phi_3,\phi_4$ also denote polynomials, not necessarily real.

Find the general form of the integral with respect to $x$ of
\[
f_1(x+\sqrt{x^2-1})/f_2(x-\sqrt{x^2-1}).
\]
\source{Oxf. I. P., 1918.}

\item Show that the area bounded by the curve
\[
x=\frac{3at^2}{1+t^3},\quad y=\frac{3at}{1+t^3},
\]
its real asymptote $x+y+a=0$, and by two lines at right angles to this asymptote through the points $t=-a$, $t=0$ of the curve, is
\[
\frac{3a^2}{4}\left\{1+\frac{u^4-1}{(u^2+u+1)^2}\right\},
\]
and find the whole area between the curve and its real asymptote. \source{Oxf. I. P., 1917.}

\item If $\phi(z)$ be a rational function of $z$ without singularities in the range $0\leq z\leq1$, prove that
\[
\int_0^{\pi/2}\phi(\sin2x)\cos^2x\cos2x\,dx=\int_0^{\pi/2}\phi(\sin2x)\cos^2x\cos^22x\,dx
\]
\[
=\int_0^{\pi/2}\phi(\sin2x)\cos^4x\cos2x\,dx.
\]
\source{Oxford I. P., 1907.}

\item Integrate

(i) $\displaystyle\int\frac{x(x-a)^{b-1}dx}{(x-b)\{(x-b)^{2b}-(x-a)^{2a}\}^{\frac12}}$,

(ii) $\displaystyle\int\frac{x^{q-1}\{px^{p+q}-qa^{p+q}\}dx}{(x^{p+q}+a^{p+q})^2+x^{2q}a^{2q}}$.

\item In the curve $\dfrac{15x}{a}=\left(\dfrac{2y}{a}\right)^{\frac23}-\left(\dfrac{2y}{a}\right)^{\frac53}$, show that $s^2=x^2+\tfrac{16}{15}y^2$, $s$ being measured from the origin.

Show that the curve is a quintic of which the $y$-axis is an axis of symmetry, and that the area of the loop $=\left(\dfrac{8}{189}\right)\left(\dfrac53\right)^{\frac13}a^2$.

\item If $2\phi$ be the eccentric angle of the point $r,\theta$ on the ellipse $c=r(1-e\cos\theta)$, prove that
\[
\{(1+e)^2-4e\sin^2\phi\}\left(\frac{d\theta}{d\phi}\right)^2=4(1-e^2\cos^2\theta).
\]
Use the fact that
\[
\int_0^\pi F(\cos^2\theta)\,d\theta=2\int_0^{\pi/2}F(\cos^2\theta)\,d\theta
\]
and the above to obtain a value of $\alpha$, such that
\[
\int_0^{\pi/2}\frac{d\phi}{\sqrt{(1+e)^2-4e\sin^2\phi}}=2\int_0^\alpha\frac{d\phi}{\sqrt{(1+e)^2-4e\sin^2\phi}}.
\]
\source{Oxford I. P., 1917.}

\item A uniform rod of mass $M$ has its extremities at the points $x_1,y_1$; $x_2,y_2$. Show that the product of inertia of the rod with respect to the axes is given by
\[
M\int_0^1\{tx_2+(1-t)x_1\}\{ty_2+(1-t)y_1\}\,dt.
\]
Hence show that the product of inertia of the rod is the same as that of three particles of masses
\[
\frac{M}{6},\quad\frac{M}{6},\quad\frac{2M}{3},
\]
placed at the extremities and the middle point of the rod respectively. \source{Oxford I. P., 1913.}

\item Show that the coordinates of any point on the curve whose intrinsic equation is $s=a(\sec^n\psi-1)$, where $n$ is an odd integer greater than unity, can be expressed rationally in terms of $\tan\psi$, and show that when $x=0$ the curve is a cubic with a cusp. \source{Oxf. I. P., 1911.}

\item Show how to evaluate the integral $\displaystyle\int f(x,y)\,dx$, where
\[
y^2=ax^2+2bx+c
\]
and $f(x,y)$ is a rational function of $x$ and $y$.

Prove that

(i) $\displaystyle\int_0^a\frac{dx}{x+\sqrt{a^2-x^2}}=\frac12\pi$,

(ii) $\displaystyle\int_0^a\frac{a\,dx}{(x+\sqrt{a^2-x^2})^2}=\frac{1}{\sqrt2}\log(1+\sqrt2)$,

the positive sign being taken for the radical in each of the subjects of integration. \source{Math. Trip., Part II., 1913.}

\item Show by means of the transformation $y=\dfrac{(x^2+1)^{\frac12}}{x+1}$ that
\[
\int_0^\infty\frac{dx}{(x+1)(x^2+1)^{\frac12}}=2\int_{1/\sqrt2}^1\frac{dy}{(2y^2-1)^{\frac12}}=\sqrt2\log(\sqrt2+1),
\]
and verify the result in an independent manner. \source{Math. Trip., Part II., 1914.}

\item Integrate $\displaystyle\int\frac{\sin x}{\sin(x-a)}\,dx$. \source{Math. Trip., Part II., 1914.}

\item Evaluate
\[
\int\frac{x+2}{(x+1)^2(x^2+4)},\quad\int\frac{dx}{(x^2+1)^4},\quad\int\frac{dx}{(5-3\cos x)^2},
\]
and the corresponding definite integrals taken between the limits $(0,\infty)$, $(0,\infty)$ and $(0,\pi)$ respectively. \source{Math. Trip., Part II., 1914.}

\item Show that

(i) $\displaystyle\int\frac{\sin4x}{\sin6x}\,dx=\frac{\sqrt3}{6}\tanh^{-1}\left(\frac{\sin2x}{\cos\dfrac{\pi}{6}}\right)$.

(ii) $\displaystyle\int\frac{\sin3x}{\sin5x}\,dx=\frac15\left[\sin\frac{\pi}{5}\log\frac{\sin\left(\dfrac{2\pi}{5}+x\right)}{\sin\left(\dfrac{2\pi}{5}-x\right)}+\sin\frac{2\pi}{5}\log\frac{\sin\left(\dfrac{\pi}{5}+x\right)}{\sin\left(\dfrac{\pi}{5}-x\right)}\right]$.

\item Prove that
\[
\int_0^{\pi/2}\frac{d\theta}{(1+\sin^2\theta)(2+\sin^2\theta)}=\frac{\pi}{\sqrt3}\sin\frac{\pi}{12}.
\]

\item Prove that
\[
\int\frac{2\cos\theta+\sin\theta}{(1+\sin\theta\cos\theta)^{\frac32}}\,d\theta=\frac{2\sin\theta}{(1+\sin\theta\cos\theta)^{\frac12}}.
\]

\end{enumerate}

\end{document}
